%%%%%%%%%%%%%%%%%%%%%%%%%%%%%%%%%%%%%%%%%%%%%%%%%%%%%%%%%%%%%%%%%%%%%%%%%%%%
%% Author template for Management Science (mnsc) for articles with no e-companion (EC)
%% Mirko Janc, Ph.D., INFORMS, mirko.janc@informs.org
%% ver. 0.95, December 2010
%%%%%%%%%%%%%%%%%%%%%%%%%%%%%%%%%%%%%%%%%%%%%%%%%%%%%%%%%%%%%%%%%%%%%%%%%%%%
\documentclass[mnsc]{informs4}

\usepackage{informs4_public}

\usepackage{eqndefns-left} % For checking the display equation width and equation environment definitions %
\RequirePackage{tgtermes}
\RequirePackage{newtxtext}
\RequirePackage{newtxmath}
\RequirePackage{bm}
\RequirePackage{endnotes}
% \OneAndAHalfSpacedXI
%%\OneAndAHalfSpacedXII % Current default line spacing
% \DoubleSpacedXII
\DoubleSpacedXI
% Optional LaTeX Packages
\usepackage{algorithm}
\usepackage{algpseudocode}
\usepackage{tikz}
% Private macros here (check that there is no clash with the style)

% Natbib setup for author-number style
\usepackage{natbib}
 \bibpunct[, ]{(}{)}{,}{a}{}{,}%
 \def\bibfont{\small}%
 \def\bibsep{\smallskipamount}%
 \def\bibhang{24pt}%
 \def\newblock{\ }%
 \def\BIBand{and}%


%% Setup of the equation numbering system. Outcomment only one.
%% Preferred default is the first option.
\EquationsNumberedThrough    % Default: (1), (2), ...
%\EquationsNumberedBySection % (1.1), (1.2), ...

%% Setup of theorem styles. Outcomment only one.
%% Preferred default is the first option.
\TheoremsNumberedThrough     % Preferred (Theorem 1, Lemma 1, Theorem 2)
%\TheoremsNumberedByChapter  % (Theorem 1.1, Lema 1.1, Theorem 1.2)
\ECRepeatTheorems  %  
% If hyperref is used, dvi-to-ps driver of choice must be declared as
%   an additional option to the \documentclass. For example
%\documentclass[dvips,mnsc]{informs3}      % if dvips is used
%\documentclass[dvipsone,mnsc]{informs3}   % if dvipsone is used, etc.

% Private macros here (check that there is no clash with the style)

%https://tex.stackexchange.com/questions/400692/different-height-of-tikz-subfigures-despite-same-axis-width-height
\usepackage{lscape}
\usepackage{booktabs}

\usepackage{array}%[=2016-10-06]
%\usepackage{tabu}
\usepackage{tabularray}
\usepackage{colortbl}

%\usepackage{parskip} % nicer spacing

\newcolumntype{P}[1]{>{\centering\arraybackslash}m{#1}}
\usepackage{multirow,calc,amsmath}
\newlength\mylenA
\newlength\mylenB
\newlength\mylenC
\newlength\mylenD
\usepackage{tikz}
\usepackage{etoolbox}
\usepackage{adjustbox}[valign=t]
\makeatother
\usepackage{fancyvrb}
\DefineVerbatimEnvironment{shadedverbatim}{Verbatim}{
  breaklines=true,
  fontsize=\small,
  bgcolor=gray!12
}
\usepackage{color, colortbl}
\usepackage{graphicx}
%\usepackage{subfigure}
\usepackage{caption}
\usepackage{subcaption}
%\usepackage{tcolorbox}

%\usepackage{amsmath}
%\usepackage{amsthm}
\usepackage{dsfont}
%\usepackage{statmath}
\usepackage{mathtools}
\usepackage{physics}
\usepackage{xcolor}
\usepackage{hyperref}
\usepackage{makecell}
\usepackage[most]{tcolorbox}

\tcbset {
  base/.style={
    arc=0mm, 
    bottomtitle=0.5mm,
    boxrule=0mm,
    colbacktitle=black!10!white, 
    coltitle=black, 
    fonttitle=\bfseries, 
    left=2.5mm,
    leftrule=1mm,
    right=3.5mm,
    title={#1},
    toptitle=0.75mm, 
  }
}

\definecolor{brandblue}{rgb}{0.34, 0.7, 1}
\newtcolorbox{mainbox}[1]{
  colframe=lightgray, 
  base={#1}
}

\newtcolorbox{subbox}[1]{
  colframe=black!30!white,
  base={#1}
}

\newcommand{\bx}{\mathbf{x}}
\newcommand{\bu}{\mathbf{u}}
\newcommand{\by}{\mathbf{y}}
\newcommand{\bA}{\mathbf{A}}
\newcommand{\bB}{\mathbf{B}}
\newcommand{\bX}{\mathbf{X}}
\newcommand{\bY}{\mathbf{Y}}
\newcommand{\bZ}{\mathbf{Z}}
\newcommand{\bI}{\mathbf{I}}
\newcommand{\bH}{\mathbf{H}}
\newcommand{\bW}{\mathbf{W}}
\newcommand{\bxi}{\boldsymbol{\xi}}
\newcommand{\bXi}{\boldsymbol{\Xi}}
\newcommand{\bS}{\boldsymbol{\Sigma}}
\newcommand{\bG}{\boldsymbol{\Gamma}}
\newcommand{\mD}{\mathcal{D}}
\newcommand{\mG}{\mathcal{G}}
\newcommand{\mM}{\mathcal{M}}
\newcommand{\mE}{\mathcal{E}}
\newcommand{\gray}{\rowcolor[gray]{.90}}



\def\checkmark{\tikz\fill[scale=0.4](0,.35) -- (.25,0) -- (1,.7) -- (.25,.15) -- cycle;} 

% \newtheorem{theorem}{Theorem}
% \newtheorem{corollary}{Corollary}[theorem]
% \newtheorem{lemma}[theorem]{Lemma}
\newtheorem{thm}{Theorem}
\newtheorem{remark2}[thm]{Remark}
% \newtheorem{proof}{Proof}


\usepackage{bm}
\usepackage{soul}
\usepackage{float}

\usepackage{algorithm}
\usepackage{algorithmicx}
\usepackage{algpseudocode}

\usepackage{placeins}

\MANUSCRIPTNO{MS-0001-1922.65}

%%%%%%%%%%%%%%%%
\makeatletter
\renewenvironment{thebibliography}[1]
     {\section*{\refname}%
      \footnotesize  % or \scriptsize, \tiny, \footnotesize
      \list{\@biblabel{\@arabic\c@enumiv}}%
           {%
            % --- Flush left, no indent ---
            \leftmargin0pt
            \labelwidth0pt
            \labelsep0.5em % space between [n] and text; set 0pt if you want tight
            \itemindent0pt
            \listparindent0pt
            \usecounter{enumiv}%
            \let\p@enumiv\@empty
            \renewcommand\theenumiv{\@arabic\c@enumiv}}%
      \sloppy\clubpenalty4000\widowpenalty4000%
      \sfcode`\.\@m}
     {\def\@noitemerr
       {\@latex@warning{Empty `thebibliography' environment}}%
      \endlist}
\makeatother

% Neutralize the old hook
\makeatletter
\providecommand{\@openbib@code}{}%
\makeatother

%%%%%%%%%%%%%%%%
\begin{document}
%%%%%%%%%%%%%%%%

% Outcomment only when entries are known. Otherwise leave as is and
%   default values will be used.
%\setcounter{page}{1}
%\VOLUME{00}%
%\NO{0}%
%\MONTH{Xxxxx}% (month or a similar seasonal id)
%\YEAR{0000}% e.g., 2005
%\FIRSTPAGE{000}%
%\LASTPAGE{000}%
%\SHORTYEAR{00}% shortened year (two-digit)
%\ISSUE{0000} %
%\LONGFIRSTPAGE{0001} %
%\DOI{10.1287/xxxx.0000.0000}%

% Author's names for the running heads
% Sample depending on the number of authors;
% \RUNAUTHOR{Jones}
% \RUNAUTHOR{Jones and Wilson}
% \RUNAUTHOR{Jones, Miller, and Wilson}
% \RUNAUTHOR{Jones et al.} % for four or more authors
% Enter authors following the given pattern:
%\RUNAUTHOR{}

% Title or shortened title suitable for running heads. Sample:
% \RUNTITLE{Bundling Information Goods of Decreasing Value}
% Enter the (shortened) title:
\RUNTITLE{The Risk of Extrapolation of Post hoc Explanations}

% Full title. Sample:
% \TITLE{Bundling Information Goods of Decreasing Value}
% Enter the full title:
% %change title!!!
% \TITLE{From Model Explanation to Data Misinterpretation: The Misuse of Post Hoc Explainers for Machine Learning Models in Business Research}
\TITLE{From Model Explanation to Data Misinterpretation: A Cautionary Analysis of Post Hoc Explainers in Business Research}
% \TITLE{The Curse of Rashomon Effects - From Model Explanation to Data Misinterpretation}
% \TITLE{The Risk of Inferring Data Insights from Post hoc Explanations of Machine Learning Models}

% Block of authors and their affiliations starts here:
% NOTE: Authors with same affiliation, if the order of authors allows,
%   should be entered in ONE field, separated by a comma.
%   \EMAIL field can be repeated if more than one author
\ARTICLEAUTHORS{%
\AUTHOR{Tong Wang}
\AFF{School of Management, Yale University, New Haven, CT 0620, \EMAIL{tong.wang.tw687@yale.edu}} %, \URL{}}
\AUTHOR{Ronilo Ragodos}
\AFF{Peter T. Paul College of Business and Economics, 
 University of New Hampshire, Durham, NH 03824, \EMAIL{Ronilo.Ragodos@unh.edu}}
\AUTHOR{Lu Feng}
\AFF{School of Economics and Management, 
 Tsinghua University , Beijing, China, \EMAIL{fenglu@sem.tsinghua.edu.cn}}
\AUTHOR{Yu Jeffrey Hu}
\AFF{Mitch Daniels School of Business 
 Purdue University, West Lafayette, IN 47907, \EMAIL{yuhu@purdue.edu}}
% Enter all authors
} % end of the block

% palette for colorblindness https://davidmathlogic.com/colorblind/#%23000000-%23E69F00-%2356B4E9-%23009E73-%23F0E442-%230072B2-%23D55E00-%23CC79A7

\ABSTRACT{%
Post hoc explainers such as SHAP and LIME are used widely in business research to interpret complex machine learning models. Although they were designed to explain model predictions, there has been an increasing trend in which the explanations they generate are treated as evidence about underlying data relationships. Based on a systematic review of 181 studies, including 56 published in leading journals, we document that this explanation interpretation is widespread and examine its validity.
We also evaluate how well post hoc explanations reliably recover the direction and relative importance of features in true data-generating process. We introduce two metrics—direction alignment and strength alignment—and assess SHAP and LIME using simulated data with known ground truth. Although explanations often appear reasonable on average, they exhibit substantial heterogeneity in their alignment and can fail to be aligned even when their predictive accuracy is high. High predictive performance is therefore necessary but insufficient for reliable explanation.
We further show that feature correlation and the Rashomon effect (where many equally accurate models rely on different feature attributions) are key drivers of misalignment. Agreement in explanations across such models provides a practical diagnostic of reliability. Overall, our findings caution against using post hoc explainers for hypothesis validation and instead position them as exploratory tools that generate, rather than confirm, substantive insights.
%\textcolor{red}{add somewhere that we prove theoretical evidence that SHAP is unsuited for this task}
}%

% Sample
%\KEYWORDS{deterministic inventory theory; infinite linear programming duality;
%  existence of optimal policies; semi-Markov decision process; cyclic schedule}

% Fill in data. If unknown, outcomment the field
\KEYWORDS{post hoc explanation, machine learning, Rashomon Set} %\HISTORY{This paper was first submitted on April 12, 1922 and has been with the authors for 83 years for 65 revisions.}

\maketitle

% use statistics of the citation of the post hoc papers
% show how difficult it is to solve the problem, e.g., cite cynthia's paper
\section{Introduction}
    % Mention Marketing Science Institute "consistently identified attribution as a topmost research priority over the last few years" (from SHAP meets uniform paper)

% =========================== Tong ===========================


Over recent decades, the growth in data volume and complexity has made machine learning (ML) a mainstay of business analytics, particularly for prediction tasks \citep{kleinberg2015prediction,mullainathan17applied}. In these settings, researchers typically solve a supervised learning problem - predicting a target ($Y$) using a set of explanatory variables or features ($X$). However, in business applications, prediction alone is often insufficient. Researchers and practitioners are also interested in gaining insights about their data, in particular, \textit{how explanatory variables ($X$) are associated with the target variable ($Y$)}. Unfortunately, most ML models with high accuracy operate as black boxes\footnote{It is worth acknowledging that recent years have seen the emergence of inherently interpretable models that can attain comparable performance to black box models, sans the reliance on post hoc explanations. However, these models, which lie beyond the purview of this paper's discussion, will not be further delved into here.} that produce predictions without revealing how features relate to the targets.



In response to this opacity, post hoc explainers have naturally garnered significant attention from business researchers. These tools provide various ways to interpret black box models, such as feature importance \citep{ribeiro16lime,lundberg17shap}, partial dependence plots \citep{friedman01pdp}, counterfactual examples \citep{mothilal20dice}, etc. Many recent studies in business research follow a two-stage pipeline that we will refer to as the \emph{Explanation Pipeline}. The Explanation Pipeline, illustrated in Figure \ref{fig:pipeline}, consists of a first stage where a black box ML model is employed for a prediction task, and a second stage where the ML model is explained by a post hoc explainer. This pipeline can illuminate model behavior by, for example, revealing which features appear most influential for the model’s predictions. 


Among post hoc explainers, two are most widely used in empirical research: Shapley Additive Explanations (SHAP) \citep{lundberg17shap} and Local Interpretable Model-Agnostic Explanations (LIME) \citep{ribeiro16lime}. SHAP is theoretically grounded in cooperative game theory and assigns each feature a Shapley value which represents its contribution to the model’s prediction relative to a baseline. Intuitively, SHAP asks how much a feature contributes to the prediction when averaged over all possible subsets of other features, thereby providing a principled attribution of importance. LIME approximates the behavior of a complex model locally around a specific instance with a sparse linear model. It perturbs the input features in the neighborhood of the instance and fits a simple, interpretable surrogate model—typically a linear regression—using proximity-weighted samples. The coefficients of this local surrogate model are then interpreted as measures of feature influence for that particular prediction. Both SHAP and LIME provide information about how features influence a model’s predictions, in terms of both the direction of influence (whether a feature increases or decreases the predicted outcome) and the strength of influence (the magnitude of that effect). In applied research, these explanations are often examined across many observations and aggregated to reflect broader patterns in model behavior.% in terms of both \emph{direction} and \emph{strength}.

Although post hoc explainers are designed to explain predictive models, a growing body of research interprets explanations they produce as evidence about underlying data-generating processes. In particular, these studies generalize explanations that characterize relationships between features and model predictions ($X \mapsto \hat{Y}$)—which are inherently model-dependent—to infer relationships between features and the true outcome ($X \mapsto Y$) in the data. This perspective is increasingly visible across recent business research. For example, \citet{park2024extracting} interpret SHAP explanations as indicators of which smartphone features customers prefer; \citet{jiang24churn} suggest that customers with high values of features associated with large SHAP values are more likely to churn; and \citet{sanchez2023role} argue that SHAP explanations reveal drivers of hotel guest satisfaction and use these explanations to motivate managerial actions. Some recent work has even advocated combining machine learning models with SHAP as an ``alternative to OLS" \citep{berger2023supply}.
To assess how prevalent this practice is, we conducted a manual review of 181 papers that made substantive SHAP or LIME. We find that approximately 42.5\% of these papers interpret post hoc explanations as evidence of data-level relationships, indicating that this practice is widespread in business research.%, we find that .... \textcolor{red}{show prevelance }

\begin{comment}
    \textcolor{red}{we may need to update this paragraph to reflect recent survey}
To gain a deeper understanding of the prevalent practice of employing post hoc explanations in business research, we conducted a thorough review of the literature. Specifically, we surveyed more than 700 papers from the business literature published between 2018 and 2024 \textcolor{red}{do we survey till 2024 or 2025?}  that cited SHAP or LIME  and manually identify 150 that actually used SHAP or LIME to obtain post hoc explanations. 109 of them have been published in top journals. The remaining 41 papers include papers from SSRN that have yet to complete the long review process of business journals. We have meticulously analyzed these 150 papers and found that among them, 42\% used the post hoc explanations to draw inferences about the \textit{data}, rather than the ML \textit{model} being explained, including qualitative claims such as `$X_j$ has a positive/negative impact on $Y$,'' when the explainer provides a positive/negative feature importance for feature $X_j$ or more quantitative claims such as ``$X_1$ is more important to $Y$ than $X_2$ is" if the feature importance of $X_1$ is larger than that of $X_2$. Such use of explainers suggests an emerging paradigm shift in which, through post hoc explanations, ML models are used not only for prediction but also as a means to gain understanding of data. However, the validity of using explainers for this purpose remains unjustified. 

\end{comment}

These interpretations raise important conceptual questions about the scope of validity of post hoc explainers. Specifically, \textbf{to what extent can explanations derived from machine learning models ($X \mapsto \hat{Y}$) reflect ($X \mapsto Y$) in the underlying data-generating process?} To address this question, we conduct extensive experiments to assess whether post hoc explanations reliably recover the direction and relative strength of true $\mathbf{X}\mapsto Y$ relationships. 


Our study proceeds in four steps. First, we identify problematic yet prevalent ways post hoc explainers are interpreted in prior research. Across a wide range of applications, researchers implicitly rely on two broad types of explanatory claims, which we term \textbf{direction} and \textbf{strength} interpretations. To formalize these interpretations, we give mathematical definitions of direction alignment and strength alignment that are applicable to any additive post hoc explainer. Direction alignment evaluates whether the direction of outcome change implied by an explainer under feature perturbations matches the true direction of change implied by the data-generating process. Strength alignment evaluates whether the relative importance ranking induced by an explainer matches the true importance ranking in the data-generating process. %: for SHAP, direction is reflected in how Shapley values change under feature perturbations, and strength is measured by mean absolute Shapley values; for LIME, direction corresponds to the sign of local surrogate-model coefficients, and strength corresponds to their magnitudes aggregated across instances.

Second, we conduct evaluations using simulated datasets in which the ground-truth $\mathbf{X} \mapsto Y$ relationships are known and precisely controlled. We systematically vary multiple data-generation factors to construct a rich set of ground truths spanning a wide range of complexities.
We find that, although explanations often appear accurate on average, their alignment with $\mathbf{X} \mapsto Y$ exhibits substantial heterogeneity with pronounced long-tailed deviations, even when predictive performance is high. This dispersion in outcomes undermines the practical utility of post hoc explanations, as researchers typically lack diagnostics to assess \textit{when or to what extent} an explanation faithfully reflects the true $X \mapsto Y$ relationship.

Third, we investigate the drivers of explanation misalignment—i.e., when explanations of $X\mapsto\hat{Y}$ fail to recover the direction and strength of the underlying $\mathbf{X}\mapsto Y$ relationships. We focuses on:  (i) data-level factors such as correlation structure among features, the interaction of features, etc, (ii) the predictive performance of the trained model, and (iii) the Rashomon Effect \citep{semenova22rashomon, paes23inevitable}. Crucially, we find that high predictive accuracy is \textbf{necessary but insufficient} for post hoc explanations to align with $\mathbf{X}\mapsto Y$. Models with nearly identical predictive performance can yield markedly different explanations, reflecting the presence of large Rashomon sets. Among data characteristics, feature correlation emerges as the dominant driver of misalignment, whereas other forms of data complexity have limited practical impact.

Finally, we identify diagnostic signals that help assess this misalignment. %We find that the disagreement among equally accurate models can serve as a diagnostic signal for when post hoc explanations are likely to misalign with the true $\mathbf{X}\mapsto Y$ relationships in the data. \
We operationalize it through two measurements that we define called the \emph{prediction agreement} and \emph{explanation agreement} within Rashomon sets, and examine how these measures relate to direction and strength alignment of post hoc explanations. We find that explanation-based agreement is strongly correlated with alignment (especially for recovering strength) whereas prediction-based agreement is  less informative. This result clarifies that it is not disagreement in predictions, but disagreement in feature attributions among observationally equivalent models, that signals heightened risk of misinterpreting $\mathbf{X}\mapsto Y$ from post hoc explanations. In this way, Rashomon agreement provides a practical, dataset-level diagnostic for assessing when explanations of $X\mapsto\hat{Y}$ are unlikely to recover the underlying $X \mapsto Y$ structure.



% Their conclusions and putative insights can be roughly categorized into \textcolor{blue}{uses of four} dominant types \textcolor{blue}{of feature importance}. 

% \textcolor{blue}{These first three} claims aim to describe or partially describe the marginal effect of $X$ on $Y$ in the data. 
% First, papers make claims about the qualitative relationships between $X$ and $Y$, most commonly about whether an $X$ has a positive or negative relationship with $Y$. They use language like ``$X_j$ has a positive/negative impact on $Y$,'' when the explainer provides a positive/negative feature importance for feature $X_j$. Second, papers infer the relative importance of the features from post hoc explanations. They use language like ``$X_1$ is more important to $Y$ than $X_2$ is" if the feature importance of $X_1$ is larger than that of $X_2$. Third, some papers would use post hoc explanations to identify the most influential features and use language like ``$X_{j_1}, X_{j_2}$, and $X_{j_3}$ are the most important to $Y$" when the feature importance for $X_{j_1}, X_{j_2}$, and $X_{j_3}$ are greater in magnitude than the others. 

% \textcolor{blue}{In addition to interpreting \emph{feature importance} in terms of marginal effects, researchers also interpret/understand the importance of a feature in terms of how it affects different outcomes related to ML models. Specifically, a feature can be viewed as important because predictive performance seems to be sensitive to the inclusion of that feature. This is the fourth type of claim. The fifth and final type of claim occurs when research assert that a feature is important because varying it allows them to construct counterfactual instances.}


  % \citep{lorentzen2020peeking}. 
      \begin{figure}[t]
        \centering
        %\fbox{\includegraphics[width = \textwidth,trim=250 500 150 250 mm, clip=true]{Figures/pipeline.pdf}}
        \includegraphics[trim={4.75cm 7.5cm 1.125cm 4.25cm},clip, width = 0.9\textwidth]{Figures/pipeline_flow_new.pdf}
        \caption{Explanation Pipeline: the pipeline of using post hoc explainers to explain a machine learning model. Additionally inferring information about the actual marginal effects of the features constitutes mistaken usage of the pipeline.}
        \label{fig:pipeline}
    \end{figure}
     % They build a black box ML model for a complicated task, which could have many variables, or strong interactions of variables, or include unstructured data, such that canonical econometric methods do not easily apply. Then a post hoc explainer is used to generate feature importance for the ML model,  from which they draw some conclusions on the features as if they are using $\boldsymbol{\beta}$s. 
 

 
 % In addition, there are many types of feature importance \citep{AchenChristopherH1982Iaur}. However, except for \citep{lorentzen2020peeking}, the works cited in the previous paragraph did not provide a formal definition of feature importance and justify that the explainer(s) used were actually providing that type of feature importance. Similarly, the basic question of why Shapley values are appropriate feature importance measures is never addressed in these works\footnote{It is argued in \citep{kumar20shap} that Shapley values are not conducive to human understanding of feature importance.}. 


%Therefore, the goal of this paper is to urge the research community to approach post hoc explanations with skepticism and a more critical eye. 
% We argue in this paper that there exist fundamental flaws in using post hoc explanations to elicit information on $\boldsymbol{\beta}$s. 
% The goal of this paper is to conduct a thorough investigation into this new and emerging trend of research methodology. We aim to directly answer the question, 
% ``\emph{Can SHAP and LIME explanations \textcolor{red}{correctly characterize the relationship between $X$ and $Y$ in the data in terms of different popular notions of feature importance?}}''
% with both empirical and theoretical evidence.  

% \textcolor{blue}{To begin, we investigate whether the SHAP and LIME explanations align with the true marginal effects using simulated data from a simulated linear ground truth model. We do this because in such cases the ground truth marginal effects are easy to derive and understand.} 
% For this purpose, we design three \textit{new} metrics that assess post hoc explanations in terms of capturing correct information about the true marginal effects. %The first metric evaluates whether an explanation's feature importance values have the same signs as the corresponding marginal effects. We call it the \textbf{directionality} of an explanation, since it represents how well an explanation identifies the directions in which features influence the target. The second metric evaluates how much the ranking of features by a post hoc explanation agrees with the ranking induced by the marginal effects. We call it the \textbf{concordance} of an explanation. The last metric evaluates whether a post hoc explanation identifies the most important features. We call it the \textbf{relevance} of the explanation.
% These metrics assess the extent to which the explanations are consistent with the marginal effects in terms of the sign, the ranking of features, and the identification of the most important features. %We call the evaluation of these three metrics the test of \textbf{data-alignment}.
%We evaluate the data-alignment of post hoc explanations using simulated data where we know and can control the ground truth marginal effects in the data. We set our ground truth models to be linear because the extent an explainer can produce data-aligned explanations of simple linear models is  indicative of its potential for data-alignment in more general scenarios -- if an explainer cannot give data-aligned explanations of simple linear phenomena, it is very likely unable to be able to explain the complex processes underlying practical datasets. 
% Unfortunately, our findings reveal that SHAP and LIME explanations typically \textit{do not align} with the marginal effects, not even qualitatively (directionality), and the alignment further drops as the feature dimension increases. This misalignment is evident even in straightforward business scenarios with a linear ground truth model and when the black box model being explained demonstrates high predictive performance on data containing all relevant features to $Y$, all of which are mutually independent. 

% \textcolor{blue}{We then investigate more general cases where the ground truth may be a nonlinear function and the intended notion of feature importance may not be in terms of marginal effects. Specifically, we consider the possibilities that \emph{feature importance} may be given in terms of 1) global ranking 2) instance-level ranking 3) predictive power 4) usefulness for counterfactual generation. We delve into the causes of this misalignment, discovering that various factors within the explanation pipeline affect the explainers' ability to extract information about each of the four kinds of feature importance.} These include having correlated features in the data, the predictive performance of the black box model, and specific settings of explainer model parameters. %We corroborate the initial presentation of our results by showing that misalignment pervades an additional 100 simulations where we explain models produced using AutoML that were trained on datasets with varying numbers of features, rows, and label noise.

% \textcolor{blue}{We also offer various theoretical justifications for data misalignment}.
% \textcolor{blue}{Our central claim is that the \emph{Rashomon effect}, in which there exists a multiplicity of distinct models with equivalent (near-optimal) performance, makes use of the explanation pipeline questionable at the outset.} Additionally, we show that SHAP explanations inherently fail to align with the marginal effects, making Shapley values fundamentally inappropriate for any definite claims regarding marginal effects. We show that SHAP estimates situational importance~\citep{StrumbeljErik2014Epma}, 
 % which expresses the consequence of a feature taking on a value different from its mean. Although the computation of situational importance for a feature involves its marginal effect, the information it contributes can be destroyed in the process. Furthermore, we present theoretical grounds for conditions under which LIME might achieve a high degree of data-alignment.

% \textcolor{red}{I think we should say SHAP and LIME separately. }

% We then demonstrate another critical issue of post hoc explanations: inconsistency. Inconsistency refers to the difference in explanations produced by different configurations within the Explanation Pipeline. Using the same simulated data, we show that varying choices along this pipeline can lead to significantly divergent, and occasionally opposing, explanations. It then follows from the view that the purpose of science is to identify phenomena that occur without fail that post hoc explainers have very limited utility to researchers.


% \textcolor{red}{rewrite this part}
% Finally, we propose a set of robustness checks for explanations provided by SHAP and LIME. Although our theoretical and empirical findings suggest that it is not possible to fully eliminate the issue of data-misalignment, they nevertheless reveal opportunities for improvement, which we address. For instance, for SHAP, the most immediate way to increase data-alignment is to divide its Shapley values by the differences between the corresponding feature values and their means; for LIME, increasing some hyperparameters will significantly increase the alignment with the data.
% Another strategy that showed promise involves averaging the explanations derived from a family of accurate models. We also observe that when the training data used by a black box model include all relevant features, all of which are independent, the alignment of explanations of the black box's predictions with the ground truth tends to increase. This suggests that the underlying data quality and structure play a significant role in the effectiveness of post hoc explanations.
% However, other intuitive strategies we tested, such as combining explanations from both SHAP and LIME for a single black box model, did not yield effective results. This highlights the complexity of the issue and the need for more nuanced solutions.
% \textcolor{blue}{It is also important to note that our robustness checks degrees are limited by XYZ...} Consequently, post hoc explainers such as SHAP and LIME should be used with great caution by business researchers, and the appropriate uses of such tools need to be clarified.


Overall, our paper identifies and cautions against a growing trend of the inappropriate use of post hoc explanations to draw inferences about data, leading to data misinterpretation. We argue that post hoc explanations should not be used to \textbf{validate} hypotheses about data; rather, their proper role in the research process is to \textbf{propose} hypotheses that can subsequently be tested through more rigorous methods. Making the distinction between exploration and validation is critical to avoid overstating the reliability or generalizability of insights derived from post hoc explainers and to ensure that business research continues to rest on sound empirical and theoretical foundations.

  The paper is organized as follows. Section~\ref{sec:related_work} reviews related work on SHAP, LIME, and their known limitations. Section~\ref{sec:use_in_b_research} surveys common problematic uses of post hoc explanations and illustrates these misinterpretations using a demonstration dataset.
Section~\ref{sec:main_alignment} presents empirical analyses showing how much post hoc explanations align with the true $\mathbf{X}\mapsto Y$, and Section \ref{sec:factors} investigates the factors that drive the misalignment. Section~\ref{sec:robustness} reports a series of robustness checks.
Finally, we examine whether agreement among approximately equally accurate models can be used to diagnose explanation reliability in Section~\ref{sec:rashomon_agreement}. The paper concludes in Section~\ref{sec:discussion} with a summary of our findings and a discussion of the implications for the appropriate use of post hoc explanations in business research.



\section{Related Work}
    \label{sec:related_work}
    This section introduces SHAP and LIME, reviews critiques of SHAP and LIME in the computer science literature, and then provides a discussion of the use and reception of SHAP and LIME in the business literature.
    To serve as a point of contrast with SHAP and LIME, we also review a standard set of econometrics tools used in business literature to extract data-level insights. Finally, we explain how the \emph{Rashomon Effect} puts the worth of using post hoc explainers to make inferences about $X \mapsto Y$ relationships into great doubt. 

    
   
    \subsection{Post Hoc Explainers - SHAP and LIME}
SHAP \citep{lundberg17shap} and LIME \citep{ribeiro16lime} are the most widely used post hoc explainers. We use them as representative methods because they dominate applied and empirical research across domains.


SHAP (SHapley Additive exPlanations) is a unified framework for interpreting machine-learning predictions through feature attribution.\footnote{It is ``unified" in the sense that its framework generalizes LIME's.} It builds on the concepts of game theory, particularly the Shapley value, to assign an importance value to each feature for a particular prediction. The values represent the contribution of each feature to the difference between the actual prediction and the average prediction over the dataset. A large part of SHAP's appeal is due to theoretical properties of its explanation model, the \emph{consistency} and \emph{local accuracy} properties in particular. Consistency guarantees that if a model changes in a way that makes a feature more important, the SHAP value of that feature will not decrease. Local accuracy ensures that the contributions of all features to a prediction sum up to the actual predicted value minus the average prediction, providing a detailed and accurate decomposition of the prediction. %%There are several variants of SHAP, such as Kernel SHAP \citep{lundberg17shap}, TreeSHAP \citep{lundberg2018consistent} and DeepSHAP \citep{lundberg17shap}. Each variant retains the core principles of SHAP while adjusting the computational strategy to suit the specific model architecture, ensuring that the interpretability benefits of SHAP can be realized across a wide range of machine learning models.
    
    % SHAP's additive feature importance model differs from LIME's in that the coefficient of each feature represent the Shapley value~\citep{shapley97shapley} of that feature. SHAP surrogate models are not fit on sampled training data instances directly like LIME's, but rather on sampled coalition vectors (binary vectors indicating presence of each feature). Shapley values were originally a game-theoretic notion meant to quantify how much each player contributed to the final score of a game. In the context of SHAP, the Shapley values (surrogate model coefficients) represent how much the value of the feature to which they correspond contributes to the final value of the black box model's prediction. 

    
 % Underlying LIME explanations are ridge regressors fitted on data sampled around the input datum being explained. 
Although not used as widely as SHAP in business research, LIME (Local Interpretable Model-agnostic Explanations) is the earliest model-agnostic post hoc explainer, introduced by \citet{ribeiro16lime}. LIME provides local, model-agnostic explanations by approximating a complex predictive model with a simple, interpretable surrogate model—typically a sparse linear regression—fitted in a neighborhood around a focal observation. The coefficients of this locally fitted surrogate capture how variations in each feature within this neighborhood are associated with changes in the model’s prediction. As such, LIME coefficients are commonly interpreted as \textit{approximate local marginal effects} of the underlying model, defined relative to the locality induced by the perturbation and weighting scheme. %Briefly, it explains a black box's prediction for a given input using the following steps. First, it generates synthetic dataset by perturbing the input and sampling in its local neighborhood, with samples weighted by their proximity to the input. A sparse linear model, typically ridge regression, is then fitted to this weighted dataset, and the resulting coefficients are interpreted as local feature importance scores.While LIME is flexible and model-agnostic, its explanations are strictly local and may be sensitive to the choice of sampling and hyperparameters.
%%LIME discretizes numerical features into quartiles to avoid dealing with problems associated with differing scales of features and ensure users do not get confused when interpreting negative weights for negative valued features. The ridge regression's optimization weighs the sampled data with an exponential kernel and includes an $L_1$ (LASSO) penalty. LIME uses the $L_1$ regularization so that the regressor it produces is sparse in the sense that, rather than including all features, only a subset of features that are most influential for predictions near the instance being explained are included. This makes LIME's explanations more concise than they would be otherwise. The ridge regressors are called surrogate models because they are trained to approximate the predictions of the black box model.
    
    %%LIME's ridge coefficients are meant to represent marginal effects, but depending on the LIME variant and the model being explained, they are not the gradient of the black box model at the input being explained. \citet{garreau2020looking} proved that the original LIME formulation's coefficients did not, in general, equal the gradient of the black box model at the input being explained. Instead, each feature importance value is proportional to the partial derivative of the black box at the input datum and the proportionality constants differ by feature. \citet{agarwal21clime} showed that a simpler variant of LIME called C-LIME that does not discretize input features and does not perform weighted ridge regression produced coefficients that are proportional to the gradient of the black box model. %Removing discretization allows C-LIME to recover the gradient because otherwise, local information about the black box's decision boundary is destroyed.
    
    %%In discretizing features, the data used in LIME's regression problem consists of binary features indicating which bins (quantiles) the feature values of sampled data fall into, which destroys the gradient since the local information about the black box's decision boundary is lost. Moreover, C-LIME drops the ridge penalty because enforcing sparsity opposes the goal of recovering the gradient. In the specific case where the model being explained is linear, C-LIME recovers the coefficients of the linear model. This being the case, in this paper, we choose to disable the discretization setting of LIME throughout the paper to improve the quality of LIME explanations.



\subsection{Critiques in Computer Science Literature}
\label{sec:cs_critique}

Concerns about post hoc explanation methods such as SHAP and LIME originated in the machine learning literature, where researchers have examined their limitations along dimensions including prediction accuracy, stability, robustness, and human interpretability. While our work shares the broad objective of rigorously evaluating post hoc explainers, it differs fundamentally in \emph{what} is being evaluated and \emph{for what purpose}.

A central distinction is that most critiques in the computer science literature assess post hoc explainers with respect to the \emph{predictive model} being explained. From this perspective, the model is treated as the primary object of explanation, and the goal is to assess how well an explainer reflects properties of the mapping $X \mapsto \hat{Y}$ learned by the model. In contrast, the perspective adopted in this paper is explicitly \emph{data-level}. We ask whether explainers can recover a minimum amount of information about the \emph{data-generating process} (DGP) $X \mapsto Y$ that is necessary to support common forms of inference made in empirical research—specifically, inference about the \emph{direction} and \emph{relative strength} of feature effects. Accordingly, our focus is not on whether explainers provide comprehensive or human-ideal representations of a model, but on whether they satisfy basic validity conditions implicitly assumed when researchers generalize SHAP or LIME explanations to data-level.

Within the model-level literature, different dimensions of post hoc explanations are evaluated. One stream assesses explanations' ability to reproduce a model’s predictions or decision boundary, effectively treating the explainer as a surrogate function whose outputs should closely match those of the model itself \citep{han22which}. Other work probes explanation correctness indirectly through stress tests, such as sensitivity to irrelevant features or robustness to perturbations \citep{blunsom2019can,nguyen2020quantitative,deyoung20eraser,luss2021leveraging,carmichael21framework,li23xaibenchmark}. A related but more extreme line of research considers adversarial settings in which explanations are deliberately manipulated without altering model predictions \citep{slack19fooling}. %These studies provide important insights into failure modes of post hoc explainers under demanding or adversarial conditions, but they are not designed to assess whether explainers recover even coarse properties of the underlying DGP in benign settings.

A second, complementary strand of work emphasizes \emph{human-centered interpretability}. Common desiderata include \emph{stability}, \emph{completeness}, \emph{simplicity}, and related concepts\citep{doshi17rigorous,markus2021role,agarwal22openxai}. Stability, a requirement that similar inputs--output pairs to receive similar explanations, is typically evaluated locally and pairwise, making it a stronger condition than global notions such as directional consistency \citep{alvarez18robustness,yeh2019fidelity,visani21stability,kelodjou2024shaping}. Completeness requires explanations to contain sufficient information to reconstruct the model’s input--output mapping \citep{nguyen2020quantitative,bassan2023towards}, while broadness concerns the scope over which an explanation applies \citep{nguyen2020quantitative,zhou21evaluating}. Simplicity-based criteria deliberately trade off completeness and broadness to enhance human comprehensibility \citep{nguyen2020quantitative,markus2021role,makridis2023towards}. The most demanding proposals further require explanations to align with psychological models of human understanding or to be interactive and personalized \citep{jacovi2023diagnosing,kaur2022sensible,ehsan2023charting,sokol2020one,dazeley21conversational}.

Our contribution is orthogonal to both strands. We ask whether attributive explainers satisfy two \emph{minimal validity conditions} required for data-level inference: (i) whether the sign of feature attributions reliably reflects the direction of a feature’s impact on $Y$, and (ii) whether the relative magnitudes of attributions meaningfully reflect the relative importance of features in the data. These requirements are substantially weaker than those commonly studied in the computer science literature. Nonetheless, we show that they are frequently violated even under favorable conditions, e.g. when models are well-specified and explainers are used exactly as intended.



\begin{comment}
Concerns about post hoc explanation methods like SHAP and LIME first emerged in the machine learning literature, where researchers questioned various aspects such as incompleteness, stability, interpretability, etc. While our work shares the broad goal of rigorously evaluating post hoc explainers, it differs fundamentally in the object of evaluation. 


A central distinction is that most existing critiques in computer science evaluate post hoc explainers primarily with respect to the \emph{predictive model} being explained, guided by the perspective that an explainer should fully capture the predictive mapping $X \mapsto \hat{Y}$ implemented by a trained model. The perspective adopted in this paper is that explainers only need to extract a minimum amount of information about the \emph{data-generating process} (DGP) $X \mapsto Y$ necessary to support inference about the direction and strength of the features. This distinction is critical; as we shall see, the properties demanded of explainers in prior critiques either overdetermine or are orthogonal to those necessary to support the direction and strength based inferences business researchers want to make. 

Most critiques impose relatively strong desiderata on explanations. One prominent line of work relates to ideas of \emph{fidelity/faithfulness}.  We are interested in how faithfully explanations convey very basic attributes of the DGP, whereas the works in question are interested in how faithfully explainers convey rather comprehensive pictures of the model's inner-workings. Such works often treat the trained model itself as ground truth and require explanations to exhibit good local function approximation properties \citep{garreau20theoretical,garreau2020looking,agarwal2021unification,han22which,rashid2024using,taimeskhanov2025feature}. Other approaches evaluate correctness through proxy criteria related to how (in)sensitive explainers are to (un)important features being present and/or perturbed\citep{blunsom2019can,nguyen2020quantitative,deyoung20eraser,luss2021leveraging,carmichael21framework,li23xaibenchmark}.  %In settings where no clear ground truth explanation exists, some studies substitute pseudo or synthetic ground truths to benchmark explainers \citep{lin2020seeevaluationexplainableartificial,donnelly2023rashomon,li2021crossmodelconsensusexplanationsimage}. Additionally, "Model-Agnostic Interpretability with Shapley Values"
 Emphasis has also been placed on the important yet atypical scenario of explanation corruption due to adversarial attacks meant to corrupt explanations of a model without affecting that model's predictions \citep{slack19fooling}.



A second, complementary strand of work focuses on facets of \emph{human-centered interpretability}. Standard desiderata include \emph{stability}, \emph{completeness}, \emph{simplicity}, as well as synonymous and related concepts \citep{doshi17rigorous,markus2021role,agarwal22openxai}. Stability, the property of producing similar local explanations of the outputs of a model given similar inputs, is desirable for the sake of user trust \citep{alvarez18robustness,yeh2019fidelity,visani21stability,kelodjou2024shaping}. Stability being evaluated at the level of pairs of inputs rather than globally like direction and strength makes it a stronger condition. An explanation of a prediction being complete means that it contains all information necessary to reconstruct the input-prediction mapping \citep{nguyen2020quantitative,bassan2023towards}.The \emph{broadness} of an explanation refers to how broad the scope of its applicability is \citep{nguyen2020quantitative,zhou21evaluating}. Concepts synonymous with simplicity oppose completeness and broadness, requiring that explanations should not contain so much content/information such that they are cannot be easily comprehended by a human \citep{nguyen2020quantitative,markus2021role,makridis2023towards}. Finally, we mention two of the most demanding requests of explainers in the human-centered interpretability literature. First, that the form and content of explanations should be in compliance with psychological models for human understanding and/or personalized at different levels of human organization (e.g. user-level, society-level) \citep{jacovi2023diagnosing,kaur2022sensible,ehsan2023charting}. Second, and relatedly, that explainers be \emph{interactive} \citep{sokol2020one,dazeley21conversational,jacovi2023diagnosing}.% These critiques highlight important limitations of post hoc explainers as tools for model transparency and user trust.




Our contribution is orthogonal to both strands. Rather than asking whether post hoc explanations are ideal or faithful descriptions of a model, we ask whether they satisfy two \emph{minimal validity conditions} required for data-level inference: (i) whether the sign of feature attributions reliably reflects the direction of a feature’s impact on $Y$, and (ii) whether the relative magnitudes of attributions meaningfully reflect the relative importance of features in the data. These requirements are substantially weaker than the desiderata commonly studied in the computer science literature, yet we show that they are frequently violated even under favorable conditions.

% Viewed this way, existing critiques document that post hoc explainers are imperfect descriptions of models; our results demonstrate a more fundamental limitation—namely, that post hoc explanations often fail to meet the basic conditions necessary to support inference about $X \mapsto Y$ relationships in the data. This distinction clarifies why our findings are not implied- by prior critiques and why evaluating direction and strength alignment addresses a logically prior question.

\end{comment}
%     \subsection{Critique in Computer Science Literature}
%         \label{sec:cs_critique}
%         The alarm for post hoc explanations had first gone off in the field of machine learning, as researchers questioned the correctness and validity of post hoc explanations. %However, the notions of correctness and validity in the ML literature differ from the point of view we take in this paper. 
%         However, the ML community has set different standards for post hoc explainers than we do in this paper; the goals and metrics we set for the evaluations in this paper are unique and complementary to those appearing in prior ML literature. The most relevant distinction we would like to point out is as follows. Although this paper and prominent critiques we will cite share the broad goal of (quantitatively) evaluating notions of explanation correctness, the CS critiques focus in on failures by SHAP and LIME to satisfy rather strong assumptions on the properties post hoc explanations should ideally possess such as completeness (that it should fully capture all information about the black box). \textbf{In contrast, our evaluation is much more pragmatic because we emphasize two rather weak properties that explanations would need to have \emph{at minimum} to facilitate investigations of $X\mapsto Y$. Our two properties require that an explanation's feature importance values 1) indicate the features' actual directions of impact on $Y$, and 2) have relative magnitudes that reflect the relative importance of the features in the data.}
%         %The ML community's attention has mainly focused on explanation quality \textbf{with respect to the ML model} since the primary goal of the machine learning field is to develop new models that make better or faster predictions.
%         Overall, there are two main themes in their critiques. They judge explainers by their ability to produce explanations that are 1) \emph{correct}, and 2) able to be easily understood by human users \citep{doshi17rigorous,markus2021role}. %%The first (and more technical of the two) concerns the ability of an explainer's surrogate model to approximate $\mM$. The second consists in the identification of properties of an explainer that lead to the explanations it generates being difficult for humans to make gainful use of in practice. 
%         %%In addition, they consider factors relating the ability of humans to interpret the explanations, most of which do not fall within the scope of our critique. 

%      \textcolor{red}{there are only five papers cited in this part. Ideally, we want to cite multiple papers discussing the same thing, for example some papers study the stability, some papers study the sufficiency of explanations ,etc}

%         Representing the first theme is a research program focused on the \emph{correctness} or \emph{fidelity} of explanations. A great challenge in evaluating explanation correctness is that there is generally not a single, unimpeachable ground truth explanation (that we have access to). One response to this issue is to consider $\mM$ itself the ground truth and require that SHAP and LIME should have good local function approximation properties \citep{garreau20theoretical,garreau2020looking,agarwal2021unification,han22which,rashid2024using,taimeskhanov2025feature}. Alternatively, explainers can be evaluated with respect to metrics that formalize intuitive notions of correctness such as that features a model does not depend on should be assigned zero importance  and that removing or adding features said to be important to a model prediction by an explainer should change the prediction to change substantially \cite{blunsom2019can,nguyen2020quantitative,deyoung20eraser,luss2021leveraging,carmichael21framework,li23xaibenchmark}. Finally, ``the" ground truth can be substituted with either a pseudo ground truth \citep{li2021crossmodelconsensusexplanationsimage} that approximates ``the" ground truth or a synthetic ground truth \citep{lin2020seeevaluationexplainableartificial,donnelly2023rashomon}. 
%         %Critiques of LIME coming from this perspective reveal that failure modes of LIME include failure to approximate $\mM$ due to its hyperparameters, sampling procedure, or default feature discretization step and/or because $\mM$ is either non-linear or non-additive \citep{garreau20theoretical,garreau2020looking,agarwal2021unification,rashid2024using,taimeskhanov2025feature}. In the case that $\mM$ is linear, it has been shown that neither LIME nor SHAP recover the gradient of $\mM$ as feature importance values \citep{han22which}. Another approach judges explainers using proposed evaluation metrics that  formalize intuitive notions of model explanation correctness such as that features a model does not depend on should be assigned zero importance  and that removing or adding features said to be important to a model prediction by an explainer should change the prediction to change substantially \cite{blunsom2019can,nguyen2020quantitative,deyoung20eraser,luss2021leveraging,carmichael21framework}. 
        
    
%         %Works by Garreau and Luxburg provide technical exhibitions of specific cases in which LIME fails, showing for example how the feature discretization process LIME performs by default result in failure modes when explaining non-additive models like decision trees \cite{garreau2020looking,garreau20theoretical}. Similarly, \citep{han22which} uses the view of SHAP and LIME surrogate models as kernel-weighted linear regressions to show that when explaining a linear model neither explainer recovers the gradient of that model as feature importance values.
        
%         % THIS PARAGRAPH CONTRADICTS MY POINT!!!
%         %Other works with less emphasis on the function approximation perspective identify sets of explanation desiderata and define new/custom evaluation metrics accordingly. For example, some works emphasize that SHAP and LIME explanations should assign large (zero) importance to features that do (not) influence a black box model's output \citep{deyoung20eraser,nguyen2020quantitative,luss2021leveraging}.

        
        
%         %%%That post hoc explainers misrepresent the black boxes they explain is a pillar of the ML community's critique. The terms \emph{fidelity} and \emph{faithfulness} are used interchangeably and generally represent the maxims of completeness and soundness outlined by \cite{zhou21evaluating}. An explanation is complete if it is a full characterization of the black box's behavior. An explanation is sound if its characterization is correct.\footnote{In this paper, however, we characterize faithfulness in reference to the ground truth instead of the black box model because we are interested in obtaining insights about reality, not about black box models. In Section \ref{sec:unified_metrics}, we define metrics that capture notions of faithfulness in terms of two interpretations of feature importance.} 
%     %%%\cite{garreau2020looking} and \cite{han22which} consider soundness from a function approximation perspective, directly measuring the difference between an explainer  and the black box model. \cite{garreau2020looking} shows that LIME's approximation error is always bounded away from 0 while \cite{han22which} shows that the data sampling procedures used by kernelSHAP and LIME prevents their surrogate models from properly approximating linear models.%%relying on a view of KernelSHAP and LIME as kernel weighted linear regression models whose input data are sampled by perturbing the input by random noise.They remark that the choice of noise distribution made by kernelSHAP and LIME prevents them from being able to recover the gradient of linear models they explain as feature importance values. It follows that KernelSHAP and LIME cannot necessarily be expected to reliably perform their intended task of interpreting black box model behaviors. 
% %Other fidelity/faithfulness metrics measure particular approximation properties such as the extent to which an explainer assigns zero importance to features the model does not depend on are used as metrics in \cite{nguyen2020quantitative}. Metrics such as Deletion and Insertion metrics are based on the idea that    


%  The second theme is represented by a body of human-centric XAI literature concerned with how easily ``interpretable" post hoc explanations are to end-users. It is accepted that there are two main aspects of explanation interpretability, named by terms synonymous with \emph{correctness} and \emph{simplicity}. There is an ongoing effort to identify salient aspects of the subjective notions of correctness and interpretability, and determine formal definitions of those aspects that can be used to evaluate explanations of black box models in different contexts \citep{nguyen2020quantitative,agarwal22openxai}. %Explainers that satisfy desiderata related to interpretability are said to be able to provide a description of the behavior of a black box model in a form which can be easily understood by humans. That is, the explanations they generate should be simple enough to be intelligible but have sufficient complexity for the sake of completeness and soundness. 
%  The well known critique of LIME's \emph{instability}, i.e.~ capacity for similar inputs leading to dissimilar explanations is a part of this effort \citep{alvarez18robustness,yeh2019fidelity}. %~\cite{alvarez18robustness}, and \emph{deterministic} so that inputs do not have a multiplicity of explanations. 

%  % make concluding distinction between cs and us
 
%  %%Parsimony \citep{nguyen2020quantitative} and clarity \citep{lakkaraju2017interpretable,montavon2018methods,alvarez18robustness} metrics, for example, measure the simplicity and ambiguity of explanations, respectively\footnote{We use the terminology of \cite{markus2021role} here.}. \cite{jacovi2023diagnosing} further requires that the form of explanations match what the psychology literature identifies as natural for human understanding.  Simplicity of explanations beyond the per-feature importance values given by attributive explainers is not desirable in our application because, e.g. artificial enforcement of sparsity could encourage explanations to omit information about the marginal effects and become incomplete and unsound as a result. %For example, LIME's use of a LASSO penalty to make its explanations sparse is undesirable because it could cause a feature to have zero importance when the $\beta_j$ for that feature is non-zero. %Futhermore, stability metrics quantify how different explanations of nearby inputs are. The size of the local Lipschitz constant of the function underlying the explainer is one such example \citep{alvarez18robustness}. Stable explanations can be desirable because an explainer's explanations for two similar inputs being very dissimilar could engender mistrust of the explainer or arouse unwarranted suspicion of unfairness in the underlying decision process. These factors, however, are not within the scope of our discussion.

% %https://www.scopus.com/results/results.uri?s=REFEID%282-s2.0-85044542379%29&sot=record&sdt=cl&origin=resultslist&src=s&sort=cp-f&limit=10&sessionSearchId=66e4328933ff18c7a00c8494bd5ffb1e&cluster=scosubjabbr%2C%22BUSI%22%2Ct%2C%22DECI%22%2Ct%2C%22ECON%22%2Ct%2Bscosrctype%2C%22j%22%2Ct&offset=10
% %https://www.sciencedirect.com/science/article/pii/S0306437922001405
\subsection{SHAP and LIME Use in the Business Literature}
SHAP and LIME have become increasingly adopted in the business literature as practical tools for interpreting complex predictive models. %\citep{meng21online_review,yang21financial,hanauer22boosting,zhou2023user,tchuente2024real}.
Reflecting the applied orientation of the field, existing critiques in the business literature tend to focus on the downstream consequences of using post hoc explanations rather than on their theoretical properties. Specifically, one branch of the literature studies the (negative) ways in which the use of explainable AI can impact use behavior in general \citep{bauer23infoproc,von2025knowing} and in specific contexts like  healthcare \citep{hassan25unlocking} and the identification of false information \citep{lee24false}. Legal regulations requiring model explanations has led to research on the resulting challenges for regulatory compliance \citep{mohammadi2025regulating} in domains such as credit scoring \citep{goodman17gdpr,bucker2022transparency,chen2024interpretable,debock24xaisurvey,ballegeer25stability}. These studies emphasize the practical and institutional implications of post hoc explainers in real-world decision-making contexts.

A second line of work uses identified limitations of SHAP and LIME as motivation for developing alternative explanation approaches. For instance, \citet{fernandez20explaining} contrast feature-importance methods with counterfactual explanations and characterize settings in which counterfactuals more directly capture how features influence decisions. \citet{kim2023rolex} proposes a new explanation method that  offers improvements/generalizations to each of the three stages of LIME explanation (sampling, local model fitting, explanation generation).




%Similarly, \citet{michiels24shapdecomp} study the problem of computing Shapley-value-based explanations when input features are missing and propose methodological refinements to address this issue. \textcolor{red}{they still explain the model, not the data right?}


However, these papers still focus on scenarios involving the intended use of post hoc explainers: explaining the behavior of the model to which they are applied. In contrast, our work studies scenarios in which model-level explanations are used to used as a basis for data-level conclusions. Specifically, we ask whether model explanations \textit{can be used to reliably infer fundamental properties of the underlying data-generating mechanism $\mathbf{X}\mapsto Y$}.





% \textcolor{red}{\citet{slack19fooling} demonstrate that post hoc methods like LIME and SHAP are vulnerable to adversarial manipulation, allowing model explanations to be deliberately distorted through targeted attacks. In contrast, our study investigates how these explainers are used in business research even in non-adversarial settings. Moreover, \citet{fernandez20explaining} focus on explaining AI decision-making processes, contrasting counterfactual explanations with feature-importance methods like SHAP and LIME, and highlighting conditions under which the former better captures how features influence decisions.}




%%A few domain-specific XAI survey papers are also worth mentioning. One published in \emph{Information Sciences}, when listing present challenges in the field, cites causal explanation as barrier to adoption of XAI systems, but does not specifically treat our problem of (causal) explanation of $X\mapsto Y$ rather than $X\mapsto \hat{Y}$\citep{ding22xaisurvey}. Another, on XAI in operations research, mentions... \citep{debock24xaisurvey}. 

 %%However, all previously mentioned papers study the fidelity with respect to the \textit{ML model}, rather than the data, as we do in this paper. Nevertheless, the aforementioned studies partially answer our research question -- if a black box explainer has difficulty even being faithful to an ML model, it would be even more challenging for it to provide correct data insights. 

%%\textcolor{red}{A small but growing body of research has begun to question the reliability of post hoc explanation methods such as LIME and SHAP. Among them, \citep{slack19fooling} highlight their vulnerability to adversarial manipulation, showing that model explanations can be deliberately distorted through attacks. However, such work typically assumes that explanations are valid under normal conditions.  We argue that many papers misinterpret model explanations as insights about the underlying data-generating process—such as marginal effects or causal relationships—thus misapplying SHAP and LIME beyond their intended scope.}
%The underlying supposition of Deletion and Insertion metrics is that removing or adding features said to be important to a prediction by an explainer should change the prediction to change substantially \citep{deyoung20eraser,nguyen2020quantitative,alvarez18robustness,luss2021leveraging}. %Some works evaluate explainers by taking advantage of contexts where feature importance information of a model is known. \cite{blunsom2019can} evaluates explainers of a white box recurrent model by quantifying the explainers' error rates with respect to identifying important tokens in inputs. \cite{carmichael21framework} evaluates additive explainers by calculating the Jaccard index between the sets of effects present in neural networks and GAMs and the effects and associated contributions of features given by explainers for those models. 
    
%  \subsection{Classical Approaches to Investigating $X \mapsto Y$ in Econometrics}

% In contrast to much of machine learning, econometrics places central emphasis on causal interpretation \citep{wold1954causality,haavelmo1944probability,pearl2009causality,heckman15casual}. The primary objective is not only to model predictive relationships, but to use data to support claims about how changes in $X$ affect $Y$.

% A large body of econometric work relies on regression-based methods applied to observational data to identify treatment effects. Common strategies include addressing endogeneity through instrumental variables \citep{angrist01iv} or conditioning on observed confounders via control variables \citep{winship1999estimation}. These approaches typically support statements of the form: ``\emph{Instrumenting for variation in $\langle X_j\rangle$, an $A\%$ increase (decrease) in $\langle X_j\rangle$ leads to a $B\%$ decrease (increase) in $\langle Y\rangle$}'' \citep{baqaee25surplus}. Relatedly, panel-based designs such as difference-in-differences are widely used to approximate experimental variation by comparing treated and control units before and after an intervention \citep{chamberlain1984panel,snow1856mode,baker2025difference}. These methods illustrate how econometric analyses rely on explicit modeling assumptions and identification strategies to support claims about $X \mapsto Y$ relationships in observational settings.

% Beyond reduced-form regressions, econometric research also employs counterfactual and structural approaches that encode causal structure directly through modeling assumptions. These methods use carefully specified models to reason about how outcomes would change under hypothetical interventions, often requiring substantial domain knowledge and manual specification \citep{chou2019linear,duncan2014introduction}. While powerful, such approaches are typically more context-specific and less automated than standard predictive modeling pipelines.

% While machine learning methods are usually designed to learn accurate predictive mappings $X \mapsto \hat{Y}$ based on observed correlations,  econometric tools are designed to support causal interpretation by making identification assumptions explicit and by incorporating substantive knowledge through model specification. This distinction highlights the differing goals and methodological commitments of predictive versus causal analysis.

\subsection{Classical Approaches to Investigating $\mathbf{X}\mapsto Y$ in Econometrics}

Econometric research is primarily concerned with inferring relationships between covariates and outcomes at the level of the data-generating process $\mathbf{X}\mapsto Y$, typically under explicit causal and identification assumptions \citep{haavelmo1944probability,wold1954causality,pearl2009causality,heckman15casual}. Rather than explaining the behavior of a fitted predictive model, these methods aim to estimate how changes in $\mathbf{X}$ affect $Y$ directly from observed data.

Classical econometric approaches rely on regression-based, quasi-experimental, and structural modeling techniques—such as instrumental variables, control strategies, and panel designs—to support interpretable claims about $\mathbf{X}\mapsto Y$ in observational settings \citep{angrist01iv,winship1999estimation,chamberlain1984panel}. Crucially, these inferences are drawn without passing through an intermediate machine-learning model. In contrast to post hoc explanation pipelines—where conclusions about $\mathbf{X}\mapsto Y$ are mediated by both the learned model and the explainer—econometric analyses seek to avoid this additional layer of abstraction, thereby reducing the risk that a model's inductive biases distort inference about the underlying data-generating process.


\subsection{Rashomon Sets and the Rashomon Effect}
A main contributing factor to the misalignment between post hoc explanations and $\mathbf{X}\mapsto Y$ is the Rashomon Effect.
Originally articulated in the context of statistical learning by Breiman, the \emph{Rashomon effect} refers to the existence of a large set of models that achieve comparably strong predictive performance, yet rely on substantially different internal representations or feature usage to generate predictions \citep{breiman2001statistical}. Recent work formalizes this idea through the notion of a \emph{Rashomon set}: a collection of all models whose predictive performance lies within a small tolerance of the optimal model \citep{semenova22rashomon,rudin2024amazingthingscomehaving}. Research stemming from the idea of a Rashomon effect emphasizes that predictive accuracy alone is insufficient to pin down a unique explanation of the data-generating process, as many structurally distinct models may be observationally equivalent.

It is important to recognize that the presence of a Rashomon set in a prediction problem is not a pathological edge case. Intuitively, when a data-generating process is noisy, many distinct functional relationships can fit the observed outcomes almost equally well, especially when only finite samples are available. Noise flattens the empirical risk landscape, reducing the extent to which competing models that rely on different features or structures can be distinguished. Formalizing this intuition, \citet{paes23inevitable} show that the Rashomon effect is inevitable when learning from finite data, implying that ambiguity across well-performing models (in terms of how explanatory variables relate to the target) is a fundamental property of empirical modeling rather than a failure of estimation. Moreover, as highlighted by \citet{rudin2024amazingthingscomehaving}, the resulting Rashomon sets can be large and diverse even in seemingly well-behaved prediction problems, with meaningful consequences for interpretability and scientific inference.

Crucially for our setting, the same data that permit multiple near-optimal models also permit multiple, potentially conflicting explanations. Because each model in the Rashomon set encodes a different plausible mapping from $X$ to $Y$, a post hoc explanation of any single model reflects only one of many admissible interpretations. As a result, even when a model achieves high predictive accuracy, its explanations need not align with the underlying data-generating relationship. This multiplicity of admissible explanations is a central challenge we address in this paper.

Consistent with this perspective, feature importance, or more generally, model explanations, can vary widely across models within a Rashomon set. Even when models exhibit near-identical predictive accuracy, they may disagree on individual predictions and on how input features contribute to those predictions, making post hoc explanations particularly fragile on unseen or unlabeled data \citep{hsu2022rashomon}.



    \section{Current Uses and Interpretations of SHAP and LIME in Business Research}
        \label{sec:use_in_b_research}
        To assess how post hoc explanation methods are currently applied and interpreted in business research. We first describe the typical \emph{Explanation Pipeline} employed in prior studies, which involves training a predictive model and subsequently applying an explainer to interpret it. Then, we analyze the types of claims researchers draw from explanations.
 Finally, we conduct a comprehensive review of the prior studies that made substantive use of SHAP or LIME to show the prevalence of the identified interpretations. %This section first summarizes the results of our literature review, documenting the prevalence of studies that use post hoc explanations to infer data-level relationships. 

\subsection{The Explanation Pipeline in Business Research}
    \label{sec:pipeline_demo}


% We formalize an Explanation Pipeline that reflects how post hoc explainers are commonly used in the business research literature and illustrate this pipeline using a synthetic dataset.\footnote{Real datasets are less suitable for the present study because evaluating the faithfulness of post hoc explanations requires access to the true underlying relationships.}
    
An Explanation Pipeline consists of two stages. In the first stage—the prediction stage—a machine learning model $\mathcal{M}$ is trained on data $\mathcal{D}$ to predict an outcome of interest, and in the second stage—the explanation stage—an explainer model $\mathcal{E}$ (e.g., SHAP or LIME) is applied to $\mathcal{M}$ to generate post hoc explanations.
    
To demonstrate how the Explanation Pipeline is typically applied, we simulate an app-downloading dataset \footnote{Real datasets are less suitable for the present study because evaluating the faithfulness of post hoc explanations requires access to the true underlying relationships.}. The dataset includes two demographic features (``Income" and ``Age") and eight features describing users’ app-store behavior such as ``Time Spent on Phone," ``Data Plan Unlimited," ``Peer Influence," etc.
 We generate the labels using a linear ground-truth model that simulates individuals’ decisions about whether to download a certain mobile app. 
We then create the dataset $\mathcal{D}$ with 5000 data points by sampling from the distributions of features to simulate user profiles ($\mathbf{X}$) and using $\mathcal{G}$ to create the labels ($Y$). See Appendix \ref{apx_sec:app_details} for more details of the data generation. 

Following the Explanation Pipeline, we first train an XGBoost classifier $\mathcal{M}$ using 5-fold cross-validation and a grid search for hyperparameter tuning. The resulting model achieves an accuracy of approximately 72\% on the held-out test set. We then apply SHAP and LIME to $\mathcal{M}$ to generate explanations. 






\subsection{Types of Interpretations Drawn from SHAP and LIME}
\label{sec:interpretations}
In this subsection, we illustrate how SHAP and LIME explanations are interpreted.
\subsubsection{SHAP Interpretations}\label{sec:shap_interpretation}
SHAP provides feature attributions that represent the marginal contributions of each feature, i.e., how much a feature adds to the model’s prediction when considered across all possible combinations of features. When aggregated across observations, SHAP attributions are often visualized using a beeswarm plot (Figure \ref{fig:shap_beeswarm}), which displays the distribution of SHAP  values for each feature. Each dot represents a single observation’s SHAP value for a given feature, with color indicating the corresponding input value (e.g., red for high, blue for low). The horizontal position of a dot shows both the direction and magnitude of the feature’s contribution; points to the right indicate positive contributions to the predicted outcome (higher $\hat{Y}$), whereas points to the left indicate negative contributions (lower $\hat{Y}$).

From this visualization, researchers typically draw two types of interpretations: \emph{direction interpretation} and \emph{strength interpretation}.


\paragraph{Direction Interpretation — Larger Shapley values correspond to larger model predictions ($\hat{Y}$)}
Under this interpretation, direction refers to how feature values are ordered with respect to the model’s predictions: feature values that correspond to larger Shapley values are interpreted as leading the model to produce larger predicted outcomes $\hat{Y}$, and vice versa. For example, if higher values of ``Peer Influence" are associated with larger Shapley values, researchers interpret this as indicating that greater peer influence leads to larger predicted likelihoods of downloading the app. Conversely, if higher ``Privacy Concern" corresponds to smaller Shapley values, this is interpreted as indicating lower model predictions.

\paragraph{Strength Interpretation — Mean absolute Shapley values indicate a feature’s importance to the model}
Under this interpretation, a feature’s strength is quantified by its aggregated contribution, measured as the mean absolute Shapley value across observations. This aggregation reflects how strongly the model relies on that feature, on average, when generating predictions. Features with larger mean absolute Shapley values are therefore interpreted as being more influential in determining $\hat{Y}$. For example, ``Peer Influence" has the largest mean absolute Shapley value, indicating that it contributes the most, on average, to the model’s predictions across the dataset. 
\begin{figure}
\centering
\begin{subfigure}{.6\textwidth}
  \centering
  \includegraphics[width=.9\linewidth]{Rebuttal-R1/Rebuttal_Figures/shap_ex0p.png}
  \caption{SHAP Beeswarm}
  \label{fig:shap_beeswarm}
\end{subfigure}%
\begin{subfigure}{.45\textwidth}
  \centering
  \includegraphics[width=.85\linewidth]{Rebuttal-R1/Rebuttal_Figures/lime_ex0p.png}
  \caption{LIME Local}
  \label{fig:sub2}
\end{subfigure}
\caption{SHAP and LIME Demo}
\label{fig:sec_3_shap_lime_demo}
\end{figure}




\subsubsection{LIME Interpretations}
LIME provides local, model-agnostic explanations by approximating a complex predictive model with a sparse linear regression fitted in a neighborhood around a focal observation. %The coefficients of this locally fitted surrogate capture how variations in each feature within this neighborhood are associated with changes in the model’s prediction. As such, LIME coefficients are commonly interpreted as \textit{approximate local marginal effects} of the underlying model, defined relative to the locality induced by the perturbation and weighting scheme.

Similarly, researchers typically draw two types of interpretations analogous to those used for SHAP:

\paragraph{Direction Interpretation — The sign of a coefficient indicates whether increasing the feature raises or lowers the model’s prediction locally.}
A positive coefficient implies that, in the local region around the focal observation, an increase in that feature moves the surrogate model’s prediction upward (i.e., increases the predicted likelihood of the outcome). For example, ``Peer Influence" has a positive LIME coefficient for an individual user, then higher peer influence increases the locally predicted likelihood of app downloads for that specific user. ``Privacy Concerns" has a negative coefficient, then within that local neighborhood, greater privacy concerns reduce the predicted likelihood of downloading the app.

\paragraph{Strength Interpretation — The magnitude of a coefficient reflects the strength of the feature’s local influence on the model.}
The absolute magnitude of a LIME coefficient reflects how influential the feature is in shaping the surrogate model’s prediction for that specific observation. Larger absolute coefficients indicate stronger local effects. Thus, for a given observation, the most influential feature is the one with the largest absolute LIME coefficient, followed by features with progressively smaller magnitudes. 

As we will show in the next section, direction and strength interpretations summarize the primary ways in which SHAP and LIME outputs are used to infer model behavior in applied business research.  


\subsection{The Controversial Generalization to Interpreting $X \mapsto Y$ in the Literature}
The interpretations described in the previous subsection are formally statements about how a trained model maps inputs to predictions ($X \mapsto \hat{Y}$). The controversial practice we investigate arises when these model-level interpretations are generalized into claims about the underlying relationship between features and outcomes in the data-generating process ($X \mapsto Y$). Under this generalization, quantities defined to describe model behavior are implicitly treated as evidence about how features influence outcomes in the data.

In the applied literature, this shift in interpretation is typically not made explicit. Instead, language that is formally appropriate for describing model behavior is framed as describing how features are associated with the outcome in the data. As a result, direction interpretations are often read as indicating whether a feature corresponds to larger or smaller values of the outcome in the data, while strength interpretations are read as indicating the relative prominence or importance of features with respect to the observed outcome.

Boxes below provide prototypical examples of this linguistic shift. In each case, explanations that describe how a model’s prediction varies with features are articulated as statements about corresponding variation in the outcome $Y$, rather than the predicted outcome $\hat{Y} $.
\begin{mainbox}{Examples of controversial claims about the direction of influence} 
\begin{itemize}
    \item ``\textit{Positive SHAP values indicate that the variable has a facilitative effect on $\langle Y\rangle$, while negative values indicate a suppressive effect.}''
    \item ``\textit{Individuals whose predictions received larger SHAP values for this feature are more likely to have $\langle Y=1\rangle$.}''
    \item ``\textit{Choosing features with positive mean SHAP values is recommended because they have positive effects on $\langle Y\rangle$.}'' 
\end{itemize}
\end{mainbox}

\begin{mainbox}{Examples of controversial claims about the strength of influence}
\begin{itemize}
    \item ``\textit{This study provides insights for planning the design of $\langle X\rangle$ by assessing the relative importance and rankings of the $\langle X_j\rangle$ in predicting $\langle Y\rangle$.}''
    \item ``\textit{Shapley values enable us to identify and rank relevant $\langle X_j\rangle$ for the investigated $\langle Y\rangle$.}''
    \item ``\textit{Feature rankings and visualizations of impact patterns provide policymakers with evidence for strategies to improve $\langle Y\rangle$.}"
\end{itemize}
\end{mainbox}


% Whether such generalizations are valid is an empirical question.  In the following analysis, we use simulated data with known ground truth to evaluate the extent to which SHAP- and LIME-based interpretations accurately recover the underlying $X \mapsto Y$ relationships.




% We conducted two waves of literature search and review. In the first wave, we used Web of Science to identify publications with either ``SHAP'' or ``LIME'' appearing in the body of the paper within business-related subject areas, including management, economics, supply chain \& logistics, and operations research \& management science. This search yielded over one thousand papers published between 2017 and January 2026. Because business journal publications typically involve long review cycles, we complemented this search with SSRN to capture working papers. We then manually examined each paper to identify studies that substantively used SHAP or LIME for model explanation, rather than merely citing or mentioning these methods. We stopped once we had identified 150 such papers. 

% In the second wave, to ensure comprehensive coverage of papers published in leading journals, we focused on journals listed in the \textit{UTD 24}, the \textit{FT50}, and journals published by \textit{INFORMS}. Using the same search keywords (``SHAP'' or ``LIME''), we manually examined all returned papers and retained those that made substantive use of SHAP or LIME. This screening process resulted in 56 papers. 

% After removing overlaps between the two waves, our final sample consists of 181 unique papers that make substantive use of SHAP and/or LIME. Among these papers, 94\% used SHAP and 13\% used LIME, with 7\% using both methods.

% We next classified papers based on how post hoc explanations were interpreted and used. We screened papers based on whether authors explicitly or implicitly framed explanations as describing properties of the trained model, or instead extrapolated them to make claims about relationships in the underlying data. Papers exhibiting such extrapolation were coded as making data-level inferences. Based on this criterion, we categorized papers into three groups: controversial papers that clearly extrapolated post hoc explanations to data-level claims; ambiguous papers whose wording could plausibly be interpreted as making such extrapolations; and proper papers that used explanations strictly for model interpretation, consistent with the intended scope of these methods.


% We then characterize how post hoc explanations were interpreted and used. Our screening criterion was whether authors explicitly or implicitly framed explanations as describing the \emph{models}, rather than properties of the underlying data. For papers that failed this screening, we identified language suggesting extrapolation from model explanations to claims about relationships in the underlying data.  We observed two recurring types of misinterpretation. First, some papers interpreted the \emph{direction} of SHAP or LIME values as reflecting the direction of influence of features on the outcome in the data, implicitly inferring that larger (or positive) explanation values correspond to larger values of $Y$. Second, other papers interpreted the \emph{magnitude} of explanation values as reflecting the strength or importance of features in the data-generating process, often treating features with larger explanation values as primary levers for intervention. We recorded representative quotations illustrating these problematic interpretations.

% Based on these distinctions, we categorized papers into three groups: \emph{controversial} papers that clearly extrapolated post hoc explanations to make data-level inferences; \emph{ambiguous} papers whose wording could plausibly be interpreted as making such extrapolations; and \emph{proper} papers that used explanations strictly for model interpretation, consistent with the intended use of these methods. Overall, we found that 44\% of the sampled papers fell into the controversial category with an additional 15\% classified as ambiguous.

\paragraph{Prevalence of Data-Level Inference}

To assess the prevalence of the interpretive practices described above, we conducted a systematic review of the literature.

We compiled a sample of 181 papers published in journals included in the \textit{UTD 24}, the \textit{FT50}, and \textit{INFORMS} lists, as well as papers indexed in Web of Science and SSRN (to account for the long publication timelines of academic journals). All selected papers made substantive use of SHAP or LIME. Our objective was to quantify (i) the prevalence of directional and strength-based interpretations of feature attributions and (ii) the prevalence of data-level inference based on post hoc explanations.

To mitigate potential human bias in coding, we complemented manual annotation with an LLM-assisted review of each paper. Final classifications were determined through an adjudication process that reconciled disagreements between human coders and LLM outputs. Detailed documentation of the review protocol, coding scheme, and adjudication procedure is provided in Online Appendix \S 1. 

Across the full sample,   94\% used SHAP and 13\% used LIME, with 7\% using both methods. 75.4\% of papers made directional interpretations of feature attributions, and 97.8\% made strength-based interpretations. These two dimensions constitute the core components of our coding framework and serve as the conceptual basis for identifying subsequent data-level inference.

Overall, 42.5\% of the sampled papers engaged in what we classify as problematic interpretations of post hoc explanations. The prevalence of such interpretations is substantially lower in leading journals: 14.3\% among \textit{UTD 24} journals, 14.0\% among \textit{FT50} journals, and 17.4\% among \textit{INFORMS} journals. Across the union of these three journal lists, the overall rate is 16.1\%.


% Among these papers, we specifically focus on papers published in three major journal sets— \textit{FT50}, and \textit{INFORMS}. %—and summarize their characteristics in Table~\ref{tab:literature_review}. As shown in the table, the incidence of controversial usage among top-journal publications remains substantial, though lower than in the broader, journal-agnostic sample. At the paper level, 40.5\% of UTD 24 papers, 40.0\% of FT50 papers, and 39.1\% of INFORMS papers exhibit controversial usage. These rates are below the 44\% prevalence observed in the full sample, indicating that conditioning on journal prestige is associated with a lower—but still nontrivial—incidence of data-level inference.

% \begin{table}[h]
% \centering
% \small
% \caption{Detailed Summary of SHAP and LIME Usage Across Journal Sets}
% \label{tab:literature_review}
% \begin{tabular}{lccc}
% \toprule
% \textbf{Journal Set} & \textbf{UTD24} & \textbf{FT50} & \textbf{INFORMS} \\
% \midrule

% \multicolumn{4}{l}{\textbf{Paper-Level Statistics}} \\
% \addlinespace[2pt]
% \quad Total papers using SHAP and/or LIME & 42 & 50 & 23 \\
% \addlinespace[2pt]
% \quad Controversial LIME usage / LIME papers & 1 / 2 & 2 / 5 & 1 / 1 \\
% \quad Controversial SHAP usage / SHAP papers & 16 / 40 & 18 / 46 & 8 / 22 \\
% \addlinespace[2pt]
% \quad Papers with any controversial usage & 17 & 20 & 9 \\
% \quad Controversial usage (\% of papers) & 40.5\% & 40.0\% & 39.1\% \\

% \midrule
% \multicolumn{4}{l}{\textbf{Journal-Level Statistics}} \\
% \addlinespace[2pt]
% \quad Journals represented in review & 15 & 22 & 7 \\
% \quad Journals with LIME usage & 2 & 5 & 1 \\
% \quad Journals with SHAP usage & 15 & 21 & 7 \\
% \addlinespace[2pt]
% \quad Journals with any controversial usage & 11 & 13 & 5 \\
% \quad Controversial usage (\% of journals) & 73.3\% & 59.1\% & 71.4\% \\

% \bottomrule
% \end{tabular}

% \smallskip
% \footnotesize
% \emph{Notes.} 
% Entries of the form $a/b$ indicate the number of papers with controversial usage divided by the total number of papers using the corresponding method. 
% \end{table}\normalsize


% As shown in Table~\ref{tab:literature_review}, the incidence of controversial usage among papers published in top journals is substantial but clearly lower than in the broader, journal-agnostic sample. At the paper level, 40.5\% of UTD~24 papers, 40.0\% of FT50 papers, and 39.1\% of INFORMS papers exhibit controversial usage. These percentages are markedly below the over-44\% rate observed in the randomly sampled set of 181 papers, indicating that conditioning on journal prestige is associated with a lower—but still nontrivial—prevalence of controversial interpretations of post hoc explanations.

% From the perspective of \emph{journals}, the pattern remains striking. A majority of journals represented in each group have published at least one paper with controversial usage—73.3\% in the \textit{UTD24} set, and 59.1\% in the \textit{FT50} set, and  71.4\% in the \textit{INFORMS} set. Thus, while top journals exhibit lower paper-level incidence than the randomly sampled literature does, controversial interpretations are nonetheless widespread across outlets rather than isolated to a small subset of journals.

% \begin{table}[htbp]
% \centering
% \caption{Statistical Summary of LIME and SHAP Usage in 74 Reviewed Papers Across Journal Sets}
% \label{tab:literature_review}
% \begin{tabular}{lccc}
% \toprule
% \textbf{Journal Sets} & \textbf{UTD24} & \textbf{FT50} & \textbf{INFORMS} \\
% \midrule
% \multicolumn{4}{l}{\textbf{Paper Statistics}} \\
% \quad Total Papers that used LIME and SHAP & 64 & 75 & 42 \\
% % \quad Papers Using LIME & 13 & 16 & 12 \\
% % \quad Papers Using SHAP & 54 & 63 & 35 \\
% \quad Papers with Controversial Use of LIME / Papers using LIME & 1 / 14 & 2 / 16 & 1 / 12 \\
% \quad Papers with Controversial Use of SHAP / Papers using SHAP & 20 /54  & 22 / 63 & 11 / 35\\
% \quad Papers with Any Controversial Usage & 21 & 24 & 12 \\
% \quad Controversial Usage (\%) & 24.7\% & 28.2\% & 14.1\% \\ %stopped here
% \midrule
% \multicolumn{4}{l}{\textbf{Journal Statistics}} \\
% \quad Journals Represented & 18 & 28 & 10 \\
% \quad Journals with LIME Use & 8 & 11 & 7 \\
% \quad Journals with SHAP Use & 16 & 25 & 8 \\
% \quad Journals with Controversial Usage & 12 & 14 & 5 \\
% \quad Controversial Usage (\%) & 66.7\% & 50.0\% & 50.0\% \\
% \bottomrule
% \end{tabular}

% \smallskip
% \footnotesize
% Note: Statistical summary based on a comprehensive review of 83 papers across three journal sets. 
% "Journals Represented in Review" indicates the number of distinct journals from each set that were included in the review.
% \end{table}





% Building on these findings, we further examine how SHAP and LIME are used in the controversial papers identified above. Specifically, we observe that their interpretations of feature importance values often fall into two categories: (1) identifying the directionality of each feature’s impact (i.e., the sign of importance), and (2) ranking features by relative importance.% and (3) selecting a subset of the most important features. 
% Table 2 presents representative quotations illustrating these three usage patterns. In Section 6.1, we introduce quantitative metrics to assess the degree to which explainers can reliably support these interpretive tasks. We refer to these metrics as directionality and concordance, respectively.




%         As part of our literature review process, we cataloged papers at the intersection of business and explainable machine learning.
        
%         We searched for papers from three sources. First,  we used Web of Science %and SSRN\footnote{We searched on SSRN because business journal papers usually take a few years to get published so there are many preprints available on SSRN.} 
%         to search for publications citing at least one of LIME~\citep{ribeiro16lime} or SHAP~\citep{lundberg17shap} and have either ``LIME" or ``SHAP" appear in the body of the paper. The results were filtered by topics related to business, such as management, economics, supply chain \& logistics, operations research \& management science, etc. More than 700 papers were identified. See Figure \ref{apx_tab:citation_count_graph} for a breakdown of the number of papers by year.
%                \begin{figure}[ht]
%             \centering
%             \includegraphics[width=.5\textwidth]{Figures/citation_count_new.pdf}
%             \caption{Year-by-year breakdown of number of papers citing either LIME or SHAP that belong to business-adjacent Web of Science categories.}
%             \label{apx_tab:citation_count_graph}
%         \end{figure} %\tw{I don't like the bar plot. Change the bar plot to gray (black edges and gray fills) }
%         The second source of papers is INFORMS journals. We searched with the same keywords on the INFORMS journal portal. Finally, we searched on SSRN for papers that are not published yet. Because business journal papers usually take a few years to get published, there are many preprints available on SSRN. The total number of papers from these three sources that matched our keywords was over 700. 
        
%         We then iterated through the resulting papers one by one, collecting papers that actually used SHAP or LIME to explain a machine learning model (many papers only mentioned SHAP or LIME in the context of related work or future work) until we had a total of 150 different papers. 
%         Out of these papers, 91\% used SHAP and 14\% used LIME to generate explanations for their models, including 5\% using both of them.


        
%         We took a close look at these papers and characterized their use of post hoc explanations by the following procedure. We first searched for evidence that it was understood that the intended use of SHAP and LIME is to explain model outputs (rather than the data). On those that fail this screening, we identified problematic language indicative of inappropriate use of post hoc explanations and recorded quotations of problematic sentences.
%         Not all papers contained such language; in some papers it was clear or made explicit that authors used SHAP and LIME only within the scope of model explanation. Having identified any quotations which we deemed problematic, we then divided the papers into three categories: the \textit{controversial} papers that clearly extrapolated explanations to glean information about the marginal effects, the \textit{ambiguous} papers whose wording could suggest extrapolation, and the \textit{proper} papers that only used explanations for model interpretation (the intended use).% Note again, though, that both SHAP and LIME fall short for their indented purpose and face intense scrutiny from the ML community \cite{kumar20shap,garreau2020looking}.

% %\textcolor{red}{add a summary of the findings \&statistics from the 116 papers}
% %\textcolor{red}{double check and update the numbers}
%         We found that 42\% of the papers are controversial - they clearly extrapolated explanations to gain insight into the marginal effects in the data. In addition to that, there is another 15\% ambiguous papers that at least used unclear language such that readers could interpret their findings as being made about the marginal effects. This means that \textbf{more than half} of the papers we sampled had some problematic language. Further, \underline{76\% of the controversial papers were from top journals}.\footnote{A journal was considered a ``top journal" if it is an INFORMS journal or its most recent 5-year impact factor according to Clarivate's Journal Citation Reports (JCR) was at least 5.0 at the time of writing (June 2024).}
%         %Given the popularity of SHAP and the idea that the presence of this problem in publications from top journal is indicative of the pervasiveness of the problem, we also separately surveyed papers from the INFORMS family of journals that used SHAP to explain a machine learning model. About 25\% of these papers had clear extrapolations, and 17\% could be construed as having extrapolated. Thus it was still the case that about half of the papers using SHAP contained problematic language.
%         We believe that this result is indicative of a rising trend of using SHAP or LIME to gain data insights.

%         The interpretation of SHAP and LIME feature importance values in the controversial papers fell into three categories: 1) identifying the directionality of impact of each feature (i.e., feature importance sign), 2) ranking features by relative importance, and 3) identifying a subset of the most important features. Table \ref{tab:metric_examples} below records prototypical ways in which authors expressed their desire to perform these tasks. In Section \ref{sec:evaluating_faithfulness}, we define metrics to characterize the ability of an explainer to perform these tasks. We call them directionality, concordance, and relevance, respectively. 



% ANONYMIZED VERSION
% \begin{table}[htbp]
% \small
% \caption{Correspondence Between Paper Language and Our Metrics. Statements appearing in angle braces $\langle \cdot \rangle$ 
% are meant to make the statements more general. The $X_j$ represent data features and $Y$ represents the outcome of interest.}
% \begin{tabu} to 1\linewidth {  P{3cm} | X[l,m]  }
% \toprule
% \textsc{Purpose} & \textsc{Examples} \\
% \hline
% \textbf{Identifying How Features Impact the Target (Directionality)}
% & ``...Positive SHAP values indicate that the variable has a facilitative effect on $\langle Y\rangle$, while negative values indicate a suppressive effect..." \\ 
% & \vspace{-.8cm}``...Individuals whose predictions received larger SHAP values for this feature are more likely to have $\langle Y=1\rangle$..."\\
% & ``...Choosing features with positive mean SHAP are recommended options because they have positive effects on $\langle Y\rangle$..." \\
% \hline
% \textbf{Ranking the Features by Relative Importance (Concordance)}
% & ``...This study provides insights for planning the design of $\langle X\rangle$ by assessing the relative importance and associated rankings of the $\langle X_j\rangle$ in predicting $\langle Y\rangle$..."\\\smallskip
% & \vspace{-1.3cm}``...The use of Shapley values enables us to identify and rank relevant $\langle X_j\rangle$ for the investigated $\langle Y\rangle$.."\\\smallskip
% & \vspace{-.7cm}``...Our feature rankings and visualizations of impact patterns provide policy makers with evidence for strategies to improve $\langle Y\rangle$..." \\
% \hline
% \textbf{Identifying the Most Important (Top-$k$) Features (Relevance)\hspace{3cm}}
% & ``...That $\langle X_1$, $X_2$, and $X_3\rangle$ have the top-3 SHAP values and that their SHAP values are positive reveals that increasing $\langle X_1$, $X_2$, and $X_3\rangle$ would lead to larger $\langle Y\rangle$..."\\\smallskip
% & \vspace{-.8cm}``...We use SHAP explanations of an XGBoost Classifier to identify the key features which contributing to the risk $\langle Y\rangle$..."\\\smallskip
% & ``...By using XAI to identify the most important factors and indicators that contribute to $\langle Y\rangle$, their values and fluctuations can then be monitored over time..." \\
% \hline
% \textbf{Other Usage}
% & (\textbf{OLS Alternative}) ``...The combination of a black box algorithm and SHAP provides similar economic insights as provided by pooled OLS..."\\\smallskip
% & (\textbf{Theory Confirmation}) ``...SHAP is used to interpret the quantitative associations found by XGBoost; our results provide empirical support for prior theoretical findings..."\\\smallskip
% & (\textbf{Recourse Advocacy}) ``...Since $\langle X_1$, $X_2$, and $X_3\rangle$ have high SHAP importance and substantial effects on $\langle Y\rangle$, $\langle X_1$, $X_2$, and $X_3\rangle$ should be given attention by policy makers to increase $\langle X_1$, $X_2$, and $X_3\rangle$..." \\
% \bottomrule
% \end{tabu}
% \label{tab:metric_examples}
% \end{table}







%\section{Misalignment of Post Hoc Explanations}
\section{Do Post hoc Explanations Align with Data?}
\label{sec:main_alignment}

In this section, we assess whether post hoc explanations on $X\mapsto\hat{Y}$ align with true $X\mapsto Y$ relationship in the data. We first define metrics that capture directional alignment and strength alignment between explanations and ground truth. We then discuss how we simulate datasets with different properties such that the $X\mapsto Y$ relationships can be precisely controlled and systematically varied, enabling a reliable assessment of alignment and the relationships between explanation alignment and data properties. Finally, we evaluate SHAP and LIME explanations using these metrics.

\subsection{Evaluation Metrics for Direction and Strength Alignment}\label{sec:metrics}

We propose a  framework for measuring \textit{direction alignment} and \textit{strength alignment} that unifies SHAP and LIME explanations. The framework is defined at a level of generality that applies to any post hoc explanation method that assigns feature-level attribution values, any predictive model $\mathcal{M}$, and any ground-truth data-generating process $\mathcal{G}$. Method-specific evaluations, such as those for SHAP and LIME used in our experiments, arise as instances of our definition schema.


\subsubsection{Direction Alignment}
Direction alignment evaluates, \textbf{for each feature}, whether the \textit{direction of outcome change implied by an explainer} matches the \textit{true direction of outcome change in the data}. Intuitively, if an explainer implies that increasing feature $j$ increases (or decreases) the model prediction $\hat{Y}$, direction alignment checks whether the same perturbation produces a corresponding directional change in the true outcome $Y$.

The evaluation is performed at the instance level and then aggregated across instances and perturbations to obtain an overall measure of directional consistency.

Let $e_{\mathcal{E}}(\mathbf{x}^{(i)}, j)$ denote the explanation assigned by method $\mathcal{E}$ to feature $j$ for instance $i$. To compare a DGP with what an explainer can be understood as implying about that DGP, we consider perturbations of a single feature while holding all other features fixed.

For an instance $\mathbf{x}^{(i)}$, feature $j$, and perturbation magnitude $\delta$, we define the perturbed instance
\begin{equation}
    \mathbf{x}^{(i)}_{j,\delta} 
    = 
    (x^{(i)}_1,\ldots,x^{(i)}_j+\delta,\ldots,x^{(i)}_d).
\end{equation}
and consider a finite set of allowable perturbations to choose $\delta$ from:
\[
    \Delta = \{m \cdot d \mid m \in \{-M,\ldots,-1,1,\ldots,M\}\},
\]
where $d$ is a step size controlling perturbation magnitude and $M$ determines the extent of the evaluation range. 
The choice of $\Delta$ is independent of the explanation method and defines the evaluation protocol used to compare explainer-implied and true directional responses.


To formalize direction alignment, we decompose the evaluation into three components: 
(i) the direction of change implied by the explainer when a feature is perturbed, 
(ii) the true direction of change induced by the same perturbation under the data-generating process, 
and (iii) an agreement indicator that compares the two and aggregates the comparison across instances and perturbations.



\noindent\textbf{Explainer-indicated direction of change} The explainer-indicated direction of change for feature $j$ is defined as 
\begin{equation}
    \Delta e^{(i)}_{j}(\delta)
    =
    e_{\mathcal{E}}(\mathbf{x}^{(i)}_{j,\delta}, j)
    -
    e_{\mathcal{E}}(\mathbf{x}^{(i)}, j).
\end{equation}
This quantity captures whether the explainer believes the contribution of feature $j$ increases or decreases when $x_j$ is perturbed. This definition applies to any explanation method that produces feature-level attributions and is evaluated over the shared perturbation set $\Delta$.


\noindent\textbf{True direction of change}
The true direction of change induced by perturbing feature $j$ is defined as
\begin{equation}
    \Delta Y^{(i)}_{j}(\delta)
    =
    \mathcal{G}(\mathbf{x}^{(i)}_{j,\delta})
    -
    \mathcal{G}(\mathbf{x}^{(i)}),
\end{equation}
which represents the actual effect of modifying feature $j$ on the outcome generated by the ground-truth process $\mathcal{G}$. Evaluating this quantity requires access to $\mathcal{G}$, allowing us to compute the true outcome at the perturbed input.

\noindent\textbf{Agreement indicator and aggregation}
An explainer is directionally aligned for instance $i$, feature $j$, and perturbation $\delta$ if the explainer-indicated and true changes point in the same direction, i.e.,
\begin{equation}\label{eqn:sign}
    \operatorname{sign}\!\left(\Delta e^{(i)}_{j}(\delta)\right)
    =
    \operatorname{sign}\!\left(\Delta Y^{(i)}_{j}(\delta)\right).
\end{equation}

Requiring strict sign agreement for all perturbations, however, is undesirable when the true effect is extremely small, as the ground truth then provides no reliable directional signal. To avoid penalizing explainers for such cases, we relax the criterion by treating small-magnitude true changes as directionally uninformative.

We define the agreement indicator
\begin{equation}\label{eqn:dir_indicator}
\small
    I^{\text{dir}}_{ij}(\delta)
    \triangleq
    \begin{cases*}
            \mathbb{I}\!\left(
        \Delta e^{(i)}_{j}(\delta)\,\Delta Y^{(i)}_{j}(\delta) > 0
    \right) &if $|\Delta Y_j^{(i)}(\delta)| > \epsilon_1$ \\
    \mathbb{I}\!\left(
       |\Delta e^{(i)}_{j}(\delta)| < \epsilon_2
    \right) & otherwise.
    \end{cases*}    
\end{equation}
which enforces sign agreement when the ground truth exhibits a meaningful directional effect, and requires only a small explainer response when the true effect is negligible.

Aggregating across instances $i$ and perturbations $\delta \in \Delta$, the \emph{direction alignment score} for feature $j$ is defined as
\begin{equation}\label{eqn:direction}
    \rho_{\text{dir},j}
    =
    \frac{1}{|S|\cdot|\Delta|}
    \sum_{i=1}^{|S|}
    \sum_{\delta\in\Delta}
    I^{\text{dir}}_{ij}(\delta),
\end{equation}
where $S$ denotes the evaluation set. The score $\rho_{\text{dir},j} \in [0,1]$ measures how reliably the explainer captures the correct \emph{direction of change} of $Y$ along feature $j$ across the dataset.

\subsubsection{Instantiations: SHAP and LIME}

The above definitions are method-agnostic. We now describe how they are instantiated for SHAP and LIME in our experiments.

\paragraph{SHAP}
For SHAP, $e_{\mathcal{E}}(\mathbf{x}^{(i)}, j)$ corresponds to the Shapley value of feature $j$ for instance $\mathbf{x}^{(i)}$. For each $\delta \in \Delta$, we evaluate the SHAP explanation 
$e_{\mathcal{E}}(\mathbf{x}^{(i)}_{j,\delta}, j)$ and the corresponding true outcome 
$\mathcal{G}(\mathbf{x}^{(i)}_{j,\delta})$, yielding 
$\Delta e^{(i)}_j(\delta)$ and $\Delta Y^{(i)}_j(\delta)$ as defined above.
This instantiation directly tests whether the SHAP interpretation (that larger Shapley values correspond to larger model predictions) extends beyond the model's learned relationship $\mathbf{X} \mapsto \hat{Y}$ to the true data-generating relationship $\mathbf{X} \mapsto Y$.

\paragraph{LIME}
For LIME, the explainer represents the local influence of feature $j$ through the coefficient $w^{(i)}_j$ of a locally fitted surrogate model around instance $\mathbf{x}^{(i)}$. We therefore consider small perturbations $\delta$ so that the evaluation remains within the local neighborhood of $\mathbf{x}^{(i)}$. In this case,
\[
    \Delta e^{(i)}_{j}(\delta) = w^{(i)}_j \cdot \delta,
\]
where $w^{(i)}_j$ is the LIME coefficient for feature $j$ at instance $i$. In our experiments, we set $\delta = 0.1\,\texttt{Std}$.


% \paragraph{SHAP} In the case of SHAP, $e_{\mathcal{E}}(\mathbf{x}^{(i)}, j)$ is the Shapley value for feature $j$ of instance $\mathbf{x}^{(i)}$.  To operationalize the metric, we construct, for each instance $i$ and feature $j$, a \emph{perturbation grid} and choose $\delta$ from it: 
% \[
%     \Delta = \{m\cdot\delta\mid m\in\{-M,\ldots,-1,1,\ldots,M\}\},
% \]
% where $\delta$ is the step size, which can take, for example, the empirical standard deviation of feature $j$.  
% For each perturbation $\delta\in\Delta$, we evaluate the SHAP explanation $e_\mathcal{E}(\mathbf{x}^{(i)}_{j,\delta},j)$ and the true outcome $\mathcal{G}(\mathbf{x}^{(i)}_{j,\delta})$, which directly supply the quantities $\Delta e^{(i)}_j(\delta)$ and $\Delta Y^{(i)}_j(\delta)$ used in Eq.~\eqref{eqn:dir_indicator}. Empirically, $M$ can be set to 1 or a bigger integer.

% Substituting these values into the correctness indicator in Eq.~\eqref{eqn:dir_indicator} determines whether SHAP correctly captures the direction of change for feature $j$ at instance $i$ under perturbation $\delta$. This metric therefore directly tests whether the SHAP interpretation—that larger Shapley values correspond to larger model predictions—extends from the model mapping $\mathbf{X} \mapsto \hat{Y}$ to the true data-generating relationship $\mathbf{X} \mapsto Y$.


% \paragraph{LIME} For LIME, the explainer represents feature $j$'s local influence through the coefficient $w^{(i)}_j$ of the locally fitted surrogate model around instance $\mathbf{x}^{(i)}$. We therefore consider perturbations where $\delta$ is chosen to be sufficiently small, so that the evaluation remains confined to the local neighborhood of the focal point $\mathbf{x}^{(i)}$.  For sufficiently small $\delta$,
% \[
% \Delta e^{(i)}_{j}(\delta) = \omega_j^{(i)}\cdot \delta,
% \]
% where $\omega_j^{(i)}$ is the LIME coefficient for feature $j$ of a focal point $\mathbf{x}^{(i)}_j$. We choose $\delta$ to be $0.1 \texttt{Std}$ in our experiments.


% \noindent\textbf{True direction of change} The true direction of change induced by perturbing feature $j$ is
% \begin{equation}
%             \Delta Y^{(i)}_{j}(\delta)
%         =
%         \mathcal{G}(\mathbf{x}^{(i)}_{j,\delta})
%         -
%         \mathcal{G}(\mathbf{x}^{(i)}),
% \end{equation}
% representing the actual effect of modifying feature $j$ on the true outcome. Thus, evaluating the true direction of change requires access to the ground truth data generation model $\mathcal{G}$, where we can plug in the perturbed data point to compute the true change. 



% An explainer is directionally aligned at instance $i$ for feature $j$ under perturbation $\delta$ if
% \begin{equation}\label{eqn:sign}
%         \operatorname{sign}\!\left(\Delta e^{(i)}_{j}(\delta)\right)
%     =
%     \operatorname{sign}\!\left(\Delta Y^{(i)}_{j}(\delta)\right),
% \end{equation}
% i.e., if both changes point in the same direction.

% A strict interpretation would require this sign agreement for all perturbations. However, such strictness is undesirable when the true effect is extremely small, as the ground truth provides no reliable directional signal. To avoid penalizing explainers for missing negligible effects, we relax the criterion by treating small-magnitude true changes as directionally uninformative.


% We define the agreement indicator
% \begin{equation}\label{eqn:dir_indicator}
%     I^{\text{dir}}_{ij}(\delta)
%     \triangleq
%     \begin{cases*}
%             \mathbb{I}\!\left(
%         \Delta e^{(i)}_{j}(\delta)\,\Delta Y^{(i)}_{j}(\delta) \ge 0
%     \right) & if $|\Delta Y_j^{(i)}(\delta)| > \epsilon_1$ \\
%     \mathbb{I}\!\left(
%        |\Delta e^{(i)}_{j}(\delta)| < \epsilon_2
%     \right) & otherwise.
%     \end{cases*}    
% \end{equation}
% which requires strict sign agreement whenever the ground truth provides a meaningful directional effect, and only requires the explainer’s response to remain small when the true effect is negligible.


% Aggregating across all instances $i$ and allowable perturbations $\delta \in \Delta$, the \emph{direction alignment score} for feature $j$ is
% \begin{equation}\label{eqn:direction}
%     \rho_{\text{dir},j}
% =
% \frac{1}{|S|\cdot|\Delta|}
% \sum_{i=1}^{|S|}
% \sum_{\delta\in\Delta}
% I^{\text{dir}}_{ij}(\delta),
% \end{equation}
% where $S$ is the evaluation set of instances.  
% Thus $\rho_{\text{dir},j}\in[0,1]$ measures how reliably, across the dataset, the explainer captures the correct \emph{direction of change} of $Y$ along feature $j$.


\subsubsection{Strength Alignment}

Strength alignment evaluates whether an explainer correctly captures the relative and undirected importance of features in the data-generating process. Unlike direction alignment, which is assessed at the feature level, strength alignment is a global metric that compares feature-importance rankings aggregated across the dataset.


Let $I_{\mathcal{E}}(j)$ denote the global importance of feature $j$ as implied by explainer $\mathcal{E}$, and let $I_{\mathcal{G}}(j)$ denote the ground-truth importance of feature $j$ under the data-generating process $\mathcal{G}$. We rank features under each scheme to obtain rankings $R_{\mathcal{E}}(j)$ and $R_{\mathcal{G}}(j)$, respectively. The \emph{strength alignment} score is defined as the Spearman Rank Correlation \citep{zar1972significance} between these two rankings:
\[
\rho_{\text{strength}}(R_{\mathcal{E}}, R_{\mathcal{G}})
=
1 - 
\frac{
6\sum_{j=1}^d (R_{\mathcal{E}}(j) - R_{\mathcal{G}}(j))^2
}{
d(d^2 - 1)
},
\]
which takes values in $[-1,1]$. A value close to $1$ indicates strong agreement between the explainer and the ground truth, while values near $0$ or negative indicate weak or systematically incorrect alignment. Because this metric compares importance rankings rather than directional effects, it evaluates exclusively the explainer's ability to recover the \emph{relative} contributions of features in the underlying model.


\paragraph{SHAP}
For SHAP, global feature importance is typically defined as the \textit{mean absolute SHAP value} of each feature across the dataset. To maintain consistency with the unified framework, we define the ground-truth importance $I_{\mathcal{G}}(j)$ using the mean absolute Shapley values computed from the ground-truth model $\mathcal{G}$. Therefore, a strength alignment score for SHAP compares whether the explainer ranks features correctly (based on feature rankings from the true data-generation process). 

\paragraph{LIME}
For LIME, %an explanation for an instance comes from a local linear approximation of $\mM$ with coefficients $w^{(i)}_j$. We define global importance as $I_{\mathcal{E}}(j) = \mathbb{E}_i[\,|w^{(i)}_j|\,]$. 
The ground-truth importance $I_{\mathcal{G}}(j)$ is defined analogously by applying LIME to the ground-truth model $\mathcal{G}$ and aggregating absolute coefficients across instances, which are compared against LIME applied to the model $\mathcal{M}$. Substituting these quantities into the general formulation above yields the strength alignment score for LIME.

\subsection{Data Simulation}
\label{sec:data}

To evaluate how well explanation-based measures align with the true relationship between $\mathbf{X}$ and $Y$, we use simulated data in all experiments.

Each dataset is generated by sampling feature values for 5{,}000 instances and applying a ground-truth data-generating process $\mathcal{G}$ to produce the outcome variable. The process $\mathcal{G}$ combines untransformed feature terms, nonlinear (squared) terms, and pairwise interaction terms in an additive index, with randomly sampled coefficients.

We systematically vary four scenario-defining factors—\texttt{Number of features}, \texttt{Correlation strength}, \texttt{Nonlinear terms}, and \texttt{Interaction terms}—while holding sample size and noise variance fixed. Each factor takes three levels, yielding a total of \textbf{81 datasets}, as summarized in Table~\ref{tab:data_factors}. The complete data-generation procedure is provided in Algorithm~\ref{alg:synth-data} in Appendix~\ref{apx_sec:81_details}.



% To examine how predictive accuracy relates to alignment with the data-generating process, we train multiple XGBoost models for each of the 81 datasets. Model accuracy is systematically varied by training models on progressively smaller subsets of the available training data.%Corresponding to each of the datasets are sets of 4-5 models with test set accuracy generally between 70\% and 90\%. In total, there are 350 models. 




For each dataset, we follow a standard explanation pipeline by first training a predictive model $\mathcal{M}$ on 80\% of the data and evaluating it on the remaining 20\%. The best model is obtained via a systematic search over multiple model families, including gradient-boosted decision trees (e.g., XGBoost, CatBoost, LightGBM, Gradient Boosting) and other commonly used learners (e.g., Random Forests, multilayer perceptrons, and support vector machines). Within each model class, hyperparameters are tuned using cross-validation, and the model achieving the highest validation performance is selected. %Across the 81 datasets, the test accuracies of the resulting best models range from 87.6\% to 95.2\%.



\subsection{Evaluation Results — \textit{Can Post hoc Explainers Recover Feature Direction and Strength from Data}?}

We evaluate SHAP and LIME using the direction and strength alignment metrics defined in Section~\ref{sec:metrics}. For each dataset–model pair, we summarize performance by averaging alignment scores over the five features with the largest mean explanation values, and then examine the distribution of these  alignment scores across all datasets. The probability density functions (PDFs) of direction and strength alignment are shown in Figure~\ref{fig:sec4_best81}. The results yield two main findings.

        \begin{figure}[ht!]
        \centering
  \begin{subfigure}[t]{.35\textwidth}
    \centering
    %\includegraphics[width=\linewidth]{Rebuttal-R1/sec4_shap_dir_best.png}
    \includegraphics[width=\linewidth]{Rebuttal-R1/feng_lu_figures/sec4_best81_shap_dir_v2.pdf}
    \caption{SHAP direction alignment}
    \label{fig:shap_sec4_direction}
  \end{subfigure}
   % \hfill
  \begin{subfigure}[t]{.35\textwidth}
    \centering
    %\includegraphics[width=\linewidth]{Rebuttal-R1/Rebuttal_Figures/shap_filtering_best_sec5_vary_none_strength_ecdf.png}
    \includegraphics[width=\linewidth]{Rebuttal-R1/feng_lu_figures/sec4_best81_shap_strength_v2.pdf}
    \caption{SHAP strength alignment}
    \label{fig:shap_sec4_strength}
  \end{subfigure}
  % \hfill
  % \begin{subfigure}[t]{.45\textwidth}
  %   \centering
  %   \includegraphics[width=\linewidth]{example-image-c}
  %   \caption{SHAP strength recall}
  % \end{subfigure}
  \smallskip

  \begin{subfigure}[t]{.35\textwidth}
    \centering
    \includegraphics[width=\linewidth]{Rebuttal-R1/feng_lu_figures/sec4_best81_lime_dir_step3_v2.pdf}
    \caption{LIME direction alighment}
    \label{fig:lime_sec4_direction}
  \end{subfigure}
   % \hfill
  \begin{subfigure}[t]{.35\textwidth}
    \centering
    %\includegraphics[width=\linewidth]{Rebuttal-R1/Rebuttal_Figures/lime_filtering_best_sec5_vary_none_strength_ecdf.png}
    \includegraphics[width=\linewidth]{Rebuttal-R1/feng_lu_figures/sec4_best81_lime_strength_v2.pdf}
    \caption{LIME strength alignment}
    \label{fig:lime_sec4_strength}
  \end{subfigure}
  % \hfill
  % \begin{subfigure}[t]{.45\textwidth}
  %   \centering
  %   \includegraphics[width=\linewidth]{example-image-c}
  %   \caption{LIME strength recall}
  % \end{subfigure}
  \caption{Evaluation of direction and strength metrics for explanations of the best models for all 81 datasets by SHAP and LIME }\label{fig:sec4_best81}
\end{figure}

\noindent\textbf{Finding 1: High Average Performance.}
Across datasets, post hoc explainers exhibit high average alignment with the ground-truth data-generating process in both direction and strength. In particular, SHAP consistently outperforms LIME, indicating a stronger ability to recover the direction of feature effects on average.


\noindent\textbf{Finding 2: Long Left Tails Undermine Dataset-Level Reliability.}
For both direction and strength alignment, the full distributions exhibit pronounced long left tails, indicating substantial heterogeneity in explainer performance across datasets. Although average alignment is high, a nontrivial subset of dataset–model pairs exhibits markedly poorer alignment.

The nature of these failures differs across explainers and metrics. For SHAP, strength alignment displays a substantially longer and heavier left tail than direction alignment, with values dropping to approximately 0.5 in the worst cases, indicating severe mis-ranking of feature importance for certain datasets. In contrast, for LIME, direction alignment exhibits the more pronounced left tail, with values also falling to approximately 0.5, corresponding to near-random agreement under a binary directional criterion. Strength alignment for LIME, while variable, is comparatively less extreme.

These long left tails imply that high average alignment masks meaningful risks at the dataset level. In particular, explainer outputs can be unreliable for specific dataset–model pairs, even when overall performance appears strong. Because users typically lack reliable ex ante indicators of when an explainer is operating in these low-alignment regimes, the presence of heavy left tails raises concerns about the robustness of post hoc explanations for drawing conclusions about the underlying $X \mapsto Y$ relationship.

\section{What Factors Influence Alignment?}\label{sec:factors}
The preceding section establishes that alignment between post hoc explanations and the true data-generating process varies substantially across dataset–model pairs, even when average performance appears strong. We now turn to understanding the sources of this variation. In particular, we examine which factors systematically influence whether explanations recover the correct direction and strength of feature effects.

Our analysis focuses on three classes of factors. First, we study the predictive performance of the learned model $\mathcal{M}$, which serves as a practical indicator of how closely $\mathcal{M}$ approximates the underlying data-generating process $\mathcal{G}$. Second, we examine the role of model multiplicity through the Rashomon effect, which highlights how multiple, equally accurate models can encode {fundamentally different internal representations} of the same task. Finally, we investigate how intrinsic properties of the data—such as feature correlation, nonlinearity, and interaction structure—shape alignment by constraining or expanding the space of plausible models.
These analyses clarify when alignment failures arise from limitations of the predictive model, when they stem from  ambiguity among high-performing models, and when they are driven by structural properties of the data itself.

\subsection{Predictive performance of $\mathcal{M}$ }
        \label{sec:role_of_predictive_power}
Post hoc explainers such as SHAP and LIME explain the predictive model $\mathcal{M}$, not the underlying data-generating process $\mathcal{G}$ directly. For explanations to reflect the true $\mathbf{X} \mapsto Y$ relationship, the model $\mathcal{M}$ must approximate $\mathcal{G}$ sufficiently well. Predictive accuracy is therefore a necessary condition for alignment: when predictive accuracy is low, $\mathcal{M}$ fails to capture key aspects of $\mathcal{G}$, making alignment with $\mathcal{G}$ unattainable for any post hoc explainer.



To examine how alignment varies with model quality, we construct multiple predictive models for each dataset spanning a wide range of test accuracies. These models are generated by varying both hyperparameter configurations and the size of the training data. We categorize test accuracy into four ranges: $[0.70,0.75)$, $[0.75,0.80)$, $[0.80,0.85)$, and $[0.85,0.90)$. For each accuracy range, we continue generating models until we obtain three models whose test accuracies fall within the corresponding interval. This procedure yields a balanced set of models across accuracy levels for each dataset. We then evaluate direction and strength alignment for each model and report the mean alignment and standard error within each accuracy range. This design allows us to isolate the relationship between predictive performance and explanation alignment while holding the underlying data-generating process fixed.

 \begin{figure}[h]
\centering
\includegraphics[width=0.7\linewidth]{Rebuttal-R1/feng_lu_figures/sec5_impact_model_accuracy_v2.pdf}
\caption{The impact of model accuracy on explanation alignment.}\label{fig:sec5_accuracy}
\end{figure}

Figure~\ref{fig:sec5_accuracy} shows that model accuracy has a statistically significant impact on explanation alignment, particularly for SHAP’s direction alignment. The number reported in the top-right corner of each plot is the coefficient from a univariate OLS regression, where each accuracy bucket is coded as an ordered category (1, 2, 3, and 4). Accordingly, the coefficient represents the average change in the alignment metric associated with a one-bucket increase in accuracy, corresponding to an accuracy improvement of 0.05. For example, a 0.05 increase in test accuracy is associated with an average increase of 0.104 in SHAP’s direction alignment.
           










\subsection{Rashomon Effect}\label{sec:rashomon_effect}

While model accuracy has significant impact, it is \textit{necessary but insufficient} to  guarantee alignment between the learned model $\mathcal{M}$ and the true data-generating process $\mathcal{G}$. This possibly unintuitive limitation arises from the presence of a \emph{Rashomon set}: a collection of models that achieve similarly strong predictive performance on the same task. When a Rashomon set is large, many distinct models map $X \mapsto \hat{Y}$ with comparable accuracy while relying on substantially different internal decision paths or feature relationships. In such settings, the accuracy of $\mathcal{M}$ does not uniquely identify the underlying structure of $\mathcal{G}$. Consequently, even a highly accurate model may yield post hoc explanations that reflect only one of many plausible decision-making mechanisms consistent with the observed data, rather than the true $X \mapsto Y$ relationship. This multiplicity poses a fundamental challenge for inference. When many viable $X \mapsto \hat{Y}$ mappings exist, explanations derived from any single predictive model cannot reliably recover the true $\mathbf{X} \mapsto Y$. As a result, post hoc explanations may vary substantially across models that are  equivalent in terms of predictive performance.

\begin{figure}
\centering
  \begin{subfigure}[t]{.47\textwidth}
    \centering
    %\includegraphics[width=\linewidth]{Rebuttal-R1/Rebuttal_Figures/sec5_vary_nf_shap_dir_best.png}
    \includegraphics[width=\linewidth]{Rebuttal-R1/feng_lu_figures/nf20_cs100_nl5_in0_v2.pdf}
    \caption{The maximum difference of ranks by SHAP for models in a Rashomon Set in one dataset}\label{fig:dataset_ranking}
  \end{subfigure}
\hfill
  \begin{subfigure}[t]{0.47\textwidth}
    \centering
    \includegraphics[width=\linewidth]{Rebuttal-R1/feng_lu_figures/sec5-range_dist_v3.pdf}
    \caption{The distribution of differences of feature ranks}
  \end{subfigure}
  \caption{The evidence of Rashomon Effect on Post hoc explanations}\label{fig:range_dist}
\end{figure}

To empirically illustrate explanation multiplicity induced by the Rashomon effect, we obtained 10 best-performing models for each of our 81 datasets. For each dataset, the selected models achieve test accuracies within 3\% of one another. We then applied SHAP to each model and ranked features according to their mean absolute Shapley values. This procedure yields 10 feature-importance rankings per dataset, one for each predictive model. 

We plot in Figure \ref{fig:dataset_ranking} the range of feature rankings for a representative dataset. Each line spans from the minimum to the maximum rank assigned to a feature across the 10 models, with features ordered by their median rank. If the models’ internal decision-making processes were highly consistent, we would expect minimal variation across these rankings. However, Figure~\ref{fig:dataset_ranking} shows that while many features exhibit modest variability, others show substantial disagreement. For instance, in the demo dataset we picked from the 81 datasets, Feature~0 is ranked as both the most important feature (rank 1) and relatively unimportant (rank 11 out of 20) by different models, despite all models achieving nearly identical predictive accuracy.

To provide a holistic view of this phenomenon, we compute the ranking variation for all features across all datasets and visualize the resulting distribution. To ensure comparability across datasets with different numbers of features, we normalize the ranking range by the total number of features, yielding values between 0 and 1. Figure~\ref{fig:range_dist} shows that, on average, feature rankings vary by approximately 10\%, with a pronounced right tail. This indicates that for a nontrivial subset of features, ranking disagreement can be substantial, further underscoring the instability of post hoc explanations within large Rashomon sets.
As a consequence, high accuracy alone cannot guarantee alignment between explanations and the true data-generating mechanism. The Rashomon effect ensures that this misalignment persists even in best-case predictive scenarios, as multiple distinct decision rules remain equally compatible with the observed data.

 



Nevertheless, \emph{high predictive accuracy is necessary while not sufficient for alignment}. While high accuracy does not imply that $\mathcal{M}$ aligns with $\mathcal{G}$, low accuracy provides strong evidence that $\mathcal{M}$ fails to represent $\mathcal{G}$, and therefore that explanations of $\mathcal{M}$ cannot align with the truth. This asymmetry motivates our investigation into how varying levels of model accuracy influence alignment. 


    
\subsection{Data Properties}

In addition to predictive performance and model multiplicity, we examine how intrinsic properties of the data influence explanation alignment. Data properties are directly observable characteristics that empirically predict when and how severely explanation misalignment occurs.

To visualize the effect, we plot the direction and strength alignment for datasets with different factors, and the coefficients in the OLS regression. Figure~\ref{fig:alignment_data} shows that feature correlation strength, the number of nonlinear terms, and the number of interaction terms are statistically associated with reduced alignment. %However, their practical impact differs substantially. Among these factors, correlation strength exhibits a markedly larger and more consistent negative effect, particularly on strength alignment.

\begin{figure}[h]
\centering
  \begin{subfigure}[t]{.95\textwidth}
    \centering
    \includegraphics[width=\linewidth]{Rebuttal-R1/feng_lu_figures/sec5_SHAP_DGP_factors_v3.pdf}
    \caption{Impact of data characteristics on SHAP alignment}
  \end{subfigure}

  \begin{subfigure}[t]{0.95\textwidth}
    \centering
    \includegraphics[width=\linewidth]{Rebuttal-R1/feng_lu_figures/sec5_LIME_DGP_factors_v3.pdf}
    \caption{Impact of data characteristics on LIME alignment}
  \end{subfigure}
  \caption{Impact of data characteristics on post hoc explanation alignment. Numbers on the top right are the coefficients of the univariate OLS analyses.}
  \label{fig:alignment_data}
\end{figure}



When features are highly correlated, multiple variables can act as substitutes or proxies for one another, enabling models with substantially different internal attributions to achieve comparable predictive accuracy. Empirically, this is associated with greater disagreement in post hoc explanations. A plausible explanation is that feature substitutability increases the number of near-optimal models consistent with the data, making the learned model’s internal attribution more sensitive to arbitrary modeling choices.

A similar mechanism applies to the number of nonlinear terms and the number of interaction terms. As the data-generating process becomes more complex, a larger set of functional forms can fit the observed data equally well, increasing model flexibility without a commensurate increase in identifiability. As a result, learned models that achieve comparable predictive performance may rely on qualitatively different internal representations, leading to greater variability in post hoc explanations.

Overall, these results suggest that explanation misalignment is not driven solely by explainer choice or model performance, but is systematically related to intrinsic properties of the data. In particular, feature correlation, nonlinearities, and interactions amplify ambiguity in how predictive relationships are encoded, making reliable recovery of the underlying $X \mapsto Y$ relationship more difficult even when predictive accuracy is high.

One counterintuitive observation is that the number of features is associated with a small positive effect on explanation alignment. However, the magnitude of this effect is negligible (below 1\%) and substantially smaller than the effect of feature correlation strength. We therefore do not interpret this result as evidence that higher dimensionality improves explanation faithfulness. %Instead, it likely reflects metric- or explainer-specific artifacts rather than a meaningful structural phenomenon.

\section{Robustness Checks}\label{sec:robustness}

The preceding results show that data misalignment in post hoc explanations arises systematically rather than as an artifact of a particular modeling choice. We now examine whether this conclusion depends on specific implementation or evaluation decisions, or whether misalignment persists even under conditions designed to be favorable to alignment. To do so, we conduct three robustness checks that probe aggregation across models, sensitivity to explainer hyperparameters, and sensitivity to evaluation metrics.

Each check is motivated by a natural stabilizing intuition: that averaging across equally accurate models may recover a more faithful signal; that careful explainer tuning may reduce approximation noise; and that alternative metric specifications may yield different conclusions. Across all settings, we find that these intuitions fail to overturn our main result. Even under favorable configurations, data misalignment remains persistent. Full experimental details are reported in Appendix~\ref{apx_sec:robustness_results}.

\subsection*{Robustness Check 1: Averaging over a Rashomon Set}

The Rashomon effect implies that many distinct models can achieve near-identical predictive performance while relying on different internal decision logic. This raises the possibility that while any single model’s explanation may be misaligned, averaging explanations across multiple equally accurate models could yield a more faithful representation of the data-generating process.

To test this idea, we construct, for each dataset, a sample of a Rashomon set \footnote{Note that a Rashomon set could contain tens of thousands of models. Here we construct only a small set to illustrate the idea.} consisting of the ten highest-performing models drawn from diverse model families. All selected models lie within a 3\% absolute accuracy margin of the best-performing model. For each test instance, we average instance–feature attributions across models and evaluate alignment using the direction and strength metrics defined in Section~\ref{sec:main_alignment}. This test assesses whether consensus across plausible models mitigates misalignment.

\subsection*{Robustness Check 2: Sensitivity to Explainer Hyperparameters}
\label{sec:lime_params}

Post hoc explainers introduce additional layers of approximation of $\mG$ (beyond that of $\mM$) through sampling schemes, surrogate models, and algorithmic hyperparameters. We examine whether misalignment is driven by such explainer-induced noise.
For SHAP, we vary the number of bootstrap samples. For LIME, we vary kernel bandwidths and neighborhood sizes, motivated by theorems in Appendix \ref{apx_sec:lime_thms}. Because LIME relies on randomized perturbations, we also average explanations across multiple random seeds while holding all other parameters fixed. This analysis isolates whether stabilizing explainer-level variability improves alignment independently of model behavior.

\subsection*{Robustness Check 3: Sensitivity to Metric Parameters}

Our direction alignment metric depends on a perturbation magnitude $\delta$ that determines the scope of feature perturbations. Smaller perturbations emphasize local behavior, while larger perturbations probe broader directional effects. We vary $\delta$ while holding all other components of the evaluation fixed and recompute alignment scores. This test assesses whether our conclusions hinge on a particular perturbation scale or persist across reasonable metric specifications.

% \section{Robustness Checks}\label{sec:robustness}

% So far, we have demonstrated the phenomenon of data misalignment in post hoc explanations of black-box models and provided insights into why such misalignment arises. In this section, we conduct a set of robustness checks to examine how our findings behave under a range of intuitive and practically motivated variations in the explanation and evaluation pipeline. Some of these checks are designed to explore whether straightforward strategies can improve alignment, while others assess the sensitivity of our conclusions to reasonable modeling and evaluation choices.

% These robustness checks are guided by three considerations. First, aggregating explanations across a collection of models or runs may yield explanations that are more aligned with the underlying data-generating process than any single explanation. Second, certain choices of LIME hyperparameters can improve alignment, as theoretically motivated and formally justified in Appendix \ref{apx_sec:lime_thms}. Third, we vary hyperparameters in our evaluation metrics to assess whether our conclusions depend on particular metric specifications.

% Across all settings, these robustness checks provide a more comprehensive view of the conditions under which misalignment arises. Importantly, even when employing strategies that are favorable to improved alignment and when varying evaluation choices, our main conclusions remain unchanged: data misalignment persists and continues to pose a fundamental challenge for post hoc explanations of black-box models. Due to the space limit, we summarize the main proposal and findings and leave all experiments results in Appendix \ref{apx_sec:robustness_results}.

% \subsection*{Robustness Check 1 - Averaging Explanations within a Rashomon Set}
% Section~\ref{sec:rashomon_effect} demonstrated that the Rashomon effect—the existence of many distinct models with near-identical predictive performance —naturally gives rise to a large family of plausible explanations for any given data sample. This raises a natural question: while an explanation derived from a single trained model $\mathcal{M}$ may be misaligned with the data-generating process, could averaging explanations over a set of highly accurate, yet distinct, models mitigate this issue?

% This idea is motivated by a classical modeling philosophy widely used in machine learning, in which aggregating multiple expert opinions (e.g., ensemble models \citep{ho1995random, dietterich2000ensemble}, mixtures of experts \citep{yuksel2012twenty}) often leads to more stable and reliable predictions. Analogously, when models are sampled from a Rashomon set, each capturing a distinct mapping $\bx \mapsto \hat{y}$, averaging their explanations may capture a form of consensus among equally plausible hypotheses for the mapping $\bx \mapsto y$, potentially improving alignment with the true data-generating process.


% For each dataset, we first train a diverse collection of models via grid search over hyperparameters across multiple model families, including gradient-boosted decision trees (e.g., XGBoost, CatBoost, LightGBM, GradientBoost) and other commonly used learners (e.g., RandomForest, MLP, SVM, DecisionTree). From this pool, we select the 10 best-performing models according to test accuracy. Empirically, the predictive performance of these models lies within a 3\% absolute accuracy margin of the highest-performing model on each dataset.
% This procedure yields a sample from the Rashomon set of 10 distinct, high-performing models for each dataset, all of which are nearly indistinguishable in terms of predictive accuracy yet may implement different internal representations and decision logic.
%  Given a test-set instance from a dataset, we generate explanations for each model and then take the average to obtain average instance-feature level explanation values. We then follow the definitions of direction and strength alignment metrics in Section 4 to compute the final alignment scores.

   


    
% \subsection*{Robustness Check 2: Sensitivity to Explainer Hyperparameters}
% \label{sec:lime_params}

% Even when the predictive model is held fixed, post hoc explainers like SHAP and LIME provide only approximate characterizations of model behavior and can introduce additional variability through algorithmic choices and stochastic components. This robustness check examines whether such explainer-induced variability affects explanation alignment.

% For SHAP, we vary the number of bootstrapped samples used to estimate expected values. For LIME, we vary the kernel bandwidth and the number of data points sampled in the local neighborhood of each focal instance. In addition, because LIME relies on randomized perturbations to fit local surrogate models, we generate explanations for each test instance using multiple random seeds while holding all other explainer hyperparameters fixed. We then average the resulting instance–feature attribution values across seeds to obtain a stabilized explanation. This analysis assesses whether stabilizing explainer-level approximation noise—independent of model behavior—meaningfully influences explanation alignment.



%     \subsection*{Robustness Check 3 - Sensitivity to Metric Parameters}
%        The direction alignment metric evaluates agreement between explainer-implied and true outcome directions by comparing changes induced by feature perturbations. As with any perturbation-based evaluation, this procedure depends on how the perturbations are constructed. In particular, the metric requires specifying a perturbation magnitude 
% $\delta$, which determines the set of perturbed instances $\Delta$ over which directional agreement is assessed.

% This robustness check examines how varying the perturbation extent affects the measured alignment. Intuitively, smaller perturbations emphasize highly local behavior, while larger perturbations probe directional effects over a broader range of feature values. Changing the step size therefore alters the perturbation set $\Delta$ and may influence the resulting alignment scores.

% To assess sensitivity, we vary the perturbation magnitude used in the direction alignment metric while holding all other components of the evaluation fixed. We then recompute direction alignment scores under these alternative settings. This allows us to evaluate whether our conclusions depend on a particular choice of perturbation scale.

\subsection*{Results}
\label{sec:summary_and_strategies}

Plots summarizing the results of our robustness checks are reported in Appendix~\ref{apx_sec:robustness_results}; here we highlight the main findings. Consistent with our primary analyses, we continue to observe persistent misalignment between post hoc explanations and the ground-truth data-generating process. In particular, alignment distributions exhibit pronounced left tails, indicating that severe failures in recovering both direction and strength occur with nontrivial frequency.

First, varying the perturbation step size in $\Delta$ does not materially alter the shape of the alignment distributions, suggesting that our conclusions are robust to reasonable choices of this metric parameter.
Second, certain aggregation strategies yield modest improvements in alignment, though none eliminate misalignment entirely. Averaging explanations across models within a Rashomon set improves direction recovery for SHAP, consistent with recent findings \citep{donnelly2023rashomon}. Similarly, averaging LIME explanations across multiple runs with different random seeds improves both direction and strength alignment.
Third, improved parameter choices within the explainer can lead to small gains in alignment but again do not resolve misalignment. In particular, the choice of kernel bandwidth has a statistically significant effect on alignment outcomes.

Overall, these results reinforce our central conclusion: while aggregation and careful parameter tuning can slightly improve alignment, post hoc explanations remain unreliable for inferring the true $\mathbf{X} \mapsto Y$ relationship given the long left tail in the distribution. Therefore, even under favorable conditions, explanation outputs should be interpreted with caution.


\section{Diagnosing Explanation Reliability via Rashomon Sets}\label{sec:rashomon_agreement}

The results in the previous sections demonstrate that post hoc explanations are not uniformly reliable in recovering true effects in data: while they perform well on average, they can fail in a non-negligible fraction of datasets. This raises a natural and practically important question: \textit{when should a post hoc explanation be trusted, and when should it be treated with caution}? In this section, we explore whether disagreement among equally high-performing models in a Rashomon set can serve as a signal of the risk associated with trusting a given post hoc explanation.


\subsection{Our Hypothesis: Rashomon Agreement Signals Explanation Alignment}
As discussed previously, the models comprising a Rashomon set achieve highly similar predictive performance on the observed data, making them effectively indistinguishable under standard evaluation metrics. Despite this similarity in accuracy, these models may rely on substantially different internal representations and decision pathways to map $X$ to $\hat{Y}$.  Figure \ref{fig:range_dist} illustrates such multiplicity reflected as  the disagreement in post hoc explanations for models in a Rashomon Set.

We refer to the degree of consistency among models in a Rashomon set as \textit{Rashomon agreement}. We hypothesize that, if models in a Rashomon set  exhibit low agreement in the sense that they disagree in their predictions or explanations, this disagreement provides evidence that \textbf{the space of hypotheses compatible with the data is relatively large}. In such settings, any single model represents only one of many plausible explanations for how $X \mapsto Y$, and the likelihood that its post hoc explanation aligns closely with the true data-generating process $\mathcal{G}$ is correspondingly reduced. Conversely, high Rashomon agreement among models with similar predictive performance suggests a more constrained hypothesis space, increasing confidence that the explanations produced by any one model reflect stable, data-supported structure rather than artifacts of model choice.


\subsection{Measuring Rashomon Agreement}
% Motivated by this intuition, we define and study two forms of \textbf{agreement} within a Rashomon set as indicators of the reliability of post-hoc explanations, based on model predictions and explanations, respectively. Agreement is computed pairwise between models in the Rashomon set and then averaged to obtain a dataset-level agreement score.


To operationalize Rashomon agreement, we must confront the fact that agreement among models in a Rashomon set cannot be observed directly: the models may differ in architecture, parameterization, or internal representations, making a direct comparison of their learned functions infeasible. Instead, we assess agreement through observable manifestations of the underlying $\mathbf{X} \mapsto Y$  mapping. In particular, agreement can be revealed through (i) the models’ predictions, which reflect their externally observable behavior, and (ii) their post hoc explanations, which provide a window into how each model attributes importance to features when producing those predictions. We therefore define two complementary agreement metrics—one based on predictions and one based on explanations—to capture these distinct but related notions of consistency within a Rashomon set. Agreement is computed pairwise between models in the Rashomon set and then averaged over all pairs to obtain a dataset-level agreement score.

\noindent\textbf{Prediction Agreement}
We begin with agreement metrics based on model predictions, since predictions are the primary externally observable outputs of a model and the basis on which models in a Rashomon set are deemed equally accurate. If two models implement similar $\mathbf{X}$ to $\hat{Y}$ mappings, this similarity should be reflected in how often they produce the same predicted labels on the same instances. Given two models, we define the prediction agreement metric as \textit{the fraction of instances on which two models agree}. Larger values correspond to greater agreement. 

\noindent\textbf{Explanation Agreement}
While prediction agreement captures similarity in observable behavior, it does not reveal whether models rely on similar feature attributions to generate those predictions. To assess agreement in the inferred structure of the $\mathbf{X}\mapsto Y$ mapping, we next define agreement metrics based on post hoc explanations.

Because explanation values are real-valued, and can vary substantially in scale across models and explainers, we focus on rank-based measures that are invariant to monotonic transformations. Specifically, for each model we rank features based on the mean absolute explanation values across instances, and define explanation agreement between two models as the Spearman correlation between these vectors. 

% Then, for a dataset $\mathcal{D}$, the mean agreement with respect to an agreement metric  is calculated as the average of the 45 pairwise agreement values across all pairs of models in the Rashomon set. 

\paragraph{Results}


If these agreement metrics are informative, we would expect higher agreement among models in a Rashomon set to correspond to stronger alignment between post hoc explanations and the data-generating process. In other words, Rashomon agreement should be positively correlated with explanation alignment for a given model trained on a dataset.


To evaluate this hypothesis, we examine the correlation between each Rashomon agreement metric and the explanation alignment metrics, along both the direction and strength dimensions. The results are shown in Table \ref{tab:rashomon_agreement_corr}. 


\begin{table}[ht]
\centering
\small
\caption{Correlation between Rashomon agreement and explanation alignment}
\label{tab:rashomon_agreement_corr}
\begin{tabular}{l l c c}
\toprule
\textbf{Explainer} & \textbf{Agreement Type} & \textbf{Direction Alignment} & \textbf{Strength Alignment} \\
\midrule
\multirow{2}{*}{SHAP} 
    & Explanation agreement & 0.326 & \textbf{0.792} \\
    & Prediction agreement  & 0.344 & 0.088  \\
\midrule
\multirow{2}{*}{LIME} 
    & Explanation agreement & 0.217 & \textbf{0.695} \\
    & Prediction agreement  & 0.131 & 0.023  \\
\bottomrule
\end{tabular}
\end{table}
\normalsize




Table \ref{tab:rashomon_agreement_corr} reveals two key findings. First, both Rashomon agreement metrics are positively correlated with explanation alignment, indicating that higher agreement among equally accurate models is indeed associated with closer alignment between post hoc explanations and the data-generating structure. This relationship is particularly strong for explanation-based agreement. For SHAP, explanation agreement correlates with alignment by up to approximately 0.792, while for LIME the corresponding correlations reach roughly 0.695. These high correlations suggest that Rashomon agreement provides a strong signal of explanation reliability in recovering true effects in the data.

Moreover, agreement based on explanations is consistently more informative than agreement based on predictions. This difference is expected. Prediction-based agreement captures similarity in model outputs, which can remain high even when models rely on different internal representations or feature attributions to arrive at the same predictions. In contrast, explanation-based agreement directly measures consistency in how models attribute importance to features. As a result, it is a more sensitive indicator of whether explanations reflect stable, data-supported structure rather than model-specific artifacts.


Table~\ref{tab:rashomon_agreement_corr} also highlights an important asymmetry between SHAP and LIME: alignment for SHAP is more strongly correlated with Rashomon agreement than alignment for LIME.  This suggests that SHAP alignment is more predictable from agreement-based signals among equally accurate models, whereas LIME’s alignment is noisier and harder to anticipate, even when models exhibit substantial agreement.




\subsection{Dataset-Level Limits of Post hoc Explanations}


To understand the \textit{dataset-level origins} of the disagreement, we examine how intrinsic data properties—such as dimensionality, feature correlation, nonlinearity, and feature interactions—affect both prediction and explanation agreement. We compute the average agreement for datasets with varying characteristics and fit ordinary least squares (OLS) regressions to quantify the contribution of each factor.

The results for prediction agreement, shown in Fig.~\ref{fig:prediction_agreement_data}, indicate that several data characteristics have a strong and statistically significant impact. The number of interaction terms is the most influential factor: increasing the number of interactions by one is associated with a 16.9\% decrease in prediction agreement. The number of nonlinear terms also has a substantial effect, with an estimated decrease of 12.4\% per additional nonlinear term. These findings suggest that certain forms of data complexity systematically weaken the extent to which accuracy constrains model behavior.



\begin{figure}[H]
    \centering
    \includegraphics[width=\linewidth]{Rebuttal-R1/feng_lu_figures/sec7_prediction_agreement_v2.pdf}
    \caption{The impact of data properties on prediction agreement.}
    \label{fig:prediction_agreement_data}
\end{figure}


We observe a qualitatively different pattern for explanation agreement. While explanation agreement is strongly correlated with explanation alignment, its relationship with dataset characteristics is weaker and less systematic. Among the factors considered, only feature correlation strength exhibits a statistically significant effect, and this effect is negative but small in magnitude. The remaining factors—dimensionality, nonlinearity, and interaction terms—have no consistent or significant impact. This contrast highlights that explanation agreement is not determined solely by dataset complexity, but also reflects how a particular explainer captures model behavior.

\begin{figure}[h]
    \centering
    \includegraphics[width=\linewidth]{Rebuttal-R1/feng_lu_figures/sec7_explanation_agreement_data_v3.pdf}
    \caption{The impact of data properties on explanation agreement}
    \label{fig:explanation_agreement_data}
\end{figure}


\subsection{Discussion}

% \textcolor{red}{edit}

The results in this section support a central diagnostic message: predictive accuracy is necessary but insufficient for post hoc explanations to reliably recover the true direction and strength of the $\mathbf{X} \mapsto Y$ relationship. When multiple models achieve comparable accuracy yet rely on different internal representations, post hoc explanations are inherently less reliable indicators of the data-generating structure. Rashomon agreement provides a practical way to detect precisely settings where accuracy alone fails to meaningfully constrain the underlying $\mathbf{X}\mapsto Y$ mapping.

From a dataset-level perspective, disagreement among models arises naturally when the observed data admit multiple predictive hypotheses. As data become more complex—through nonlinear relationships, feature interactions, or correlated inputs—there are more distinct ways for models to combine features to achieve comparable predictive performance. In such settings, accuracy does not uniquely pin down how predictive information is represented. Models may rely on different subsets of features, encode interactions differently, or resolve trade-offs among correlated variables in alternative ways, all while achieving nearly identical performance on standard metrics.

Crucially, this ambiguity is not merely theoretical: it is directly observable through low Rashomon agreement. When agreement is low, the space of hypotheses consistent with the data is large, and any single model’s post hoc explanation reflects only one of many plausible explanations. In these cases, misalignment between explanations and the true data-generating process is not a failure of the explainer per se, but a consequence of insufficient \textit{identifiability} in the data. As a result, explanations produced in low-agreement regimes should be interpreted with caution, regardless of a model’s predictive accuracy.

Conversely, high Rashomon agreement indicates that accuracy tightly constrains model behavior. When equally accurate models converge in their predictions and—especially—their feature attributions, the hypothesis space consistent with the data is more restricted. In these regimes, post hoc explanations are substantially more stable and more likely to align with true underlying effects. This distinction clarifies why explanation-based agreement is a particularly informative diagnostic: it directly probes consistency in the mechanisms that explainers are designed to recover, rather than surface-level agreement in outputs.

Overall, this section reframes explanation reliability as a dataset-level property rather than a property of a specific model or explainer. Rashomon agreement does not guarantee that an explanation is correct, but low agreement provides a clear and actionable warning sign: when many equally accurate models tell different stories, none of them should be taken at face value. This diagnostic perspective complements existing evaluations of post hoc explanations and provides practitioners with a principled way to identify when explanations are most likely to \emph{mislead rather than inform}.  


        

\section{Conclusion}
    \label{sec:discussion}
    As post hoc explanations are gaining popularity in business research, we found an alarming trend of using them to make claims about the direction and strength of variables at the data level rather than at the model level.  However, the validity of this use of post hoc explanations has not been studied in the literature. 
In this paper, we systematically evaluate this practice.
     In the following sections, we will summarize our findings and give some parting advice for the use of post hoc explainers for obtaining data insights in business research. 
    % We also contribute to the business literature by proposing some simple and intuitive robustness checks which show various ways in which both stages of the pipeline are sensitive to inputs and parameters. 



    \subsection{Main Findings}
        We contribute several important findings to the use of post hoc explanations in business research.

            First, through an extensive literature review, we found that the use of post hoc explanations to infer information about the data is quite prevalent. %In the 181 papers we reviewed, 80 used SHAP or LIME to draw conclusions about the $X\mapsto Y$ relationship instead of the $X\mapsto \hat{Y}$ relationship. 
            Such use deviates from the intended purpose of post hoc explainers; they are designed to explain the machine learning model. Identifying this misuse of post hoc explanations is an important contribution of this paper. %Note that while post hoc explanations have been first introduced and widely used in the machine learning community, the issue we discuss in this paper is less of a concern for that community because computer scientists are more interested in explaining the machine learning models. They use post hoc explainers to understand a model's behavior in order to diagnose and debug the model. However, the goal of business researchers is to understand the data and the underlying phenomena, making this issue more pertinent in business research.


Second, we show that SHAP and LIME explanations do not reliably recover the direction or relative strength of true feature effects, even when predictive models achieve high accuracy. %Using simulated data with known ground truth, we demonstrate that explanation misalignment arises even under highly favorable conditions, including linear data-generating processes, independent features, complete observability, and strong predictive performance. These results show that high accuracy is necessary but insufficient for reliable explanation: equally accurate models can produce explanations that disagree substantially with the true $X \mapsto Y$ relationships.

Third, we identify the fundamental sources of explanation misalignment and show that these failures are structural rather than incidental. At a conceptual level, misalignment is driven by the Rashomon effect: the existence of many distinct models with near-identical predictive performance but different internal representations and feature attributions. At the data level, misalignment arises from the complexity of the data-generating process, including feature correlation, nonlinearities, and interaction effects, which together admit multiple substitutable predictive explanations. These forces imply that even in best-case predictive settings, post hoc explanations reflect only one of many plausible hypotheses consistent with the data, rather than a uniquely supported account of the data-generating process.


        
Fourth, we show that disagreement among equally accurate models provides a powerful diagnostic of explanation unreliability, while common robustness strategies offer only limited mitigation. In particular, agreement in post hoc explanations across models within a Rashomon set is strongly correlated with alignment to the ground truth, especially for feature-importance rankings. Conversely, low agreement reliably signals settings in which explanations are most likely to mislead. Although averaging explanations across models or explainer runs can modestly improve alignment, such strategies do not eliminate long-tailed failures and cannot guarantee reliable recovery of true feature effects.

 
        


        
        
        % \textcolor{red}{what else did we *find*?}

            
        % Second, users must deal with inconsistent explanations because not only do e.g. SHAP and LIME each produce different explanations for the same outcome, but LIME may produce self-contradictory explanations of the same outcome. How to select an explanation(s) from the myriad available from different combinations of accurate black box models, explainers, and explainer hyperparameters is up to the user, but the mitigations suggested in this paper can help users make more informed decisions.%but robustness checks like those presented in this paper can be used to estimate how confident users can be in their selection(s). 



        % stuff re: prescriptive removed
        %Our main finding regarding prescriptive insights is that LIME, SHAP, and DiCE seem to only be of use for actionable recourse a minority of the time. When tasked with generating new application profiles for denied loan users, those new application profiles would be downloaded by $\mathcal{G}$ less than 50\% of the time. %\textcolor{blue}{May need to be removed!! We showed that... DiCE, a counterfactual explainer (as opposed to an attributive one like SHAP and LIME) was always able to produce new application profiles for ignored loan users that would lead to approval in our experiments (for a sufficiently small sparsity parameter). }For DiCE, if we restrict the range of feature variations (within one standard deviation) and the number of features allowed to vary, it can only provide prescriptive insights to a small number of denied users.
        
    % \subsection{On robustness checks}
    %     We proposed some intuitive robustness checks for the issues of unfaithful and inconsistent explanations. Each of them was effective with respect to at least one of our three metrics. However, none of them were fully effective solutions that would guarantee that practitioners could confidently use LIME or SHAP for $\boldsymbol{\beta}$ estimation and econometric modeling. Thus, we encourage cautionary usage of SHAP and LIME and to consider exploring the usage of our robustness checks in efforts to obtain more data-aligned explanations. 
        
    %     Our principal finding regarding faithfulness was that SHAP has no chance at data-alignment unless its Shapley values are normalized by the difference between feature values and feature means. We also found that the most important strategy for increasing the faithfulness of LIME is to disable discretization of numerical features and to make the number of samples and kernel bandwidth it uses for weighted ridge regression large. 
        
    %     %We found some success in increasing the faithfulness of an explanation of a model $\mathcal{M}$ to the ground truth $\mathcal{G}$. As expected, whether $\mathcal{M}$ and $\mathcal{E}$ can correctly predict decisions from $\mathcal{G}$ affects how faithful their explanations are to $\mathcal{G}$. In addition to ensuring low generalization error of $\mathcal{M}$, we also noted that ensuring that $\mathcal{M}$'s training data without correlated or confounded features resulted in remarkably better results than otherwise. However, it is important to note that the explanations were still usually imperfect even though the model and data were ideal. 
    
    %     To combat inconsistency, we proposed techniques to construct aggregate explanations that are better on average than individual explanations. These strategies worked by  averaging feature importance values obtained from explanations obtained from different combinations of models and explainers. The motivation behind these averaging strategies was that they address uncertainty in the choices users need to make by aggregating the many combinations of models and explainers that are the root of the uncertainty. In addition to addressing inconsistency, these strategies resulted in statistically significant improvements in at least directionality, concordance, or relevance. 

    %     In exploring robustness checks, we did not find evidence that LIME or SHAP can be used in a way that makes them suitable alternatives to traditional econometric modeling. Although our recommended strategies for SHAP and LIME compared favorably to their associated baselines in a test to reproduce results an actual econometric study of loan approval, they were unable to consistently extract the $\boldsymbol{\beta}$ that the econometricians found. 

%\textcolor{red}{add in the discussion that althgouh we do not know when the explanation aligns with beta, we can identify when they do not: just run  pipeline multiple times with different split of data and different ML model, if the results do not agree, it is a clear sign that the beta-alignment is low}
    
    
  %   \subsection{On the Use of Post Hoc Explanations}
  % We want to emphasize that this paper focuses on the issue of using post hoc explanations to make inferences about the data. We do not explore their primary/intended usage in explaining the machine learning model $\mM$ or its predictions $\mathcal{M}(X)$. However, it is important to note that even this intended use of post hoc explainers is controversial \citep{garreau2020looking,alvarez18robustness,kumar20shap,jacovi2023diagnosing}. For the scope of this paper, our discussion is solely focused on the idea of inferring ground truth information from $\mathcal{G}$ through the explanations for $\mathcal{M}$. 

% \subsubsection{How does SHAP Compare to Traditional Regression Methods?}
%     In our literature review, we found that researchers were motivated to use SHAP explanations of an ML model to obtain feature importance values because of perceived advantages XAI has over regression analysis. The most commonly cited advantage is that, unlike traditional least squares regression methods, SHAP explanations automatically include nonlinearities and interactions \cite{hanauer22boosting,berger2023supply,srivastava23energy,banulescu2023practical,li24roles}. For example, \cite{berger2023supply} describes their motivation to use SHAP by asserting that SHAP can provide similar economic insights to pooled OLS \emph{and} capture the non-linear impacts of features. 
    
%     %%\cite{srivastava23energy}, which expresses the desire to use and explain ML models with SHAP in order to find factors driving energy consumption because ML models can ``detect highly complex interactions between drivers of energy consumption," is an example of this. Furthermore, \cite{zhou23consumption} says that XAI methodology is appealing because it does not necessitate ``...the construction of more intermediate variables and complex statistical analysis." Similarly, \cite{banulescu2023practical} are drawn away from logistic regression because it suffers when a large number of predictors are used. \cite{hanauer22boosting} choose not to use traditional linear regression analysis used for fundamental stock analysis and instead explain random forest and XGBoost models because the ML models deal with non-linearities, interactions, and noisy features. \cite{li24roles} makes inferences about bike sharing by explaining their XGBoost model instead of examining the coefficients of their linear regression because the XGBoost had better predictive performance. 

%      The fact that SHAP is a local explainer makes its use for obtaining global feature importance unnatural, and this is a shortcoming that SHAP has in comparison to regression modeling. We found that authors often circumvented SHAP's status as a local explainer by computing the average of the absolute Shapley values over their test set and using the results as measures of global feature importance. \cite{lenaers23rent,maeng23vr,guenther2023useful} are all examples. This obviously comes at the cost of losing any indication of the directionality of impact by features. However, they must do this because the means of signed Shapley values should be 0 due to Equation \ref{eqn:shap_gt}.\footnote{This is usually the case empirically; \cite{lenaers23rent,maeng23vr,jiang24churn} all show scatter plots of the signed Shapley values of each feature over their dataset and the empirical distributions of the Shapley values are almost always centered at 0.} Equation \ref{eqn:shap_gt} implies that the expected absolute Shapley value for feature $X_j$ is proportional to $\beta_j$, but also a function of the mean and standard deviation of $X_j$. So, unless all of the $X_j$ are identically distributed, the expected absolute Shapley values may not be comparable for the sake of relative importance in the same way as the $\beta_j$. In other words, a ranking of the average absolute Shapley values may be different than the ranking of the $\mid \beta_j\mid$. 
     
     
     %For example, if $\beta_1=1$ and $\beta_2=2$ with $X_1\sim\mathcal{N}(\mu_1=100,\sigma_1=1)$ and $X_2\sim \mathcal{N}(\mu_2=0,\sigma_2=\frac{1}{10})$ that are viewed as important in the sense of the effect of a unit increase, the expected absolute Shapley value for $X_1$ is larger than the expected absolute Shapley value of $X_2$ even though $X_2$ is more important than $X_1$ because $\beta_2>\beta_1$:
    % \begin{equation*}
    %     \mathds{E}[\mid \beta_1(X_1-\mu_1)\mid] = \sqrt{\frac{2}{\pi}} > \mathds{E}[\mid \beta_2(X_2-\mu_2)] = \frac{1}{5}\sqrt{\frac{2}{\pi}}
    % \end{equation*}

% \subsection{Using post hoc Explainers on Non-linear Models}    
%     \label{sec:nonlinear} %use taylor series argument
%     \textcolor{blue}{ronilo: I'm worried to talk about other explainers in this section because some of them (SmoothGrad and C-LIME actually do directly approximate the marginal effect), so if I talk about them here the reviewers would criticize us for not using those methods in all of our experiments}
%     %\textcolor{red}{it's a discussion and is not supposed to be so technical. If you can write out the maths, can we just prove it? If you can only do a hand wavy argument, can we avoid using maths or try to use less maths? I also feel it is a little short. We can consider including it in other subsections as a paragraph.}
%     Intuitively, since we demonstrated that SHAP and LIME generally cannot extract the marginal effects from linear models (i.e. the weights of the linear models), it also makes sense that they should not be expected to be able to extract the marginal effects from non-linear models (the gradient at the instance being explained $\nabla \mathcal{M}(\bxi)$). This is the case because $\nabla \mathcal{M}(\bxi)$ is not the estimand of SHAP and LIME. Consider for example cases where $\mathcal{M}$ is a scalar valued real analytic function. We will discuss how LIME and SHAP behave when explaining such $\mathcal{M}$. 
    
%     If $\mathcal{M}$ is linear, then the target in LIME's regression is $\by=\bX \nabla \mathcal{M}(\bxi)$. As LIME is a linear approximator, it will then just output $\nabla \mathcal{M}(\bxi)$ in the best case (as we saw in Theorem \ref{thm:lime_thm}. However, if $\mathcal{M}$ is non-linear, we can see that LIME will generally not be able to extract $\nabla \mathcal{M}(\bxi)$ because $\by \neq \nabla \mathcal{M}(\bxi)$. Instead, $\by$ is given in terms of a power series with higher order terms such as $\nabla^2 \mathcal{M}(\bxi)$. 

%     As shown by \cite{StrumbeljErik2014Epma}, when $\mathcal{M}$ is linear the SHAP coefficients are proportional but not equal to the marginal effects. In general, in the expression of the SHAP value for the $i$th feature in the prediction $\mathcal{M}(\bxi)$, each summand %has a factor $\frac{f(h_\bxi(z'))-f(h_\bxi(z'\setminus i))}{M!}$, where $z'$ is a coalition vector and $M$ is the maximum number of members of a coalition. 
%     includes a factor which is a  rough estimate of $\pdv{\mathcal{M}}{X_i}\big|_{\bxi}$ (a first order difference of $\mathcal{M}$). So, although SHAP values are related to the gradient of $\mathcal{M}$, SHAP is not directly computing or estimating $\nabla \mathcal{M}(\bxi)$.
    
%     %%To illustrate this more concretely, let us assume the settings of Theorem \ref{thm:lime_thm} but allow for $f$ to possibly be a real analytic function (that is, one that can be described locally by a convergent power series). As in Theorem \ref{thm:lime_thm}, suppose we want to explain the prediction of $\mathcal{M}$ at a point $\bxi\in \mathbb{R}^D$ with LIME. Let $\bH$ be the projection matrix produced by LIME's regression, so that $e_{\text{LIME}}(\bxi) = \bH \mathcal{M}(\bX)$ where $\bX$ is a random matrix whose rows are sampled around $\bxi$. If $\mathcal{M}(\bx)=\boldsymbol{\beta}^T\bx$ then $e_{\text{LIME}}(\bxi) = \bH \bX \boldsymbol{\beta}$ and that in the large sample limit, $\bH \bX \mapsto \bI$ so that $\boldsymbol{\beta}$ is recovered. In the case that $f$ is non-linear, $\by$ is not $\bX \nabla f$. Instead, we consider that, viewing $f$ in a Taylor expansion using a  matrix form of Taylor's theorem, the product $\bH\by$ will generally not be $\bH \bX \nabla f$ and thus $\nabla f$ cannot be recovered. 
%     % the desired output is $\nabla f (\bxi)$ (the gradient of $f$ at $\bxi$). However, using the multivariate form of Taylor's theorem, we have $e_{\text{LIME}}(\bxi) \simeq  \bH(\tilde{\by}_1,\ldots,\tilde{\by}_N)^T$ where $\tilde{\by}_n$ is an approximation of $f(\bx^{(n)})$ for $n=1,\ldots,N$ and %nabla is 1xd
%     % %1x1
%     % %(1,1) + ()(1,D)
%     % \begin{equation*}
%     %     \tilde{\by}_n  = f(\bxi) + (\bx^{(n)} - \bxi)\nabla f(\bxi)^T + \frac{1}{2}(\bx^{(n)} - \bxi) \nabla^2 f(\bxi) (\bx^{(n)} - \bxi)^T
%     % \end{equation*}
%     % % Thus
%     % % % y is N x 1              nxd          dx1           nxd dxd 
%     % % \begin{equation*}
%     % %     \by = f(\bxi)\bI_N + (\bX - \bXi) \nabla f(\xi)^T + \frac{1}{2} (\bX-\bXi) \nabla^2 f(\xi) (\bX - \Xi)
%     % % \end{equation*}
%     % Comparing this to the linear case, we see that if $f(\bx)=\boldsymbol{\beta}^T\bx$ then $\tilde{\by}_n = \boldsymbol{\beta}^T \xi + \boldsymbol{\beta}^T (\bx^{(n)} - \bxi) = f(\bxi) + \nabla f(\bxi)(\bx^{(n)}-\bxi)= \nabla f(\bxi) \bxi + \nabla f(\bxi)(\bx^{(n)}-\bxi)$.
    

    


\subsection{General Concerns About Using Post Hoc Explainers to Get Data Insights}

    
    Although our paper focuses on the two most widely used post hoc explainers,our concerns about extracting data level $\mathbf{X} \mapsto Y$ insights through post hoc explainers are considerably more general. 
    The Rashomon effect \citep{semenova22rashomon,donnelly2023rashomon,rudin2024amazingthingscomehaving} in machine learning models shows that in any given learning problem there can be an abundance of models with similarly high predictive performance. Each of these models represents a distinct hypothesis about the relationship between the inputs X and the outputs y. As a result, the existence of many near-optimal hypotheses implies that even a perfect explanation of a single model is unlikely to provide a faithful or comprehensive summary of the true relationships encoded in the data.
    

    
    
    Explainable machine learning (XAI) is a relatively young field, having begun in the late 2010s as a subfield of machine learning and only recently being explored as a means of explaining ground truths $\mathcal{G}$ rather than models $\mathcal{M}$. A principal challenge consists in the fact that XAI has not yet figured out best practices for extracting information about $\mathcal{G}$ in a way that mitigates the Rashomon effect. After all, the first post hoc explainer that aims to directly mitigate the Rashomon effect (RID \cite{donnelly2023rashomon}) was only published in 2023.
    
    \subsection{Using Explainers as Hypothesis Generators}

While this paper cautions against using post hoc explanations to make definitive claims about the underlying data-generating process, this does not imply that post hoc explainers are without value.  Instead of serving as tools to \textbf{validate} hypotheses or as evidence about $X \mapsto Y$ relationships, post hoc explanations are best viewed as tools to \textbf{propose} hypotheses that can guide subsequent empirical investigation.

When used in this exploratory capacity, post hoc explanations can help researchers reveal potentially important variables, uncover candidate mechanisms, or suggest heterogeneity patterns that merit closer scrutiny. However, because explanations are not guaranteed to be data-aligned—and may reflect one of many plausible models consistent with the data—any hypothesis they generate must be treated as provisional. For example, examining the degree of agreement in explanations across equally accurate models within a Rashomon set can help assess whether an explanation reflects a stable pattern in the data or one that is sensitive to modeling choices. Low agreement serves as a warning that hypotheses generated from explanations should be interpreted with particular caution.

Credible inference therefore requires follow-up analyses using methods explicitly designed for identification, such as regression-based approaches, causal inference techniques, or experimental designs. Several recent studies illustrate this more appropriate workflow. For example, \cite{zhang2023can} use SHAP to identify potentially important features for restaurant survival, but then rely on causal forests to estimate treatment effects and assess whether those features exert a causal influence. \cite{hanauer22boosting} treat traditional regression as a trusted baseline against which SHAP explanations are compared, rather than as conclusions drawn directly from the explanations themselves. \cite{feng2025ai} uses SHAP to rank important predictors  of celebrity visual
potential and then immediately compare the empirical results with theories to validate the hypotheses. \cite{senoner2022using} uses nonlinear modeling with Shapley additive explanations to infer how a
set of production parameters and the process quality of a manufacturing system are related, and then validate the hypotheses with real-world application. In these cases, post hoc explainers play an exploratory role, while inferential claims are delegated to methods with clearer identification guarantees.

This distinction between hypothesis generation and hypothesis validation is critical. When post hoc explanations are used to guide inquiry rather than to justify conclusions, they can complement existing empirical tools without overstating their inferential content.
    % \tw{I think we should remove this paragraph. don't talk about psychology at all in this paper.}
    % Second, the scope of XAI has not yet widened beyond the realms of computer science and machine learning to cover psychology. Consideration of how humans use and process explanations is of paramount importance; a pipeline which overcomes all of the technical issues we have identified in this paper would nevertheless be practically useless if the explanations it generates are incomprehensible. That post hoc explainers give explanations in forms that naturally cohere with human psychology has had deleterious effects. For example, although researchers view SHAP's rigorous game theoretic foundation as a boon, the resulting highly technical meaning of Shapley values ends up being a source of confusion for end users. Human evaluation studies \citep{kaur2020interpreting,passi18trust} have shown that that because data scientists do not fully understand Shapley values, they fill in what they didn't understand about SHAP explanations with their own biases. %This general problem of explainers themselves needing to be explained so that their explanations can be understood is applies to more than just SHAP. For example, 
    
    % Plainly worded natural language explanations are the gold standard for explanations that XAI should strive for, and the current technicality of SHAP and other explainers prevents them from achieving this goal. Whereas a natural language explanation can typically be understood by a native speaker of that language without aid, SHAP and many other explanation methods have to themselves be explained prior to use. Additional examples include Accumulated Local Effect (ALE) plots \citep{apley2020visualizing}, Anchors \citep{ribeiro2018anchors},  Convolutional Neural Network Attention mechanisms \cite{itti1998model}, and counterfactuals \citep{mothilal20dice}, which all fall short of the gold standard because of the technical expertise required for their use.  

    
    
    %Many papers employing human evaluation studies have appeared in recent years demonstrating scenarios in which post hoc explainers were not helpful to users because the explainer was wrong and/or the user misunderstood the explainer~\citep{Krishna2022TheDP,kumar20shap,weerts19human,dieber20assessing,yacoby22judges}
    
    
    %\textcolor{red}{say sth to connect to our paper first and then transition to others. sth like "while LIME and SHAP produce explanations in the form of feature importance, other post hoc explainers ...", to introduce that there are various forms of explanations, which could makes interpretation challenging} . \textcolor{red}{For example... (one example would be helpful to give readers some sense of the concern)} Even SHAP, which produces explanations in the form of feature importance that is relatively easy to understand, still confuses users. 




% \subsection{Final Remarks}% and Future Work}
% This paper aims to highlight the issue with a rising use of post hoc explanations that we identify as inappropriate. To demonstrate the seriousness of this problem, we show that the misalignment occurs even in the simplest context, where the ground truth $\mathcal{G}$ is linear and features are mutually independent. Even in this simple case, SHAP and LIME produce misaligned explanations. Additionally, explanation quality seems to be limited by feature dimensionality. This being the case expect that in realistic business applications, the data-alignment of post hoc explainers can only be worse that what we have shown in our paper. Real decision processes are often non-linear and involve feature interactions and unobservable features; naturally, these added complexities should reduce the efficacy of SHAP and LIME, making it even more challenging to obtain data-aligned explanations. Given the uncertainty regarding the validity of post hoc explanations, researchers should view them skeptically and consider them as hypotheses rather than definitive claims about reality. 

%\tw{should Figure 13 be updated since we are in the middle of 2024 and there should be more papers in 2024} That is the updated version. 
 
 % When feasible, we advocate for intrinsically interpretable models to be used instead. %When data have unobserved confounders, randomized experiments should be used to identify treatment effects.


    % reword to remove word MISUSE; need to say instead that users get incorrect misunderstanding
    %Throughout this paper, we have shown how and why the use of post hoc explainers could lead to confusion for business researchers and bad decision making for businesses. Our experiments demonstrated SHAP and LIME's inability to produce consistent and faithful explanations. Possible consequences of this include businesses or customers trying to improve upon some outcome, which could waste time and resources by adjusting for an irrelevant feature that LIME or SHAP implied was important, or even worsen the outcome they were trying to improve. On the surface, model-agnostic post hoc explainers are attractive because they are easy to use and applicable to any model with any dataset. However, their propensity toward unfaithful and / or inconsistent explanations decreases their efficacy. Due to the stakes involved with understanding business problems, post hoc explainers should only be used with great caution. 
        

% Acknowledgments here
% \section*{Acknowledgements}
% We thank Zhongju Zhang, Tat Chan, Shan Huang, and Katherine Ragodos for comments which greatly improved the manuscript. We are grateful to Palle Jorgensen for sharing his wisdom and perspective on the mathematical portions of this work. We would also like to thank the participants of the 2024 Hong Kong Quant Marketing Mini-Conference, BizAI Conference, Production and Operations Management Society (POMS) Conference, Workshop on Information Technologies and Systems (WITS), as well as the participants of the 2023 INFORMS annual meeting and INFORMS Workshop on Data Mining \& Decision Analytics.

% References here (outcomment the appropriate case)

% CASE 1: BiBTeX used to constantly update the references
%   (while the paper is being written).
%\bibliographystyle{informs2014} % outcomment this and next line in Case 1
%\bibliography{<your bib file(s)>} % if more than one, comma separated

% CASE 2: BiBTeX used to generate mypaper.bbl (to be further fine tuned)
%\input{mypaper.bbl} % outcomment this line in Case 2

%If you don't use BiBTex, you can manually itemize references as shown below.


% \bibliographystyle{informs2014}
% \bibliography{refs}



% ... your document ...

\FloatBarrier
% \newpage
{
\tiny
\bibliographystyle{informs2014}
\bibliography{refs}}
% \endgroup



\newpage

% dimensions finds only 206 but indexes preprints
% scopus finds over 400 but doesn't index preprints
%% preprint not necessary
% need to ensure that the papers are actually explaining a black box
% e.g. tong will anonymize my quotations
% \section*{Appendix A: A Survey of Business Papers that Use SHAP or LIME} \label{apx_sec:lit_review}
    % The appendix contains additional details and experiments. It begins with information about how we conducted our literature review and how we judged papers for inappropriate use of explainers in Appendix Section \ref{apx_sec:lit_review}. We then provide all details relevant for Section \ref{sec:causes_of_misalignment}. First, we give proofs and experimental justification for theoretical claims in Appendix Section \ref{apx_sec:proofs}.
    % Then we detail the experimental designs used in our investigation of correlation and predictive performance in Appendix Section \ref{apx_sec:misalignment_exp_details}. %We expound upon the problem of explainer inconsistency in Appendix Section \ref{apx_sec:inconsistency} using experiments that demonstrate how and why it occurs. 
    % The next part of the Appendix is about our robustness checks. Appendix Section \ref{apx_sec:mitigation_details} details our hypothesis testing framework for investigating whether a strategy improves over a baseline pipeline. Appendix Section \ref{apx_sec:additional_mitigations} demonstrates some additional robustness checks that we tried that were less effective than the ones in the main body of the paper. Appendix Section \ref{apx_sec:econometric_details} has details about the dataset we used in our econometric case study. Finally, we give the results of a repetition of all experiments in the main body of the paper using a synthetic direct marketing example in Section \ref{apx_sec:marketing}
    
    % \subsection{}
       
        % As part of our literature review process, we cataloged papers at the intersection of business and explainable machine learning. We searched through each of the INFORMS journals and the Web of Science for publications citing at least one of LIME~\citep{ribeiro16lime} or SHAP~\citep{lundberg17shap} and have either ``LIME" or ``SHAP" appear in the body of the article. On Web of Science, the results were filtered by topics related to business, such as management, economics, supply chain \& logistics, operations research \& management science, etc. We also searched on SSRN because business journal papers usually take a few years to get published, so there are many preprints available on SSRN.  We plot the number of papers citing LIME or SHAP by year in Figure \ref{apx_tab:citation_count_graph}. %the following Web of Science topics: transportation, management, economics, supply chain \& logistics, economic theory, and operations research \& management science. Marketing does not appear in the table because Web of Science does not include it as a distinct category among business adjacent fields. Figure \ref{fig:citation_count_graph} shows how the frequency of such papers has increased over time. 
        % Figure \ref{fig:treemap} shows the breakdown of the papers we identified by category. %index any such papers that cited WachterCF~\citep{wachter17wachtercf} or DiCE~\citep{mothilal20dice}.

        % make paragraph or table with prototypical language forms
        % labeled with relevant metric
        % 10 examples of each
        % quotations without citation 
        %We took a close look at these papers and characterized the papers' use of explainers by following two steps. First, we searched for evidence that it was understood that SHAP and LIME explain model outputs, not the data. Next, we searched for problematic language in four forms and recorded quotations of each type (if available). Table \ref{apx_tab:mistakes} provides examples of each language form. The first form was text where the authors either asserted that features which were important to their predictive model(s) per LIME or SHAP were also drivers of their target variable in reality. The second form consisted in recommendations or implications that individuals, groups, managers, or policy-makers should change or guide their behavior on the basis of LIME/SHAP feature importance. In the third form, authors claim that SHAP and LIME explanations are evidence of (or even confirm) an existing hypothesis from prior related work.  The final form was unclear wording that could lead readers to think that the authors were making an assertion of the first form. %Having searched a paper for each form of language, we labeled each paper with 1) why the authors used LIME or SHAP in their paper and 2) whether the paper made an error in using LIME or SHAP (by using language of the first three forms), contained potentially misleading wording, or was error-free. 
        
            % \begin{table}[]
            % \centering
            % \begin{tabular}{l|c}
            % \toprule
            % \textbf{Web of Science Category}         & \textbf{Number of Papers} \\ \hline
            % Operations Research / Management Science & 159                     \\ \hline
            % Management                               & 55                      \\ \hline
            % Economics                                & 47                      \\ \hline
            % Business Finance                         & 42                      \\ \hline
            % Business                                 & 31              \\ \bottomrule
            % \end{tabular}
            % \caption{Number of papers citing either LIME or SHAP that belong to business-adjacent Web of Science categories.}
            % \label{tab:citation_count}
            % \end{table}

        % \begin{figure}[ht]
        %     \centering
        %     \includegraphics[width=.5\textwidth]{AppendixFigures/citation_count.png}
        %     \caption{Year-by-year breakdown of number of papers citing either LIME or SHAP that belong to business-adjacent Web of Science categories.}
        %     \label{apx_tab:citation_count_graph}
        % \end{figure}
 % We went through each paper to find those that that actually used actually SHAP or LIME to explain a machine learning model instead of just citing them. We stopped at finding 150 such papers.  Out of the 150 papers, 91\% used SHAP and 14\% used LIME.
        % \begin{figure}
        %     \centering
        %     \includegraphics[width=.5\textwidth]{Figs/treemap.png}
        %     \caption{Breakdown of number of papers citing either LIME or SHAP by web of science category.}
        %     \label{tab:treemap}
        % \end{figure}


    % Please add the following required packages to your document preamble:
    % \usepackage{booktabs}
    % \usepackage{lscape}

    % Ambiguous because it's not clear if they mean that the feature is useful for prediction of the target by their model or in reality
%     \begin{landscape}
%    \begin{table}[htbp] 
%     \tiny
%     \caption{Types of Language Used when Interpreting Explanations}
%     \begin{tabu} to 1\linewidth { | X[l] | X[l] | X[l] | X[l] | X[l] | X[l] |}
%        \hline
%        \textbf{Language Form} & \textbf{Example 1} & \textbf{Example 2} & \textbf{Example 3}& \textbf{Example 4} & \textbf{Example 5}\\
%        \hline
%        \textbf{Correct Usage} 
%        & However, as outlined in Rudin (2019), the explanations obtained may not be sufficient for a sensitive decision-making process as the results from these methods can be inaccurate representations of the original relationships and potentially misleading. 
%        & we analyze feature importance metrics to shed light on how machine learning models arrive at their decisions 
%        & We now study the importance of characteristics and their interactions for the performance of gradient boosting and random forests. 
%        & To further understand the inner workings of the language-model classifier, we employ Local Interpretable Model-agnostic Explanations 
%        & The advantage of this method is that for every prediction the RNN model makes, LIME examines nearby instances and their predictions, and tries to ascertain which parameters have influence over the output using instances in the local environment  \\ \hline 
%        \textbf{Extrapolationation}
%        & If there are more blue dots with negative SHAP values and more red dots with positive SHAP values, the
% values of that feature would be positively correlated with the prediction results. In other words, customers with higher scores for this characteristic are more likely to be churners. 
%        & Positive SHAP values indicate that the variable has a facilitative effect on bike-sharing use, while negative values indicate a suppressive effect
%        & Here, SHapley Additive exPlanations (SHAP) systematically uncover the crucial elements that govern guest satisfaction during hotel stays.
%        & XGBoost, is utilized to capture the quantitative associations, and the SHAP (Shapley Addictive exPlanation) model is used to interpret the results from XGBoost. Our results provide empirical support for prior theoretical findings based on mathematical equilibrium models  
%        & Shapley values to understand how various items explain analyst forecasts are generated by utilizing SHAP after XGBoost, when the machine learning model predicts analyst implied ETRs (using the same inputs as predicting actual future ETRs). These Shapley values measure the extent to which various financial
% statement items can explain analyst forecasts \\ \hline
%        \textbf{Advocacy for Recourse}
%        & the method used in this paper has effectively extracted the important influencing variables. On the one hand, it is convenient for investors to quickly evaluate high-polluting enterprises affected by financing, and on the other hand, it can help the government to better promote the transformation of high-polluting enterprises when introducing green credit and green finance
%        & When the product contains quad-core (blue), it has a negative impact on customer choices. This relationship implies that a higher spec is recommended for AP in the smartphone
%        & Since the number of bus stops, bus lines, bike-sharing trips, parking lots and pedestrian shed ratio have high SHAP importance and substantial effects on metro ridership, relevant last-mile transport facilities should be given sufficient attention by transport planners and operators to enhance these access modes to metro stations.
%        & Their SHAP values also raise with the increase of corresponding CV values, revealing that alternative routes with large difference in travel distances, travel speeds or signalized intersection numbers would lead to higher degrees of disequilibrium states at the road network.  
%        & The useful forecasting indicators—higher-order moments—used in this study are of great applied economic importance for helping investors, market participants, and policymakers to forecast the volatility of China’s oil futures more effectively, reduce investment risk, and achieve better returns. \\ \hline
% %        % \textbf{Unclear if Assertion Pertains To $\mathcal{M}$ or $\mathcal{G}$}
% %        &  SHAP can fairly capture the contribution of each predictor and reflect to what extent can each predictor impact the outcome variable
% %        & But credit and the slope of the yield curve, both domestically and globally are consistently the four most
% % important predictors. Together, these results strengthen the view that these key variables are robust indicators for predicting financial crises.
% %        &  According to the feature importance ranking, the top feature may be the potential indicator to trigger financial warning. 
% %        &  These information can be leveraged by the analyst to understand how to reduce the bankruptcy probability of a specific firm or to get insight in which balance sheet fields need to be improved to increase the rating, therefore providing a business instrument to actively manage clients and structured finance deals. 
% %        & By analyzing SHAP values, we identify the credit features that significantly influence borrowers' store pointss, gaining insights into the mechanisms behind these features' impact on personal store pointss. \\ \hline
%        \textbf{OLS Alternative}
%        & Hence, the applied black box algorithm in combination with Shapley values allows to unbox the black box ML approach and provides similar economic insights as provided by pooled OLS
%        & we choose to employ machine learning algorithms instead of traditional econometric techniques. The reason is that machine learning offers the following advantages: (1) Machine learning models can deal with non-linearity and complexity...
%        %, which allows better modelling of complicated relationships between variables; 
%        (2) % Traditional econometric techniques may encounter difficulties when dealing with large datasets due to computational limitations. 
%        Machine learning models are often designed to handle large datasets and can be scaled efficiently; (3) Machine learning models can learn from data without the need for explicit model specification. %They can learn patterns and relationships from the data...% and are therefore well suited for cases where the underlying relationships are not well understood and 
%        (4) Machine learning models can better capture interaction effects between variables...%, allowing a better understanding of how different factors influence each other. 
%        In this way, we then display how key features influence the severity of traffic accidents...% in the early, middle and late stages of the Covid-19 pandemic.
%        & To avoid biased estimates and capture non-linear associations, Light Gradient Boosting Machine (LightGBM) model and SHapley additive explanations (SHAP) model are combined to explore the factors affecting the competitiveness of bike-sharing to underground.
%        & ...conventional linear HAR models may falter. Additionally, the presence of nonlinearity contributing to volatility forecasting among higher-order moments cannot be identified using conventional linear models. However, machine learning (ML) methods have proven to be effective in addressing multicollinearity and accommodating more sophisticated functional forms, enabling the exploration of complex nonlinear relationships among variables 
%        & Machine learning techniques are regarded as a viable alternative to traditional models for traffic forecasting, due to their remarkable ability to capture nonlinearities and exhibit robustness against diverse distributional properties of data. For this reason, machine learning techniques are commonly perceived as a compelling approach for policy and planning analysis \\ \hline
%    \end{tabu}
%    \label{apx_tab:mistakes}
%    \end{table}
%    \end{landscape}
%\section*{Appendix O: Unorganized Temporary Holding Place}
% \setcounter{section}{0}
% \renewcommand{\thesection}{\Alph{section}}






%      Additionally, as was done in \cite{tang2024china}, plausible post hoc explanations can be used to aid classical regression analyses as a feature selection method. That is, researchers can opt to use only those features which are said to be important by a post hoc explainer(s) in their regressions but should keep in mind that those features may not actually be important to $\mathcal{G}$. \textcolor{red}{I think this is controversial. This is relevance and we said relevance is not 100\%}
    
    
    
% \textbf{How to Make Fruitful use of post hoc Explainers}.
% \cite{Usc paper and the other paper using shap to select features}
%  Model-agnostic post hoc explainers certainly have legitimate use cases, but due to the nature of real business datasets and stakes associated with business problems, we urge the business community to use them with caution and with tempered expectations. The problems of data-misalignment and inconsistency make it difficult for researchers to choose explanations in a principled fashion and ascertain their truth value. Thus, without sufficient domain knowledge or evidence, a researcher can neither believe nor trust the content of a post hoc explanation, nor use it to make inductive inferences. In such a context, the only truly principled view of a post hoc explanation is as a hypothesis. 
 

    
    %Then their explanations can be used as a basis for more rigorous analysis. For example, the treatment effects of features identified as important by the explainers can be estimated using causal forests \citep{zhang2023can}.   
    
        %The outline given in Section \ref{sec:summary_and_strategies} can undoubtably be improved and tailored to user-specific situations, and so is worth expounding. We do so below, and provide additional information that could be useful in different use-cases. 


        %During Prediction stage, before even beginning model development, users should use correlation analysis and any available domain knowledge to identify correlated or colinear features, and then decide whether to drop them from the dataset. % attempt to understand the causal graph associated with their problems and employ statistical tests of unconfoundedness assumptions. With this knowledge, users need to decide whether it is worthwhile to use post hoc explainers instead of randomized experiments. If so, the data should be pre-processed to have the desirable properties we have outlined previously. Domain knowledge and knowledge of the causal graph are useful for this purpose~\citep{zhao2021causal}.% to remove redundant features, causal descendants of a collinear pairs of causally linked features. Confounders should be controlled for if possible.
        %Then, unless users are only interested in explaining a particular black box model, we propose that many be generated and used for explanations. This is due to the Rashomon Effect\citep{semenova22rashomon}, which states that models (interpretable or otherwise) with approximately the same predictive performance will represent different relationships of the input features to the target. In other words, equally good models will have different (possibly conflicting) explanations and interpretations. In the best case, users could leverage explanations entire Rashomon sets of good models from specified model classes and training datasets. We showed that this technique was useful when averaging explanations over a rather small set of NNs with high ROC AUC using LIME explainers which correctly classify the input being explained. Additionally, because we found that the complexity and generalization error of $\mathcal{M}$ was related to post hoc explanation quality and variability, we propose that researchers gather as much training data as feasible.%.qasw11111 and take care not to overfit their models. %Overfitting is especially detrimental because the complexity of the decision boundary of $\mathcal{M}$ contributes to the instability of the explainer. After all, even a perfect explainer for a model which is highly locally nonlinear must be unstable. Thus, needless complexity in a decision boundary due to overfitting could cause an explainer of the decision function to be unstable when it should not be.  

       
        %When starting the explanation stage, users should keep in mind that post hoc explanation can be accidentally misused due to a fundamental misunderstanding of the explainer. For example, \cite{kaur2020interpreting} showed that practitioners misused explanations due to having an incorrect understanding of Shapley values. \cite{jacovi2023diagnosing} also discusses the possibility of users coming away with an incorrect understanding of incomplete explanations because they fill in missing information with their own biases. Users should ensure that they have a complete understanding of how explainers they intend to use work and how their biases may affect their understanding of the explanation. Having done this, we recommend researchers to view post hoc explanations simply as hypotheses and avoid using them to assert novel findings about $\mathcal{G}$ or claiming that they confirm existing hypotheses. 
        
        %As we have discussed, due to the problems of misalignment and inconsistency, researchers should not trust that a single explanation from any explainer applied to any black box model is true unless they can corroborate the explanation. Ideally, explanations would be corroborated using domain knowledge. Otherwise, users could consider comparing the explanations with patterns found in inherently interpretable models or use explainers only for exploratory analyses that are followed up with more rigorous analysis. For example, we recommend that researchers use robustness checks to figure out how confident they can be in the variability of explanations, and select or aggregate explanations from among the large set of explanations created during the robustness checks. 



\newpage
%\newcommand{\webappendix}{
    %\appendix
    % Redefine the section numbering: WA, WB, WC...
    \renewcommand{\thesection}{OA \arabic{section}}
    % If you want to customize subsections as WA.1, WA.2, etc.
    \renewcommand{\thesubsection}{\thesection.\arabic{subsection}}
    % Reset section counter to 1
    \setcounter{section}{0}
%}
\section*{Online Appendix}
The following online appendices provide additional information about our experimental designs, additional figures and experimental results, and some theoretical justification for the robustness checks we proposed involving LIME's hyperparameters. 

\section{Paper Review on the Uses of SHAP and LIME}
We conducted two waves of literature search and review.

In the first wave, we used Web of Science to identify publications containing either “SHAP” or “LIME” in the body of the paper within business-related subject areas, including management, economics, supply chain \& logistics, and operations research \& management science. This search yielded over one thousand papers published between 2017 and January 2026. Because publications in business journals typically involve long review cycles, we complemented this search with SSRN to capture relevant working papers. We then manually screened each paper to identify studies that substantively used SHAP or LIME for model explanation, rather than merely citing or mentioning these methods. We stopped once we had identified 150 such papers.

In the second wave, to ensure comprehensive coverage of papers published in leading journals, we focused on journals listed in the \textit{UTD 24}, the \textit{FT50}, and journals published by \textit{INFORMS}. Using the same search keywords (“SHAP” or “LIME”), we manually examined all returned papers and retained those that made substantive use of SHAP or LIME. This screening process yielded 56 papers.

After removing overlaps between the two waves, our final sample consists of 181 unique papers that make substantive use of SHAP and/or LIME. Among these papers, 94\% used SHAP and 13\% used LIME, with 7\% using both methods.

We next classified papers based on how post hoc explanations were interpreted and used. Specifically, we assessed whether authors explicitly or implicitly framed explanations as describing properties of the trained model, or instead extrapolated them to make claims about relationships in the underlying data-generating process. Papers exhibiting such extrapolation were coded as making data-level inferences. Based on this criterion, we categorized papers into three groups: (i) controversial papers that clearly extrapolated post hoc explanations to data-level claims; (ii) ambiguous papers whose wording could plausibly be interpreted as making such extrapolations; and (iii) proper papers that used explanations strictly for model interpretation, consistent with the intended scope of these methods.

To reduce potential labeling bias, we additionally employed GPT-5.2 to review each paper and generate multiple annotations. The full prompt is provided below. We applied strict criteria to classify a paper as having “misused” SHAP or LIME. A paper was given the “misuse” label only if all of the following conditions were satisfied: (1) the paper did not explicitly state that the interpretation was limited to the model level, nor caution readers against causal or data-level interpretations; (2) the LLM identified quotations suggesting an inference of $\mathbf{X} \mapsto Y$ relationships in the underlying data; (3) the LLM identified quotations that derived managerial insights or recommendations directly from the explanations; and (4) the assigned severity level was 4. These criteria are intentionally stringent, making the misuse classification conservative.

Finally, we compared human labels with LLM-generated labels, identified disagreements, and manually re-examined the corresponding papers to determine the final classification.

\begin{tcolorbox}[
  colback=gray!15,   % background
  boxrule=0pt,       % no border
  frame hidden,      % no frame
  breakable,
  left=6pt,right=6pt,top=6pt,bottom=6pt
]
\begin{Verbatim}[fontsize = \footnotesize]
You are assisting with a systematic literature review in business and management research.

Research Question
"Do authors misuse SHAP or LIME by generalizing model-level explanations (X → Ŷ) into claims 
about real-world relationships (X → Y)?"

Core Definition of Misuse
A misuse occurs when feature attributions from: - SHAP (including Shapley-value-based methods and 
variants), or -LIME (including variants) are interpreted as evidence about relationships between
X and the real-world outcome Y, rather than about the model’s prediction Ŷ.

Key Conceptual Distinction

Correct Usage (Model-Level): Statements about how features influence the model prediction (Ŷ). 
Example: "Feature X increases the predicted probability."

Misuse (Data-Level Interpretation): Statements about how features influence the real-world 
outcome Y. Example: "Feature X increases customer churn."

Important Clarifications 
1. This review is NOT about causality. 
2. Even associative language counts as misuse if it refers to real-world Y. Example: "X is 
associated with higher Y." 
3. If the paper explicitly states that explanations apply only to the model or predictions, do
NOT classify as misuse. 
4. Only evaluate language appearing in sections interpreting SHAP/LIME results (e.g., Results, 
Discussion, Implications).

Output Requirement Return a SINGLE JSON object following this EXACT schema:

{ "title": string|null, 
"journal": string|null, 
"year": number|null,
"explainer_used": "SHAP"|"LIME"|"Both"|"None"|"Other", 
"outcome_Y": string|null, 
"use_interpretation_pipeline": "Yes"|"No"|"Unclear",
"explicit_model_level_only": { "answer": "Yes"|"No", "evidence_quotes":[string] }, 
"infer_relationships_X_to_real_Y": { "answer": "Yes"|"No","evidence_quotes": \[string] },
"policy_managerial_recommendations_from_explanations": { "answer":"Yes"|"No", "evidence_quotes": 
[string] }, 
"overall_severity_of_misuse": 1|2|3,
"notes": string|null }

Field Definitions & Coding Rules

explainer_used: - "SHAP" includes Shapley-based methods and variants -"LIME" includes LIME and 
variants - "Both" if both are used - "Other" if other post-hoc explainers are used - "None" if no
explainer is used

use_interpretation_pipeline: Answer "Yes" only if ALL three occur: - The paper trains a 
predictive model - The paper applies SHAP/LIME (or variant) - The paper interprets feature 
attributions in text

explicit_model_level_only: Answer "Yes" ONLY if authors explicitly state explanations apply only
to the model or predictions, or caution against interpreting SHAP/LIME as reflecting real-world 
relationships.

infer_relationships_X_to_real_Y: Answer "Yes" if SHAP/LIME results are described using language 
that refers to the real-world outcome Y (e.g., increases Y, affects Y, drives Y, leads to higher Y).
If authors explicitly limit interpretation to predictions, do NOT code as misuse.

policy_managerial_recommendations_from_explanations: Answer "Yes" if managerial or policy 
recommendations are explicitly derived from SHAP/LIME explanations.

Evidence Rules - Evidence quotes must be verbatim - Maximum 40 words per quote - Do NOT 
paraphrase - If unavailable, use []

Misuse Severity Scale 
1 = No Y-level language 
2 = Minor or isolated loose wording 
3 = Clear, repeated data-level interpretation (X → Y)

Additional Rules - If title/journal/year are not available, set them to
null and explain in notes. - If evidence is ambiguous, explain reasoning
in notes. - Do NOT return anything outside the JSON object.
\end{Verbatim}
\end{tcolorbox}
\section{Data Generation and Model Construction Details for the Explanation Pipeline Demo}

In this first appendix we we document how the synthetic app downloading dataset and the accompanying XGBoost model trained on it, both used in Section \ref{sec:pipeline_demo}, were created.

\subsection{Details on Synthetic App Downloading Dataset}
    \label{apx_sec:app_details}

    The data generation process is based on the linear model shown below as equation \ref{eqn:gt}. This equation was written to express the idea that the probability that people will download a certain app increases with their income, phone usage, usage limits, their exposure to peer pressure and advertisements, positive opinion of the app, and tech savviness. It also expresses that the probability decreases with user age and the extent of their personal privacy concerns. 
    
\begin{equation}
\label{eqn:gt}
\begin{aligned}
\mathcal{G}(X) ={}& 0.0 - .03\cdot\text{Age}
+ 0.5\cdot\text{Income} \\
& {}+ 0.4\cdot\text{Time Spent on Phone}\\
&+ 0.15\cdot\text{Num Apps Installed} \\
& {}+ 0.8\cdot\text{Peer Influence}
+ 0.25\cdot\text{Ad Exposure} \\
& {}- 2.0\cdot\text{Privacy Concern}\\
&+ 0.6\cdot\text{App Rating Perception} \\
& {}+ 1.2\cdot\text{Tech Savviness}\\
&+ 0.7\cdot\text{Unlimited Data Plan}
+ \epsilon.
\end{aligned}
\end{equation}
    %\end{tcolorbox}
    \noindent where $\epsilon \sim \mathcal{N}(0,1)$ captures all the variation in the response variable that is not explained by the explanatory variables. The meanings of the feature names in Equation \ref{eqn:gt} as well as the statistical distributions used to simulate them when synthesizing user profiles appear below, in Table \ref{tab:app_features}.
    
    
    There are several important items about our procedure for synthesizing user profiles to note. First, after initially generating users' features according the distributions in Table \ref{tab:app_features}, the generated value ranges of ``Tech Savviness" and ``App Rating Perception" are clipped between 0 and 1. Then, ``Privacy Concern" was defined to be correlated with ``Tech Savviness" using the following definition.
    \begin{equation*}
        \text{Privacy Concern} = .7 \cdot \text{Tech Savviness} + \sqrt{1-.7^2}
    \end{equation*}
    We then clip the values of ``Privacy Concern" to be between 0 and 1. ``Privacy Concern" and ``Tech Savviness" then have a Pearson correlation of about .63. The most important is item to note is that the feature values that are input to $\mG$ are standardized into z-scores so that the so that the relative importances of the features can be compared directly using the coefficients of $\mG$. 
    
    The binary target is synthesized for users as follows. Using the logistic sigmoid function $\sigma(\cdot)$, we begin by assigning $Y = 1$ (downloaded) for cases with $\sigma(\mathcal{G}(\mathbf{x}))\geq 0.5$, and $Y = 0$ (ignored) otherwise. %A total of 5000 instances were generated, 4000 of which were used for training and the remaining 1000 used as the test set. 


  
    \begin{table}[ht]
\centering
\small
% \resizebox{\textwidth}{!}{%
\begin{tabular}{l|l|l}
\toprule
\small
\textsc{Feature} & \textsc{Meaning} & \textsc{Distribution} \\ \toprule
\small
Age        &        User Age        & $\mathcal{N}(35,10)$\\ 
Income        & User Income        & $\mathcal{N}(60000,20000)$ \\
Time Spent on Phone        & Average phone hours per day        & $\mathcal{N}(3,1)$ \\ 
Num Apps Installed         & Number of apps installed        & Uniform(0,100)\\ 
Peer Influence        & Measure of peer pressure to install        & Bernoulli(.4) \\ 
Ad Exposure        & Number of times app ad seen & Uniform(0,10) \\ 
Privacy Concern        & Measure of concern about app permissions        & see above \\ 
App Rating Perception    & Measure of impression of app store rating & Uniform(1,5) \\ 
Tech Savviness     & Measure of technological literacy  & $\mathcal{N}(.6,.2)$ \\ 
Unlimited Data Plan & Whether user has unlimited cellular data & Bernoulli(.5) \\ \bottomrule
\end{tabular}%
% }
\caption{Description of features for app downloading data.}
\label{tab:app_features}
\end{table}

    
\subsection{Details on the App Download Prediction Model}

The model whose explanations were used to produce Figure \ref{fig:sec_3_shap_lime_demo} in the main body of the paper was an XGBoost model from the Python xgboost package. It was obtained using 5 fold nested cross validation. The parameter search loop covered maximum tree depths in [3,4,5], learning rates in [.001,.01,.1], gamma in [0,.25,.5,.75,1], number of estimator in [100,500,1000], colsample\_bytree in [.5,1]. The subsample parameter was fixed at .9, and the number of early stopping rounds was fixed at 10. 

It is also worth mentioning that a root cause of this model's test set accuracy only being 72\% is the variance of the noise term $\epsilon$ in $\mG$. We have confirmed that it is the case that changing the distribution of $\epsilon$ from $\epsilon \sim \mathcal{N}(0,1)$ to, e.g., $\epsilon \sim \mathcal{N}(0,.1)$ makes it so that xgboost models trained as detailed in the previous paragraph can achieve at least 98\% test set accuracy. 
  


\section{Details of the Data Generation Procedure}
\label{apx_sec:81_details}

This section of the appendix describes how we generated the 81 datasets used in Section \ref{sec:main_alignment} of the main paper, as well as how the models we trained on those data were configured.

\subsection{Constructing the datasets}

We generated 81 binary classification datasets by varying four scenario-defining factors while holding sample size and noise fixed. The factor design is summarized in Table~\ref{tab:data_factors} of the main text. 
\begin{table}[htbp]
\centering
\small
\caption{Data Generation Factors and Configurations}
\label{tab:data_factors}
\begin{tabular}{lllc}
\toprule
\textbf{Factor} & \textbf{Variable} & \textbf{Type} & \textbf{Values} \\
\midrule
Sample size & \texttt{Num\_instance} & Fixed & 5{,}000 \\
Noise std. dev. & \texttt{Noise\_std} & Fixed & 0.1 \\
Number of features & \texttt{Num\_feature} & Varied & 10, 20, 30 \\
Correlation strength & \texttt{Corr\_strength} & Varied & 0, 0.5, 1 \\
Nonlinear terms & \texttt{Nonlinear\_num} & Varied & 0, 5, 10 \\
Interaction terms & \texttt{Interaction\_num} & Varied & 0, 5, 10 \\
\bottomrule
\end{tabular}
\end{table}\normalsize

Below, we summarize the step-by-step procedure used to generate the datasets.

\begin{algorithm}[htb]
\caption{Generating Synthetic Dataset}
\label{alg:synth-data}
\begin{algorithmic}[1]
\Require 
    \Statex $N \gets$ number of instances per class (total $= 2N$)
    \Statex $m \gets$ number of original features
    \Statex $s \gets$ noise standard deviation
    \Statex $\rho \gets$ correlation strength
    \Statex $p \gets$ number of nonlinear (squared) features
    \Statex $q \gets$ number of interaction terms
\Ensure 
    \Statex $\mathbf{X}_\text{final}$, $\mathbf{Y}$

\Statex \textbf{1. Generate Original Features:}
    \For{$i=1$ to $m-1$}
        \State $x_i \gets$ sample $2N$ values from $\mathcal{N}(0,1)$
    \EndFor

\Statex \textbf{2. Introduce Correlated Feature:}
    \State Randomly pick $k$, sample $\epsilon \sim \mathcal{N}(0,1)$
    \State $x_m \gets \rho \cdot x_k + \sqrt{1-\rho^2}\cdot\epsilon$
    \State $\mathbf{X} \gets [x_1,\ldots,x_m]$

\Statex \textbf{3. Nonlinear (Squared) Terms:}
    \State Select $p$ features, add squared terms

\Statex \textbf{4. Interaction Terms:}
    \State Select $q$ pairs, add interaction terms

\Statex \textbf{5. Coefficients:}
    \State Sample coefficients from $\text{Uniform}(-1,1)$
\algstore{myalg}
\end{algorithmic}
\end{algorithm}

\begin{algorithm}[htb]
\ContinuedFloat
\caption*{Generating Synthetic Dataset (Continued)}
\begin{algorithmic}[1]
\algrestore{myalg}

\Statex \textbf{6. Compute Probability Vector:}
    \State $z \gets$ linear + quadratic + interaction contributions
    \State $\mathbf{P} \gets \text{sigmoid}(z + \mathcal{N}(0,s))$

\Statex \textbf{7. Generate Labels:}
    \State $\mathbf{Y}[i] = 1$ if $p_i > 0.5$, else $0$

\Statex \textbf{8. Balance Dataset:}
    \State Sample $N$ positives and $N$ negatives
    \State Form $\mathbf{X}_\text{final}, \mathbf{Y}$
\State \Return $\mathbf{X}_\text{final}, \mathbf{Y}$
\end{algorithmic}
\end{algorithm}



\subsection{Constructing the Predictive Model $\mathcal{M}$}
\noindent
Beyond the synthetic data generation process, we also formalize the construction of the predictive black-box models $\mathcal{M}$. 
These models are trained on each generated dataset under varying training sizes and hyperparameters, 
ensuring a range of accuracies across pre-defined bins. 
The following algorithm provides a detailed step-by-step procedure for constructing $\mathcal{M}$.

\begin{algorithm}[htb]
\caption{Construct All Models $\mathcal{M}$}
\label{alg:construct-models}
\begin{algorithmic}[1]
\Require 
    \Statex $D = \{D_1, \ldots, D_{81}\}$ 
    \Statex $H = \{\text{max\_depth}, \text{n\_estimators}\}$
    \Statex $\text{valid\_sizes}$ 
    \Statex $A\_\text{bins} = \{[0.70,0.75), [0.75,0.80), [0.80,0.85), [0.85,0.90)\}$
    \Statex $\text{Models\_per\_bin}=3$
\Ensure $\mathcal{M} = \{\mathcal{M}_{i}^{j}\}$

\State $\mathcal{M} \gets \emptyset$
\For{each $D_i \in D$}
    \State $\mathcal{M}_i \gets \emptyset$
    \For{each bin $a \in A\_\text{bins}$}
        \State $\text{Count} \gets 0$
        \While{$\text{Count} < \text{Models\_per\_bin}$}
            \State Sample $s \in \text{valid\_sizes}$, $h \in H$
            \State Split $D_i$ into (train, test) of size $s$
            \State Train XGBoost $M_\text{temp}$ with $h$
            \State Evaluate $\text{acc}$ on test
\algstore{modelM}
\end{algorithmic}
\end{algorithm}

\begin{algorithm}[htb]
\ContinuedFloat
\caption*{Construct All Models $\mathcal{M}$ (Continued)}
\begin{algorithmic}[1]
\algrestore{modelM}

            \If{$\text{acc} \in a$}
                \State Add $M_\text{temp}$ to $\mathcal{M}_i$
                \State $\text{Count} \gets \text{Count}+1$
            \EndIf
        \EndWhile
    \EndFor
    \State Add $\mathcal{M}_i$ to $\mathcal{M}$
\EndFor
\State \Return $\mathcal{M}$
\end{algorithmic}
\end{algorithm}


% \subsection{Notes on Model Specifications}
% \label{app:constuction_models}

% As discussed in the main text (Table~\ref{tab:models}), two models are associated with each dataset. Model $\mathcal{G}$ reproduces the data-generating mechanism exactly, while Model $\mathcal{M}$ represents a black-box predictor with stratified accuracy. We omit the full table here for brevity, as it has been incorporated into the main body of the paper.

\section{Robustness Check Details and Results}
\label{apx_sec:robustness_results}
This section provides information about the motivation and design of the three robustness checks introduced in Section \ref{sec:robustness} of the main body of the paper. The plots illustrating the experimental results of the three robustness checks are also provided here. 

 The subplots in the figures in this section each show shaded PDF curves for one mode of data alignment for one explainer type. Each subplot will have two PDF curves: one representing the alignment of explanations produced baseline usage of SHAP and LIME and  one representing alignment of explanations produced by intervening on the factor relevant for a robustness check. %individual explanations of the best models on each dataset while the other represents averaged explanations over a Rashomon set of each dataset. 
Two-sample t-tests for the difference of means are used to indicate whether SHAP and/or LIME are robust to the intervention. Tests whose p-values are significant can be identified by checking for one the signs "(*)," "(**)," "(***)," or "(****)." The number of asterisks between the parentheses is suggestive of \emph{how} significant a p-value is. 

\subsection{Averaging Explanations within A Rashomon Set}

A useful way to describe the motivation of this robustness check is to relate it to the phenomenon where individuals seek out second or third opinions for high stakes medical diagnoses. We may view models in a Rashomon set like expert opinions (hypotheses) about a diagnostic problem ($X \mapsto Y$). The similarity with medicine ends there because whereas patients will be lucky to have access to 3 expert doctor opinions, data scientists can have access to myriad opinions from models and, therefore, will require a way to synthesize them in a coherent and concise manner. 

Since, for each dataset, we trained 10 highly accurate black box models, we have 10 hypotheses about each of the corresponding $X\mapsto Y$ relationships that can only be understood through post hoc explanation. As we are explaining the models using SHAP and LIME, explanations of each model can be represented as feature importance vectors. 

Our robustness check leverages one of the simplest ways to synthesize the SHAP and LIME explanations of each model in a Rashomon set: by simply averaging together each of the feature importance vectors obtained by explaining each model. This procedure produces a single feature importance vector equating the importance of each feature $X_j$ to the average of the importances assigned to that feature by explanations of the models in the Rashomon set.

We investigated whether such a routine is significant by, for each dataset, comparing the data alignment of explanations of A) the model with the highest test set accuracy with B) explanations averaged over the Rashomon set. Our results appear below in Figure \ref{apx_fig:apx_rashomon}. The two significant results we observed were that the averaged explanations had significantly \emph{lower} mean SHAP direction alignment, but significantly \emph{higher} mean LIME strength alignment. From these results we adduce that our averaging strategy cannot be used to completely remedy the misalignment of SHAP and LIME. 



   \begin{figure}[htb]
  \begin{subfigure}[t]{.45\textwidth}
    \centering
    \includegraphics[width=\linewidth]{Rebuttal-R1/appendix_figures/apx_shap_rashomon_3step_dir_v2.pdf}
    \caption{Directionality for SHAP}
    \label{apx_fig:shap_rashomon_3step_dir}
  \end{subfigure}
   \hfill
  \begin{subfigure}[t]{.45\textwidth}
    \centering
    \includegraphics[width=\linewidth]{Rebuttal-R1/appendix_figures/apx_shap_rashomon_3step_strength_v2.pdf}
    \caption{Strength for SHAP}
    \label{apx_fig:shap_rashomon_3step_str}
  \end{subfigure}
  \medskip

  \begin{subfigure}[t]{.45\textwidth}
    \centering
    \includegraphics[width=\linewidth]{Rebuttal-R1/appendix_figures/apx_lime_rashomon_3step_dir_v2.pdf}
    \caption{Directionality for LIME}
    \label{apx_fig:lime_rashomon_3step_dir}
  \end{subfigure}
   \hfill
  \begin{subfigure}[t]{.45\textwidth}
    \centering
    \includegraphics[width=\linewidth]{Rebuttal-R1/appendix_figures/apx_lime_rashomon_3step_strength_v2.pdf}
    \caption{Strength for LIME}
    \label{apx_fig:lime_rashomon_3step_str}
  \end{subfigure}
  \caption{Robustness of SHAP and LIME data alignment to averaging explanations over a Rashomon set}\label{apx_fig:apx_rashomon}
\end{figure}












    \subsection{Sensitivity to Explainer Hyperparameters}
    
    The idea of this robustness check is to check whether leveraging non-default values of explainers' hyperparameters can remedy misalignment. LIME explainers are initialized with parameters for a random seed, sample size $N$, and kernel bandwidth $\mu$. SHAP also has parameters for random seed and sample size $N$, but is known to be essentially deterministic and so SHAP's random seed is less salient. 
    
        \paragraph{Results for LIME Hyperparameters}
        
        LIME's random seed influences the regression used to build the surrogate model, and so has some effect on misalignment. We design one robustness check where the baseline uses LIME explainers with default random seed settings and the intervention uses an average over explanations created from LIME explainers initialized with different random seeds. The averaging here is done in the same manner as was done in the context of the Rashomon set. 
        
         Theoretical justification for this robustness check is provided in Theorem \ref{thm:lime_thm} in Appendix \ref{apx_sec:lime_thms}; it implies that \emph{increases} to both $N$ and $\mu$ beyond their default values have potential to yield comparatively better data alignment. For this reason, we treat explanations produced via LIME explainers with the default hyperparameter settings of $N=1000$ and $\mu\simeq 2.37$ as set $N=\mu=10000$ as the intervention we compare the baseline against. %$\hl{Is this correct? I thought we didn't actually increase $N$ for figure 11? } \textcolor{red}{-- I referred to your previous experiment, and here, both $\mu$ and N are increased at the same time, where $\mu$ refers to the kernel parameter of the LimeTabularExplainer function, and N refers to the num\_samples parameter of the explain\_instance function, both of which are set to 10000.}
        
        Figures \ref{apx_fig:apx_lime_kernel} and \ref{apx_fig:apx_lime_seed} below show that, both increasing $N,\nu$ and averaging explanations from LIME explainers initialized with different random seeds lead to the mean direction alignment was significantly \emph{lower},  but the mean strength alignment was significantly higher. 
        
           \begin{figure}[htb]
          % \begin{subfigure}[t]{.45\textwidth}
          %   \centering
          %   \includegraphics[width=\linewidth]{Rebuttal-R1/appendix_figures/apx_lime_seed_1step_dir.png}
          %   \caption{1 step Directionality}
          %   \label{apx_fig:a}
          % \end{subfigure}
          %  \hfill
          % \begin{subfigure}[t]{.45\textwidth}
          %   \centering
          %   \includegraphics[width=\linewidth]{Rebuttal-R1/appendix_figures/apx_lime_seed_1step_strength.png}
          %   \caption{1 step Strength}
          %   \label{apx_fig:b}
          % \end{subfigure}
          % \medskip
        
          \begin{subfigure}[t]{.45\textwidth}
            \centering
            \includegraphics[width=\linewidth]{Rebuttal-R1/appendix_figures/apx_lime_seed_3step_dir_v2.pdf}
            \caption{Directionality}
            \label{apx_fig:lime_seed_3step_dir}
          \end{subfigure}
           \hfill
          \begin{subfigure}[t]{.45\textwidth}
            \centering
            \includegraphics[width=\linewidth]{Rebuttal-R1/appendix_figures/apx_lime_seed_3step_strength_v2.pdf}
            \caption{Strength}
            \label{apx_fig:lime_seed_3step_str}
          \end{subfigure}
          \caption{Robustness for LIME's instability, varying its random seed}\label{apx_fig:apx_lime_seed}
        \end{figure}
        
           \begin{figure}[ht!]
          % \begin{subfigure}[t]{.45\textwidth}
          %   \centering
          %   \includegraphics[width=\linewidth]{Rebuttal-R1/appendix_figures/apx_lime_kernel_1step_dir.png}
          %   \caption{1 step Directionality}
          %   \label{apx_fig:a}
          % \end{subfigure}
          %  \hfill
          % \begin{subfigure}[t]{.45\textwidth}
          %   \centering
          %   \includegraphics[width=\linewidth]{Rebuttal-R1/appendix_figures/apx_lime_kernel_1step_strength.png}
          %   \caption{1 step Strength}
          %   \label{apx_fig:b}
          % \end{subfigure}
          % \medskip
        
          \begin{subfigure}[t]{.45\textwidth}
            \centering
            \includegraphics[width=\linewidth]{Rebuttal-R1/appendix_figures/apx_lime_kernel_3step_dir_v2.pdf}
            \caption{Directionality}
            \label{apx_fig:lime_kernel_3step_dir}
          \end{subfigure}
           \hfill
          \begin{subfigure}[t]{.45\textwidth}
            \centering
            \includegraphics[width=\linewidth]{Rebuttal-R1/appendix_figures/apx_lime_kernel_3step_strength_v2.pdf}
            \caption{Strength}
            \label{apx_fig:lime_kernel_3step_str}
          \end{subfigure}
          \caption{Robustness for LIME's instability, varying its kernel bandwidth}\label{apx_fig:apx_lime_kernel}
        \end{figure}
        
        \paragraph{Results for SHAP Sample Size}

            The relevant idea behind SHAP's sampling procedure for this section is that increasing the number of samples in the background dataset for SHAP explanations should lead to a better Shapley value approximation. Before we discuss the results however, we will address a potential point of confusion in the Python API for SHAP. There is an \textbf{nsamples} parameter that is passed to explainer objects that controls how many data samples SHAP will produce explanations for. This is \emph{not} the parameter we are changing in this robustness check. Rather, we are changing a parameter in the constructor for SHAP explainer objects controlling number of samples used for the background dataset. Specifically, setting a baseline of 500 samples, our intervention is to set the number of samples to 2000 (using the \textbf{shap.sample()} function). %Thus the intervention is actually to \emph{downsample}. Such downsampling is recommended for the sake of decreasing the computation time of SHAP.\footnote{A popular issue people have with SHAP is that it can take very long to produce explanations.} 

            Figure \ref{apx_fig:apx_shap_sample} suggests that our quadrupling of the number of samples for SHAP's background datasets was not associated with a significant change to SHAP's data alignment. More specifically, the results of our t-tests did not indicate a significant difference between the mean alignments for both direction and strength. In fact, the PDF curves are overall quite similar and for the most part coincide with each other. 

                    \begin{figure}[htb]
                   \begin{subfigure}[t]{.45\textwidth}
                    \centering
                    \includegraphics[width=\linewidth]{Rebuttal-R1/appendix_figures/apx_shap_sample_3step_dir_v2.pdf}
                    \caption{Direction}
                    \label{apx_fig:shap_sample_3step_dir}
                  \end{subfigure}
                   \hfill
                  \begin{subfigure}[t]{.45\textwidth}
                    \centering
                    \includegraphics[width=\linewidth]{Rebuttal-R1/appendix_figures/apx_shap_sample_3step_strength_v2.pdf}
                    \caption{Strength}
                    \label{apx_fig:shap_sample_3step_str}
                  \end{subfigure}
                  \caption{Robustness for SHAP's instability, varying its sample size}\label{apx_fig:apx_shap_sample}
                \end{figure}

\subsection{Sensitivity to Metric Parameters}

This third robustness check is based on our acknowledgment of the possibility that the value we assign to the free step size parameter used in the definition of our direction alignment metric influences our appraisals of SHAP LIME. As outlined in Section \ref{sec:metrics} in the main body of the paper, the direction alignment metric definition includes a perturbation grid $\Delta = \{m\cdot\texttt{Std}\mid m\in\{-M,\ldots,-1,1,\ldots,M\}\}$ depending on a number of perturbation steps $M$. By default we have $M=3$; the baseline $\Delta$ is $\{-3,-2,-1,1,2,3\}$. The interventions we consider in this robustness check are to set $\Delta$ instead to $\{-1,1\}$, $\{-2,2\}$, and $\{-,3,3\}$. 

The experimental results of this robustness check are shown below in Figures \ref{apx_fig:apx_lime_step} and \ref{apx_fig:apx_shap_step} for LIME and SHAP, respectively. Both figures have four subplots; the first subplot (a) represents the baseline while the other three are part of the intervention. Evidence for our intuition that LIME alignment could suffer as a result of larger step sizes appears in Figure \ref{apx_fig:apx_lime_step}; the mean direction alignment is at least .9 for all choices of $\Delta$ except for the most extreme: $\Delta = \{-3,3\}$. In contrast, it is not apparent from Figure \ref{apx_fig:apx_shap_step} that the choice of $\Delta$ is associated with a significant change to the mean or median direction alignment for SHAP. 

   \begin{figure}[H]
  \begin{subfigure}[t]{.24\textwidth}
    \centering
    \includegraphics[width=\linewidth]{Rebuttal-R1/appendix_figures/apx_lime_step_123_dir_v2.pdf}
    \caption{Step = 1,2,3 (default)}
    \label{apx_fig:lime_step_123_dir}
  \end{subfigure}
   \hfill
  \begin{subfigure}[t]{.24\textwidth}
    \centering
    \includegraphics[width=\linewidth]{Rebuttal-R1/appendix_figures/apx_lime_step_1_dir_v2.pdf}
    \caption{Step = 1}
    \label{apx_fig:lime_step_1_dir}
  \end{subfigure}
   \hfill
  \begin{subfigure}[t]{.24\textwidth}
    \centering
    \includegraphics[width=\linewidth]{Rebuttal-R1/appendix_figures/apx_lime_step_2_dir_v2.pdf}
    \caption{Step = 2}
    \label{apx_fig:lime_step_2_dir}
  \end{subfigure}
   \hfill
  \begin{subfigure}[t]{.24\textwidth}
    \centering
    \includegraphics[width=\linewidth]{Rebuttal-R1/appendix_figures/apx_lime_step_3_dir_v2.pdf}
    \caption{Step = 3}
    \label{apx_fig:lime_step_3_dir}
  \end{subfigure}
  \caption{LIME Step Perturbation}\label{apx_fig:apx_lime_step}
\end{figure}

   \begin{figure}[ht!]
  \begin{subfigure}[t]{.24\textwidth}
    \centering
    \includegraphics[width=\linewidth]{Rebuttal-R1/appendix_figures/apx_shap_step_123_dir_v2.pdf}
    \caption{Step = 1,2,3 (default)}
    \label{apx_fig:shap_step_123_dir}
  \end{subfigure}
   \hfill
  \begin{subfigure}[t]{.24\textwidth}
    \centering
    \includegraphics[width=\linewidth]{Rebuttal-R1/appendix_figures/apx_shap_step_1_dir_v2.pdf}
    \caption{Step = 1}
    \label{apx_fig:shap_step_1_dir}
  \end{subfigure}
   \hfill
  \begin{subfigure}[t]{.24\textwidth}
    \centering
    \includegraphics[width=\linewidth]{Rebuttal-R1/appendix_figures/apx_shap_step_2_dir_v2.pdf}
    \caption{Step = 2}
    \label{apx_fig:shap_step_2_dir}
  \end{subfigure}
   \hfill
  \begin{subfigure}[t]{.24\textwidth}
    \centering
    \includegraphics[width=\linewidth]{Rebuttal-R1/appendix_figures/apx_shap_step_3_dir_v2.pdf}
    \caption{Step = 3}
    \label{apx_fig:shap_step_3_dir}
  \end{subfigure}
  \caption{SHAP Step Perturbation}\label{apx_fig:apx_shap_step}
\end{figure}

\section{Theoretical Justification of LIME Robustness Check 2}
\label{apx_sec:lime_thms}

In Section \ref{sec:lime_params} of the main paper we proposed robustness checks involving an increases to LIME's sample size $N$. The theorem below provides some theoretical justification for this suggestion.
 
        \paragraph{$e_{\text{LIME}}$ when Explaining a Linear Model}
            \label{sec:lime_linear_case}
             %In this section we we will show that there are conditions when the LIME coefficients have the property $e(x_j)\simeq \beta$ and use this fact to justify that  $e(x_j)$ can be evaluated with respect to the ground truth $\boldsymbol{\beta}$. Those conditions are discussed in the next paragraph. 
             
             %%LIME works by sampling data points around a given instance, then using these samples to train a weighted linear model as a local surrogate model that approximates the predictions of the original complex model in the vicinity of the instance. The local surrogate model is trained as a weighted ridge regressor. To obtain interpretable and sparse explanations, LIME discretizes numerical features and includes a regularization component to promote sparsity in the local surrogate models it creates. This regularization helps in selecting a subset of features that are most influential for predictions near the instance being explained, thus ensuring that the explanations are both interpretable and concise. 
             LIME's explanation model is a weighted ridge regressor trained on discretized data sampled around the instance being explained. By discretizing numerical features and using weighted ridge regression (rather than standard least squares), the possibility of LIME's data-alignment is compromised. The use of discretized features in LIME's ridge regression model can eliminate critical information about the model being explained. Additionally, the sparsity penalty which is used in ridge regression may zero out coefficients for features that actually have a non-zero marginal effect, distorting the picture one may obtain of the marginal effects. Also, employing weighted ridge regression introduces a weight matrix that further skews the estimation of marginal effects. This distortion can be observed in the closed-form solution of the weighted ridge regression, indicating potential limitations in LIME's ability to provide accurate estimates under certain conditions.
             
              Having identified the aforementioned properties of LIME's regression design as contributors to its lack of data-alignment, we derive conditions under which LIME would be data-aligned. We begin this effort by analyzing the large sample behavior of LIME coefficients in an idealized scenario where LIME does not use discretization and the model being explained is a linear function of i.i.d. normal features. Under such idealized assumptions, our Theorem \ref{thm:lime_thm} below shows that LIME's coefficients converge to the ground truth coefficients $\boldsymbol{\beta}$ as the number of samples $N$ tends to $\infty$. %Our Corollary 1 analyzes the closed form for the coefficients in order to ascertain the conditions for LIME's alignment in this idealized scenario. 
            \begin{theorem}
                \label{thm:lime_thm}
                Suppose $f:\mathbb{R}^D\rightarrow \mathbb{R}$ that maps $\mathbf{x} \mapsto \boldsymbol{\beta}^T \mathbf{x}$ is explained at a sample $\boldsymbol{\xi}$ drawn from $\mathcal{N}(\mathbf{0}_D, \bI_{D\times D})$ with LIME. Suppose further that LIME does not discretize features and samples i.i.d. data points $\mathbf{x}^{(1)},\ldots,\mathbf{x}^{(N)}$ from $ \mathcal{N}(\boldsymbol{\xi}, \bI_{D\times D})$ for its regression.
                 Then, in the large sample limit, LIME recovers $\boldsymbol{\beta}$.  
            \end{theorem}

             % Although LIME does not recover $\boldsymbol{\beta}$ with its default settings, Theorem 2 shows that $e_{\text{LIME}}(\mathbf{x}^{(i)}_j)$ converges to $\boldsymbol{\beta}$ in the large $n$ limit when explaining a linear model $\boldsymbol{\beta}^T \mathbf{x}$ where the data $\mathbf{x}$ are sampled from a standard multivariate normal distribution and LIME does not discretize continuous features.
             Theorem 1 uses an application of the law of large numbers to show that, in the large sample limit of random design matrices $\mathbf{X}$, LIME recovers $\boldsymbol{\beta}$. The practical implication of this theorem is that, when using LIME, $N$ should be large.  %For a non-asymptotic analysis, see the Appendix. %We also show that if $\nu$ is too small the LIME coefficients all go to $\mathbf{0}$.
             Its proof, along with a that of a require lemma, are given below. 

            % \begin{remark2}
            %     \label{nu_remark}
            %     Given a fixed number of samples $N$, $\nu$ being too small makes each of LIME's coefficients close to $0$. $\nu$ being very large can lead to a better approximation of $\boldsymbol{\beta}$. 
            % \end{remark2}

            % Remark \ref{nu_remark}, which is justified by our non-asymptotic analysis, suggests that when $N$ is fixed and finite, $\nu$  should also be set to be large. Notably, \citep{garreau20theoretical} obtained these same conclusions about $\nu$ and $N$ in the case where discretization is used. The proof for Theorem \ref{thm:lime_thm} as well as justification for Remark \ref{nu_remark} are provided below. %We also provide simulation experiments that empirically confirm the the assertions of Remark \ref{nu_remark}. 
            % Required lemmas are also proven.
        
        % \subsection{Proofs}
         \subsection{Theorem  
        \ref{thm:lime_thm} Proof}
        \paragraph{Theorem Setting}
        
            To investigate the limitations of LIME, we consider the simple case of explaining a linear function $f:\mathbb{R}^D \rightarrow \mathbb{R}$ such that $f(\mathbf{x}) = \boldsymbol{\beta}^T \mathbf{x}$ at the point $\boldsymbol{\xi}\in \mathbb{R}^D$ (which itself was sampled from $\mathcal{N}(\mathbf{0}_D, \bI_{D\times D})$) using a LIME explainer with the following properties. Suppose that LIME samples random vectors $\{\mathbf{x}^{(n)}\}_{n=1}^N$ with $\mathbf{x}^{(n)} \sim \mathcal{N}(\boldsymbol{\xi},\bI_{D\times D})$ for each $n$ and, given a bandwidth parameter $\nu > 0$, weights their distances to $\boldsymbol{\xi}$ using its default exponential kernel $\pi(\xi,\mathbf{x}) = \exp(-\norm{\boldsymbol{\xi} - \mathbf{x}}_2^2/2\nu^2)$. 
            
        We assume that LIME follows its default settings to solve a ridge regression problem, minimizing the following loss function to obtain a feature importance vector $\hat{\bu}$. 
        \begin{equation*}
            (\mathbf{y} - \mathbf{X}\mathbf{u})^T \mathbf{W}(\mathbf{y} - \mathbf{X}\mathbf{u}) + (\mathbf{u} - \mathbf{u}_0)^T \Delta(\mathbf{u} - \mathbf{u}_0)
        \end{equation*}
        In this loss function, $\mathbf{X} \in \mathbb{R}^{N\times D}$ is a data matrix  whose $n^{th}$ row is $\mathbf{x}^{(n)}$ and $\mathbf{W} \in \mathbb{R}^{N\times N}$ is a diagonal weight matrix for each of the $n=1,\ldots,N$ samples with $w_{nn} = \pi(\boldsymbol{\xi},\mathbf{x}^{(n)}) = \exp(-\norm{\boldsymbol{\xi} - \mathbf{x}^{(n)}}_2^2/2\nu^2)$. The response vector $\mathbf{y}\in \mathbb{R}^N$ is defined pointwise by $y^{(n)} = \boldsymbol{\beta}^T \mathbf{x}^{(n)}$. The ridge penalty parameter $\Delta\in \mathbb{R}^{D\times D}$ has $\Delta = \lambda \bI_{D\times D} \bI_{D\times D}$ so that $\Delta$ represents the ridge penalty with hyperparameter $\lambda=1$. As we shall see, $\lambda$ does not matter because LIME does not set it to scale with $N$. There is also another hyperparameter $\mathbf{u}_0 \in \mathbb{R}^D$ which is the target that $\mathbf{u}$ is to be shrunken toward. By default, it is $\mathbf{0}_D$. 

        In the proof, $n$ will be used to index the $N$ data samples. $i$ and $j$ will be used to index the $D$ features. Individual data samples in the data matrix $\mathbf{X}$ will be represented with superscript notation, e.g., $\mathbf{x}^{(n)}$. Individual features of a sample will be indexed with subscript notation, e.g. $x^{(n)}_j$. We need not refer to whole rows or columns of $\mathbf{W}$, so it will be indexed entry-wise with $w_{nn}$ for $n=1,\ldots,N$. Finally, we will use the Dirac delta function $\delta_{i=j}$, which is 1 when $i=j$ and 0 otherwise. 

            %\begin{theorem*}[\ref{thm:lime_thm}]
            \paragraph{\textbf{Theorem \ref{thm:lime_thm}}}
                %\label{thm:lime_thm}
                Suppose $f:\mathbb{R}^D\rightarrow \mathbb{R}$ that maps $\mathbf{x} \mapsto \boldsymbol{\beta}^T \mathbf{x}$ is explained at a sample $\boldsymbol{\xi}$ drawn from $\mathcal{N}(\mathbf{0}_D, \bI_{D\times D})$ with LIME. Suppose further that LIME does not discretize features and samples i.i.d. data points $\mathbf{x}^{(1)},\ldots,\mathbf{x}^{(N)}$ from $ \mathcal{N}(\boldsymbol{\xi}, \bI_{D\times D})$ for its regression.
                 Then, in the large sample limit, LIME recovers $\boldsymbol{\beta}$.   
            %\end{theorem*}

\noindent \emph{Proof.}
    The main idea underlying the proof is to apply the law of large numbers to a closed form solution to generalized ridge regression. Using the notation previously established, we can write a closed form solution of the regression as
    \begin{equation}
        \label{eq:closed_form}
        \hat{\bu} = (\bX^T \bW \bX + \bI)^{-1}(\bX^T \bW \by)
    \end{equation}
    Here, $\hat{\bu}$ is a point estimate using a single sample of data $\bX$. By simple algebraic manipulation we have, equivalently, 
    \begin{gather*}
         \hat{\bu} =  (\bX^T \bW \bX + \bI)^{-1}(\bX^T \bW \bX) \boldsymbol{\beta}
    \end{gather*}
    Let us define $\hat{\bS} = \frac{1}{N}(\bX^T \bW \bX + \bI) $ and $\hat{\bG} = \frac{1}{N}\bX^T \bW \bX$ so that $\hat{\bu} = \hat{\bS}^{-1}\hat{\bG}\boldsymbol{\beta}$.
    
    We may write $\hat{\bS}$ and $\hat{\bG}$ entry-wise as follows.
    \begin{gather*}
        (\hat{\bS})_{ij} = \frac{1}{N}\sum_{n=1}^N( x_i^{(n)} w_{nn} x_j^{(n)} + \delta_{i=j}) \\
        (\hat{\bG})_{ij} = \frac{1}{N} \sum_{n=1}^N x_i^{(n)} w_{nn} x_j^{(n)} 
    \end{gather*}
    Above, $\delta_{i=j}$ represents the Dirac delta function which is 1 when $i=j$ and 0 otherwise. For fixed $i,j$, the sequence of random variables $x_i^{(1)} w_{11} x_j^{(1)},\ldots,x_i^{(N)} w_{NN} x_j^{(NN)}$ are mutually independent because the samples $x^{(1)},\ldots,x^{(N)}$ are i.i.d. random vectors. 
    By writing the members of the sequence in the form
    \begin{equation*}   
        x_i^{(n)} e^{-\norm{\bxi - \bx^{(n)}}_2^2/2\nu^2} x_j^{(n)}
    \end{equation*}
    we can see that are also identically distributed. %The mean of these terms is finite, as $w_{nn}\in [0,1]$ for all $n=1,\ldots,N$ means that the mean is upper bounded by the square of the largest entry in $\bxi$. 
    %https://math.stackexchange.com/questions/4153550/existence-of-expectation-of-product-of-two-random-variables
    % https://math.stackexchange.com/questions/1043082/bound-the-variance-of-the-product-of-two-random-varables
    By Lemma \ref{finitemean}, the terms have finite means. 
    
Thus, by the strong law of large numbers, the sample averages
$\frac{1}{N}\sum_{n=1}^N x_i^{(n)} w_{nn} x_j^{(n)}+\delta_{ij}$
converge to
$\mathds{E}_{\bx \sim \mathcal{N}(\bxi,\mathbf{1}_{D})}\!\left[x_i e^{-\norm{\bxi-\bx}_2^2/2\nu^2}x_j\right]$
almost everywhere. This means that the ridge adjustment vanishes (converges to a null matrix). Likewise, the averages $\frac{1}{N}\sum_{n=1}^N x_i^{(n)} w_{nn} x_j^{(n)}$ converge to the same value $\mathds{E}_{\bx \sim \mathcal{N}(\bxi,\mathbf{1}_{D})}[x_i e^{-\norm{\bxi-\bx}_2^2/2\nu^2}x_j]$ almost everywhere. Thus, in the large sample limit, $\hat{\bS}$ and $\hat{\bG}$ converge to the same matrix. \textbf{It follows that, in the large sample limit, $\hat{\bS}^{-1}$ and $\hat{\bG}$ cancel each other, and so we expect that $\hat{\bu} \rightarrow \boldsymbol{\beta}$.} That is, $\boldsymbol{\beta}$ is recovered by the regression in spite of the ridge penalty and sample weighting. 
    
    % \qed
    
    
    
   
\begin{lemma}
    \label{finitemean}
   $x_i^{(n)} w_{nn} x_j^{(n)}$ has a finite mean for all $i,j,n$. 
\end{lemma}
    \noindent \emph{Proof.} Fix some $i,j,n$. 
    Since a random variable having finite variance implies that it has a finite mean, it is sufficient to show that $x_i^{(n)} w_{nn} x_j^{(n)}$ has a finite variance. We do this using the fact that, for dependent random variables $A$ and $B$
    \begin{equation}
        \label{varbd}
        \mathrm{Var}(AB) \leq  2\mathrm{Var}(A)\lVert B\rVert_\infty^2+2(\mathbb E[A])^2\mathrm{Var}(B)
    \end{equation}
    When using this inequality, we set $A:=x_i^{(n)} x_j^{(n)}$ and $B:=w_{nn}$.

    We will establish that each term on the RHS of Equation \ref{varbd} is finite in order to show that $\mathrm{Var}(AB)$ is finite. To see that $\mathrm{Var}(A)$ is finite, note that $x_i^{(n)} x_j^{(n)}$ are Chi-square random variables with 1 degree of freedom, which have finite variance. The weights $B=w_{nn}$ are bounded in the interval $(0,1]$, so $\norm{B}_\infty^2$ is finite. Chi-square random variables have finite means, so $\mathds{E}[A]$ is finite. To see that $\mathrm{Var}(B)$ is finite, note again that $B$ is bounded in $(0,1]$. Then $\mathds{E}[B] \leq 1$. As $B$ is bounded, $B^2$ is also bounded. Hence, the variance $\mathrm{Var}(B) = \mathds{E}[B^2]-\mathds{E}[B]^2$ is a subtraction of finite quantities and therefore also finite. Evidently, each term in the RHS of Equation \ref{varbd} is finite, and so $x_i^{(n)} w_{nn} x_j^{(n)}$ has a finite variance. It follows that the mean of $x_i^{(n)} w_{nn} x_j^{(n)}$ is also finite.  \qed

%     \begin{lemma}
%         \label{norm_bd}
%         Under the assumptions of  Theorem \ref{thm:lime_thm}, and given a random  design matrix $\bX \in \mathbb{R}^{N\times D}$,
%         \[\norm{(\bX^T \bW \bX + \bI)^{-1}\bX \bW \bX))}_2 \in O\left(\frac{1}{\sigma_{\min}(\bW)^D}\right) \]
%     \end{lemma}
    
%   \noindent \emph{Proof.} For the sake of brevity we will define $\bA := \bX^T \bW \bX$ and $\bB := \bA + \bI$. To bound $\norm{\bB \bA}_2$, we will begin by bounding $\norm{\bB}$. We employ the following upper bound for the norm of a $D \times D$ matrix's inverse. We use the notation $\sigma_k(\cdot)$ for the $k$th singular value of a matrix and $\sigma_{\min}(\cdot)$ and $\sigma_{\max}(\cdot)$ for the minimum and maximum singular values, respectively. 
%   \begin{equation*}
%       \norm{\bB}^{-1} \leq \frac{\norm{\bB}^{D-1}}{\mid \mathrm{det}(B)\mid}
%   \end{equation*}
%   Moreover,
%   \begin{gather*}
%       \frac{\norm{\bB}^{D-1}}{\mid \mathrm{det}(B)\mid} \leq \frac{\norm{\bA}^{D-1} + 1}{\mid \mathrm{det}(B)\mid} 
%       =  \frac{\norm{\bA}^{D-1} + 1}{\prod_k \sigma_k(\bB)} 
%       \leq \frac{\norm{\bA}^{D-1} + 1}{\sigma_{\min}(\bB)^D} 
%       \leq  \frac{\norm{\bA}^{D-1} + 1}{(\sigma_{\min}(\bA) + 1)^D} 
%   \end{gather*}
%     The first inequality follows by Cauchy-Schwarz and the facts that  $D>1$ and power functions are monotonic. The last inequality follows from the fact that $\bA$ and $\bI$ are real and symmetric.
%     Now, since the spectral norm $\norm{\cdot}_2$ is bounded above by the spectral radius, after applying Cauchy-Schwarz and bounding the denominator we obtain that
%     \begin{gather*}
%         \frac{\norm{\bA}^{D-1} + 1}{(\sigma_{\min}(\bA) + 1)^D} \leq \frac{\norm{\bX^T}^{D-1}\norm{\bW}^{D-1}\norm{\bX}^{D-1} + 1}{\sigma_{\min}(\bA)^D} \leq \frac{\sigma_{\max}(\bX^T)^{D-1}\sigma_{\max}(\bW)^{D-1}\sigma_{\max}(\bX)^{D-1} + 1}{\sigma_{\min}(\bA)^D}
%     \end{gather*}
%     We can now bound $\norm{\bB^{-1}\bA}$ using Cauchy-Schwarz as follows.
%     \begin{equation*}
%         \norm{\bB^{-1}\bA} \leq \frac{\sigma_{\max}(\bX^T)^{D}\sigma_{\max}(\bW)^{D}\sigma_{\max}(\bX)^{D} + \sigma_{\max}(\bX^T)\sigma_{\max}(\bW)\sigma_{\max}(\bX)}{\sigma_{\min}(\bX^T)^{D}\sigma_{\min}(\bW)^{D}\sigma_{\min}(\bX)^{D}}
%     \end{equation*}
%     Now, since $0<w_{nn}\leq 1$ for all $n=1,\ldots,N$, and since singular values are invariant under transposition, we can give the following upper bound which should be viewed as a function of $\sigma_{\min}(\bW)$ since we assumed that $\bX$ is fixed. 
%     % \begin{equation*}
%     %     \norm{\bB^{-1}\bA} \leq \left(\frac{\sigma_{\max}(\bX^T)^{D}\sigma_{\max}(\bX)^{D} + \sigma_{\max}(\bX^T)\sigma_{\max}(\bX)}{\sigma_{\min}(\bX^T)^{D}\sigma_{\min}(\bX)^{D}}\right) \frac{1}{\sigma_{\min}(\bW)^D}
%     % \end{equation*}
%     % Since singular values are invariant under tranpsosition, we can simplify this into
%     \begin{equation*}
%         \norm{\bB^{-1}\bA} \leq \left(\frac{\sigma_{\max}(\bX)^{2D} + \sigma_{\max}(\bX)^2}{\sigma_{\min}(\bX)^{2D}}\right) \frac{1}{\sigma_{\min}(\bW)^D}
%     \end{equation*}

% \begin{remark2}[Connection to condition numbers]
%     \label{condition_number_remark}
%     $\norm{\bB^{-1}\bA}$ may be made larger when $\bW$ is ill-conditioned.
% \end{remark2}
%         Recall that the condition number $\kappa(\cdot)\geq 1$ induced by the spectral norm of a matrix is the ratio of its largest and smallest singular values. The bound of Lemma \ref{norm_bd} can be written in terms of condition numbers. By writing it in this form, it can be inferred that $\norm{\bB^{-1}\bA}$ is made larger when $\bW$ is ill-conditioned. 
%     \begin{equation*}
%         \norm{\bB^{-1}\bA} \leq \left(\kappa(\bX)^{2D} + \frac{\kappa(\bX)^2}{\sigma_{\min}(\bX)^{2D-2}}\right) \kappa(\bW)^D 
%     \end{equation*}
%     The ill-conditioning of $\bW$ occurs when $\bW$ is close to being singular, which can happen when $\nu$ is near $0$ since this would cause all entries of $\bW$ to be either 0 or close to it. 


    % \subsection{Details on Remark \ref{nu_remark}}
    %             Remark \ref{nu_remark} suggested that, fixing a design matrix $\bX$ consisting of $N$ samples, then if 1) $\nu$ is too small then each of the coefficients of $\hat{\bu}$ will be near $0$ and 2) making $\nu$ very large may yield a better  approximation of $\boldsymbol{\beta}$. Both of these statements have to do with the construction of the weight matrix $\bW$ using a given $\bX$. 

    %             Recalling that $\hat{\bu}=(\bX^T \bW \bX+\bI)^{-1}\bX^T \bW \bX \boldsymbol{\beta}$ and that, by Theorem \ref{thm:lime_thm} $\hat{\bu}$ tends to $\boldsymbol{\beta}$ in the large sample limit, we can infer that it is desirable for $(\bX^T \bW \bX+\bI)^{-1}\bX^T \bW \bX$ to be near the identity matrix and, therefore, to have norm $1$. Lemma \ref{norm_bd} gives an upper bound on $\norm{(\bX^T \bW \bX+\bI)^{-1}\bX^T \bW \bX}$, so we can use it to investigate what makes this upper bound large. By analyzing $\norm{(\bX^T \bW \bX+\bI)^{-1}\bX^T \bW \bX}$, we can get a rough estimate of how much $\boldsymbol{\beta}$ is distorted. % and possibly more distant from the identity matrix.
                
    %             We can deduce from the bound given in Lemma \ref{norm_bd} that setting $\nu$ too small makes it possible for $\norm{\bB^{-1}\bA}$ to be large. This is because $\sigma_{\min}(\bW)$ is just $\min_{n=1,\ldots,N} w_{nn}=\min_{n=1,\ldots,N} e^{-\norm{x^{(n)}-\bxi}_2^2/2\nu^2}$ and smaller $\nu$ increases the decay rate of the $w_{nn}$, leading to $\frac{1}{\sigma_{\min}(\bW)^D}$ being very large. On the other hand, since in the context of Lemma \ref{norm_bd} $\bX$ is fixed, larger values of $\nu$ make $\frac{1}{\sigma_{\min}(\bW)^D}$ closer to $1$ and therefore lead to smaller upper bounds for $\norm{\bB^{-1}\bA}$.

    %             % https://mathoverflow.net/questions/353085/anti-concentration-inequalities-lower-bound-on-realized-second-moment
    %             To see this more concretely, suppose that a random design matrix $\bX$ with rows sampled from $\mathcal{N}(\mathbf{0}_{D}, \bI_{D\times D})$ is fixed. Furthermore, assume that each row has $\norm{\bx^{(n)}} \leq 1$. Then, we can use the anti-concentration of the norms of the rows $\bx^{(n)}$ to show that when $\nu$ is too small, all values in $\bW$ will be smaller than a value $\epsilon > 0$. Simple computations show that $w_{nn} < \epsilon$ when $\norm{\bx^{(n)}} > t:= \sqrt{2}\sqrt{-\log(\epsilon)/\nu^2}$. If we assume that $t\leq \mathds{E}[\norm{\bx}^{(n)}]$, and by extension impose the restriction $\nu \geq \sqrt{2}\sqrt{-\log(\epsilon)/D}$, we may use an application of the Azuma-Hoeffding inequality \citep{hoeffding1994probability} to estimate the probability that $\norm{\bx^{(n)}}$ is sufficiently large that $w_{nn} < \epsilon$. 
    %             \begin{equation*}
    %                 \mathds{P}[\norm{\bx^{(n)}} > t] \geq 1 - \frac{1}{2D}\exp(-2\sqrt{2}\sqrt{-\log(\epsilon)/\nu^2}-2\sqrt{D})
    %             \end{equation*}
    %             If, for example $D=100$, $\epsilon=\frac{1}{1000}$, and $\nu$ takes on the smallest value for the inequality to be used (about $.38$), then $\mathds{P}[\norm{\bx^{(n)}} > t]=1$. That is, $w_{nn}<\epsilon$  almost surely.


    %         %https://arxiv.org/abs/0810.0800
    %         % , given $\nu < 1$ we can see from the concentration of a gaussian vector that 
    %         %fixing eps > 0, wnn will be less than epsilon whenever $\norm{x-xi}_2 > sqrt(-2log(eps)/nu^2)$. Then, setting t =   (-2log(eps)/nu^2)^1/4, we see from a subgaussian tail bound that
    %         %P[\norm{x-xi}_2 > t] > 
            
    %     \subsection{Insights From Theoretical Results}
    %         \begin{itemize}
    %             \item Although explanations using discretized features may be more easily interpretable, they make the explanations less data-aligned
    %             \item Users should generally opt to set $N$ to be large 
    %             \item Users should generally set $\nu$ to be large, otherwise the LIME coefficients may be become zero
    %             %\item LIME becomes more unfaithful as the number of features $d$ grows. 
    %             \item The ridge penalty matters only when $\lambda \in O(N)$
    %         \end{itemize}
    
    %         \subsection{Simulation Experiments}
    %             \label{apx_sec:simulation_experiments}
    %             We provide simulations to test the assertions of Theorem \ref{thm:lime_thm} and Remark \ref{nu_remark} experimentally. We adopt the settings of Theorem \ref{thm:lime_thm} and experimentally test the roles of $N$, $\nu$ on the convergence of weighted ridge regression using LIME's exponential kernel and the ridge penalty set to $\Delta=1$.% In the second simulation, we use implementation of LIME by its authors and our loan dataset. In this simulation we investigate how the discretization setting affects the actual LIME implementation. In particular, we investigate the ratios $e_{\text{LIME}}(x^{(i)}_j)/\beta_j$. 
    

                
    %             \paragraph{Role of $\nu$ and $N$ on convergence}
    %                 We give empirical confirmation of the assertions of our theoretical results and show the experimental results in Figure \ref{apx_fig:nu_N_convergence}. To do this, we create ten random design matrices $\mathbf{X}$ for each combination of $N=10,100,1000,10000,1000000$ rows and $D=5,10,100$ columns. For each of these matrices, we compute $\norm{\bB^{-1}\bA - \mathbf{I}}$ using each of $\nu=.1,1,10,10000$. We report the mean value of $\norm{\bB^{-1}\bA - \mathbf{I}}$ as $N$ is varied.


    %                 Figure \ref{apx_fig:nu_N_convergence} is a log-log plot of dimension versus the distance $\norm{\bB^{-1}\bA - \mathbf{I}}$. It shows that small $\nu$ leads to slower convergence and larger $\nu$ leads to faster convergence. It also shows that increasing $N$ generally leads to $\norm{\bB^{-1}\bA - \bI}$ being smaller. That the curves are approximately linear indicates the presence of a power law relationship. These results constitute supporting evidence for our claim that $\nu$ and $N$ should be set to large values. In particular, it is advantageous to set them to be larger than their defaults of $\nu=.75\sqrt{D}$ and $N=1000$. 
                    
    %                 % It shows that, for not too small $\nu$, $\norm{\bB^{-1}\bA - \mathbf{I}}$ converges to $1$, not needing $n$ to be unreasonably large. For example, it seems to be sufficient to use the default settings of $\nu = .75\sqrt{d}$ and $n=1000$. However, it is apparent from Figure \ref{apx_fig:nu_N_convergence} that the largest settings $nu=10000$ gives an advantage for smaller $n$.  

    %                 % Figure \ref{apx_fig:nu_D_convergence} illustrates the relations of $d$ and $\nu$ with $\norm{\bB^{-1}\bA - \mathbf{I}}$. It shows that larger $\nu$ are associated with smaller $\norm{\bB^{-1}\bA - \mathbf{I}}$ and larger $d$ are associated with larger $\norm{\bB^{-1}\bA - \mathbf{I}}$. This implies that, while increasing $\nu$ is beneficial to LIME's fidelity, larger $d$ make attaining higher fidelity more difficult (especially for small $\nu$). This backs up the results shown in Figure \ref{fig:lime_shap_vary_d}, that is, that increasing $d$ is generally associated with lower data-alignment. 
    
    %             % \begin{figure}[ht!]
    %             % \centering
    %             %     \includegraphics[height=.2\textwidth]{AppendixFigures/apx_lime_conv-figure0.pdf}%{Figures/fig21.pdf}
    %             %     \caption{Role of $\nu,n$ on $\norm{\mathbf{H}\mathbf{X} - \mathbf{I}}$, with $d=4$.}
    %             %     \label{apx_fig:nu_N_convergence}
    %             % \end{figure}

    %             % %\paragraph{Role of $d$ on convergence}
    
    %             % \begin{figure}[ht!]
    %             % \centering
    %             %     \includegraphics[height=.2\textwidth]{AppendixFigures/apx_lime_conv-figure1.pdf}%{Figures/fig21.pdf}
    %             %     \caption{Role of $d$ on $\norm{\mathbf{H}\mathbf{X} - \mathbf{I}}$ with $n=50000$}
    %             %     \label{apx_fig:nu_D_convergence}
    %             % \end{figure}

    %             \begin{figure}[ht!]
    %             \centering
    %                 %\begin{subfigure}{}
    %                 \centering
    %                 \includegraphics[width=.5\textwidth]{Figures/new_apx_lime_conv.pdf}%{Figures/fig21.pdf}
    %                 \caption{Log-log plot of the how $\nu,N$ influence $\norm{\bB^{-1}\bA - \mathbf{I}}$. }
    %                 \label{apx_fig:nu_N_convergence}
    %                 %\end{subfigure}%
    %                 % \hspace{3cm}
    %                 % \begin{subfigure}{.2\textwidth}
    %                 % \centering
    %                 % %\fbox{\includegraphics[width=\textwidth,trim=0 10 70 70 mm, clip=true]{Figures/fig21.pdf}}%{Final_Figures/fig20.png}
    %                 % \includegraphics[height=\textwidth]{AppendixFigures/apx_lime_conv-figure1.pdf}%{Figures/fig21.pdf}
    %                 % \caption{Role of $d$ on $\norm{\bB^{-1}\bA - \mathbf{I}}$ with $n=50000$}
    %                 % \label{apx_fig:nu_D_convergence}
    %                 % \end{subfigure}
    %             \end{figure}
    % %             \paragraph{Upper Bounds on $\norm{ (\mathbf{X}^T \mathbf{W} \mathbf{X} + \mathbf{I})^{-1}\mathbf{X}^T\mathbf{W}\mathbf{X}}_2$}
    % %                 We show in Figure \ref{apx_fig:nu_N_convergence} that, except when $\nu$ is too small ($\nu<.2$), increasing the number of samples $n$ generally nets a significant decrease in an upper bound for the error of LIME's ridge problem. We use the condition number of  $\norm{ (\mathbf{X}^T \mathbf{W} \mathbf{X} + \mathbf{I})^{-1}\mathbf{X}^T\mathbf{W}\mathbf{X}}_2$ as an upper bound.     
                
    % %                 \begin{figure}[ht!]
    % %                 \centering
    % %                     \includegraphics[height=.2\textwidth]{AppendixFigures/apx_lime_conv-figure2.pdf}%{Figures/fig21.pdf}
    % %                     \caption{Role of $\nu,n$ on upper bound of $\norm{\mathbf{H}\mathbf{X} - \mathbf{I}}$}
    % %                     \label{apx_fig:nu_N_convergence}
    % %                 \end{figure}

    
    % %             \paragraph{Role of Discretization}
    % %                 We provide evidence of our claim that users should not let LIME discretize numerical features if recovery of the marginal effects is their aim by comparing LIME explainers with and without discretization. For our comparison, we use LIME with and without discretization to explain a linear regressor $f(x) = .05 + .05x_1+ .2x_2+ .8x_3 -.5x_4 +x_5$ with $x_i \sim \mathcal{N}(0,1)$ for all $i$. $\nu$ is set to the default value. We pick 100 test set instances and explain each two sets of 100 LIME explainers. The first set of LIME explainers do not use feature discretization while the second set does. The $i^{th}$ explainer has its random seed set to $n$ in both sets. 

    % %             %As we claimed in Section 3, the LIME explainers with no feature discretization generally have $C_i$ closer to $1$ than the LIME explainers with discretization. The difference is especially stark with $C_{1}$, which is always 1 with no discretization and between -10 and 20 with discretization. Note that when $C_i$ is negative, LIME has misidenified how feature $i$ impacts the output of $f$. Moreover, in accordance with Theorem 1 and \citep{garreau20theoretical}, we see that $C_i$ can differ significantly across $i$, indicating that $e(x_j)$ is not directly proportional to $\boldsymbol{\beta}$ regardless of whether discretization is used. 

    % %             For each set of explainers, we plot the ratios $C_i:= e(x_j)_i/\beta_i$ for $i=1,\ldots,5$ in Figure \ref{apx_fig:c_demo}. The most striking result is that all $C_j$ are centered around $\frac{1}{3}$ and contained in $[-1,1]$ when discretization is turned off, but $C_1$ takes on values in $[-10,10]$ with a mean of $\sim 3.70$ and $C_2,C_3,C_4,C_5$ have respective means of $.52$, $.06$, $.03$, and $-.001$ when discretization is allowed. Thus, when discretization is allowed, LIME may distort the $\beta_j$ more and have severe inconsistency between how much distortion is present across each $j$. That the $C_j$ have nearly identical means when discretization is not used indicates that, on average, the correct rankings of the $X_j$ are preserved in the $e(x_j)_i$. The $C_j$ being very different when discretization is used indicates that such LIME rankings would (almost) never match the ground truth ranking. 
                
    % %             %use fbox to tailor crop
    % %             \begin{figure}[ht!]
    % %             \centering
    % %                 \begin{subfigure}{.28\textwidth}
    % %                 \centering
    % %                 %\fbox{\includegraphics[width=\textwidth,trim=0 10 70 70 mm, clip=true]{Figures/fig21.pdf}}%{Final_Figures/fig20.png}
    % %                 \includegraphics[width=\textwidth]{AppendixFigures/apx_lime_proportionality-figure0.pdf}%{Figures/fig21.pdf}
    % %                 \caption{No discretization}
    % %                 \label{apx_fig:empirical_c}
    % %                 \end{subfigure}%
    % %                 \hspace{3cm}
    % %                 \begin{subfigure}{.28\textwidth}
    % %                 \centering
    % %                 %\fbox{\includegraphics[width=\textwidth,trim=0 10 70 70 mm, clip=true]{Figures/fig21.pdf}}%{Final_Figures/fig20.png}
    % %                 \includegraphics[width=\textwidth]{AppendixFigures/apx_lime_proportionality-figure1.pdf}%{Figures/fig21.pdf}
    % %                 \caption{Discretization allowed}
    % %                 \label{apx_fig:theoretical_cs}
    % %                 \end{subfigure}%
    % %             \caption{Empirical $C_j$ for 100 users from 100 LIMEs with/without discretization, and all other parameters set to the defaults.}
    % %             \label{apx_fig:c_demo}
    % %             \end{figure}
    

    


    
    
  
 
  
        \section{Computational Cost and Runtime Accounting}
\label{apx:computation}

This section documents the computational resources and runtime required to conduct the experiments reported in the paper. The scale of computation reflects both the number of simulated datasets and the intensive nature of post hoc explanation methods, particularly SHAP.

\subsection{Model Construction}

Across all experiments, we generated 81 synthetic datasets. For each dataset, we estimated predictive models for two distinct experimental settings corresponding to Sections~5 and~7 of the paper.

\textbf{Models used in Section~5.}  
For the analysis in Section~5, we trained 12 XGBoost models per dataset with systematically varying predictive accuracy. To obtain evenly spaced accuracy levels in the range 0.7--0.9, we conducted a grid search over a $60 \times 10 \times 4$ grid spanning the number of estimators, tree depth, and training-set size. Each parameter configuration required approximately 3 seconds of runtime. Aggregated across all datasets, model construction for Section~5 required approximately
\[
60 \times 10 \times 4 \times 3 \times 81 / (3600 \times 24) \approx 6.75
\]
days of continuous computation.

\textbf{Models used in Section~7.}  
For the analysis in Section~7, we constructed a Rashomon set of 10 high-accuracy models per dataset using nine machine-learning algorithms (including XGBoost, CatBoost, and LightGBM). For each algorithm, we explored an average of 81 hyperparameter configurations using five-fold cross-validation, with each configuration requiring approximately 2 seconds. Aggregated across datasets, model construction for Section~7 required approximately
\[
81 \times 5 \times 2 \times 81 / (3600 \times 24) \approx 0.76
\]
days of computation.

All model training was conducted on a workstation equipped with an Intel i7 CPU and 32~GB RAM. In total, predictive model construction required approximately \textbf{7.5 days} of continuous computation.

\subsection{Post hoc Explanation Generation}

The computational burden increases substantially in the explanation stage. For each dataset--model pair, we generated SHAP and LIME explanations for 100 test instances. To compute the directionality metric, we additionally generated six perturbed versions of each instance ($m=-3,-2,-1,1,2,3$), resulting in 700 explanation tasks per model.

All explanation experiments were executed on a dedicated server equipped with 80 CPU cores (2.5~GHz) and 256~GB RAM, with up to 30 parallel processes per task.

\textbf{SHAP.}  
SHAP proved particularly computationally intensive. For models used in Section~5, generating a single SHAP explanation required approximately 30 seconds on average. For the high-accuracy models used in Section~7, the average runtime increased to approximately 90 seconds per instance due to heterogeneity in model structure and SHAP compatibility. Aggregating across all datasets and models, SHAP explanation required approximately
\[
\frac{81 \times (12 \times 700 \times 30 + 10 \times 700 \times 90)}{3600 \times 24 \times 30}
\approx 27.6
\]
days of continuous server-level computation.

\textbf{LIME.}  
By contrast, LIME was substantially faster, requiring approximately 0.1 seconds per instance on average. Across all datasets and models, LIME explanations required approximately 0.05 days of computation and did not materially affect total runtime.

\subsection{Total Runtime}

Taken together, our experimental pipeline required approximately \textbf{35.6 days} of computation: 7.5 days for model construction and 27.6 days for SHAP-based explanation. These figures underscore why evaluating SHAP on entire test sets would be computationally infeasible in this setting and motivate our use of a fixed subsample of test instances in the main analysis.
 
  



\end{document}
